% Generated by scripts/correspondence.py from the docstrings of the model and
% correspondence/j-points.toml. Do not edit: run `scripts/correspondence.py table`.

\chapter{The correspondence with the C code}\label{chap:correspondence}
% the long names of Lean and of the C code cannot be broken: let the spaces stretch
\sloppy

The model is checked against \texttt{j-points} v3.0.2 (commit \texttt{3633f6328451}), the files \texttt{j-points.c}, \texttt{j-sift.c}, \texttt{j-points.h}. Every line of code of them lies in exactly one recorded range: a range is either modeled by the Lean declarations named, or excluded from the model with a reason. Each recorded range is checked against the text it had when it was reviewed.

\begin{center}
\begin{tabular}{lrrrr}
file & lines of code & modeled & excluded & not recorded \\
\hline
\texttt{j-points.c} & 1844 & 1564 & 280 & 0 \\
\texttt{j-sift.c} & 846 & 729 & 117 & 0 \\
\texttt{j-points.h} & 72 & 15 & 57 & 0 \\
\hline
total & 2762 & 2308 & 454 & 0 \\
\end{tabular}
\end{center}

\section*{\texttt{j-points.c}}
\begin{itemize}
\item \csource{j-points.c}{29}{38}: excluded (other): the feature-test macro \_GNU\_SOURCE, which declares asprintf for the threads (line 2117) and makes the strtol of scan\_long glibc's \_\_isoc23\_strtol, the same in base 10 (JPoints.Cli.strtol), and the includes of the headers of the C library, of POSIX threads, of GMP and of j-points.h; the model has no headers
\item \csource{j-points.c}{44}{44}: modeled by \lname{JPoints.defaultSize}
\item \csource{j-points.c}{46}{48}: modeled by \lname{JPoints.version}
\item \csource{j-points.c}{54}{57}: excluded (instrumentation): the table prime[] of the debug build (DEBUG > 0), the five odd primes up to 13, with the \#if that selects it and the \#else of the table of the default build (lines 58-61)
\item \csource{j-points.c}{58}{61}: modeled by \lname{JPoints.primeTable}
\item \csource{j-points.c}{62}{62}: excluded (instrumentation): the end of the selection of the table prime[] by DEBUG
\item \csource{j-points.c}{63}{63}: modeled by \lname{JPoints.Plan}
  (deviation of \lname{JPoints.Plan}: a structure for the globals; natural numbers for the \texttt{long} indices of \texttt{pnn} and for the \texttt{int} \texttt{sieve\_mode}; a Boolean for the \texttt{int} flag \texttt{tube\_cut}; natural numbers for the \texttt{long}s \texttt{sieve\_primes1}, \texttt{sieve\_primes2} and \texttt{sieve\_primes3} after the choice, which leaves them nonnegative (before it they hold what \texttt{read\_input} leaves, \texttt{JPoints.Cli.Input}).)
\item \csource{j-points.c}{66}{66}: modeled by \lname{JPoints.Cli.Input}
  (deviation of \lname{JPoints.Cli.Input}: a structure for the globals; an array of seven \texttt{Int}s for the seven GMP integers of \texttt{coeffs}; \texttt{Int}s for the \texttt{long}s, whose ranges after \texttt{read\_input} later lemmas give; \texttt{Bool}s for the \texttt{int} flags; a string for \texttt{print\_format}, the empty string for the null pointer that it is before \texttt{read\_input} sets it; the flags \texttt{set} of the entries of \texttt{cost\_names} as an array of their own.)
\item \csource{j-points.c}{81}{84}: excluded (other): the defaults COST\_AND\_VW, COST\_SIZE\_VW and COST\_TEST2\_VW of the cost model for the vector build VW = 4 (AVX2), with the \#if that selects them; the build modeled, VW = 1, has the defaults of the \#else (lines 85-91)
\item \csource{j-points.c}{85}{88}: modeled by \lname{JPoints.costAnd}
\item \csource{j-points.c}{89}{89}: modeled by \lname{JPoints.costSize}
\item \csource{j-points.c}{90}{90}: modeled by \lname{JPoints.costTest2}
\item \csource{j-points.c}{91}{94}: modeled by \lname{JPoints.costAnd}
\item \csource{j-points.c}{95}{99}: modeled by \lname{JPoints.costSize}
\item \csource{j-points.c}{100}{103}: modeled by \lname{JPoints.costWord}
\item \csource{j-points.c}{104}{106}: modeled by \lname{JPoints.costTest2}
\item \csource{j-points.c}{107}{110}: modeled by \lname{JPoints.costBit}
\item \csource{j-points.c}{111}{114}: modeled by \lname{JPoints.costRow3}
\item \csource{j-points.c}{115}{117}: modeled by \lname{JPoints.costTest3}
\item \csource{j-points.c}{118}{120}: modeled by \lname{JPoints.costExact}
\item \csource{j-points.c}{121}{124}: modeled by \lname{JPoints.costRow}
\item \csource{j-points.c}{125}{127}: modeled by \lname{JPoints.costCell}
\item \csource{j-points.c}{128}{130}: modeled by \lname{JPoints.costTword}
\item \csource{j-points.c}{131}{133}: modeled by \lname{JPoints.costTtest}
\item \csource{j-points.c}{134}{136}: modeled by \lname{JPoints.costTrow}
\item \csource{j-points.c}{137}{139}: modeled by \lname{JPoints.costRrow}
\item \csource{j-points.c}{140}{143}: modeled by \lname{JPoints.costRexcl}
\item \csource{j-points.c}{144}{146}: modeled by \lname{JPoints.costRcut}
\item \csource{j-points.c}{147}{149}: modeled by \lname{JPoints.costRange}
\item \csource{j-points.c}{150}{152}: modeled by \lname{JPoints.costPass}
\item \csource{j-points.c}{153}{155}: modeled by \lname{JPoints.costTable}
\item \csource{j-points.c}{156}{159}: modeled by \lname{JPoints.costInit}
\item \csource{j-points.c}{160}{161}: modeled by \lname{JPoints.samplePairs}
\item \csource{j-points.c}{164}{170}: modeled by \lname{JPoints.Cli.Costs}
  (deviation of \lname{JPoints.Cli.Costs}: \texttt{Float}s, as the cost model runs in \texttt{Float}, which is the hardware double.)
\item \csource{j-points.c}{171}{178}: modeled by \lname{JPoints.Cli.costNames}
  (deviation of \lname{JPoints.Cli.costNames}: the pointer of an entry to its static double as the functions that read and write the field of \texttt{JPoints.Cli.Costs} that it points to; the flags \texttt{set} of the entries in the globals (\texttt{JPoints.Cli.Input.costSet}), since the names and the pointers do not change.)
\item \csource{j-points.c}{179}{179}: modeled by \lname{JPoints.Cli.numCosts}
\item \csource{j-points.c}{180}{180}: modeled by \lname{JPoints.Cli.Input}
  (deviation of \lname{JPoints.Cli.Input}: a structure for the globals; an array of seven \texttt{Int}s for the seven GMP integers of \texttt{coeffs}; \texttt{Int}s for the \texttt{long}s, whose ranges after \texttt{read\_input} later lemmas give; \texttt{Bool}s for the \texttt{int} flags; a string for \texttt{print\_format}, the empty string for the null pointer that it is before \texttt{read\_input} sets it; the flags \texttt{set} of the entries of \texttt{cost\_names} as an array of their own.)
\item \csource{j-points.c}{182}{182}: modeled by \lname{JPoints.Tables}, \lname{JPoints.Tables.initial}
  (deviation of \lname{JPoints.Tables}: flat arrays for the arrays of the C code, with its index arithmetic, in a structure, for the globals.)
\item \csource{j-points.c}{184}{184}: modeled by \lname{JPoints.Entry}
  (deviation of \lname{JPoints.Entry}: natural numbers for the \texttt{long}s \texttt{p} and \texttt{n}, which the C code sets to a prime of the table and an index of the table.)
\item \csource{j-points.c}{186}{186}: modeled by \lname{JPoints.Tables}, \lname{JPoints.Tables.initial}
  (deviation of \lname{JPoints.Tables}: flat arrays for the arrays of the C code, with its index arithmetic, in a structure, for the globals.)
\item \csource{j-points.c}{188}{188}: modeled by \lname{JPoints.Tables}, \lname{JPoints.Tables.initial}
  (deviation of \lname{JPoints.Tables}: flat arrays for the arrays of the C code, with its index arithmetic, in a structure, for the globals.)
\item \csource{j-points.c}{191}{191}: modeled by \lname{JPoints.CurveTables}, \lname{JPoints.CurveTables.initial}
  (deviation of \lname{JPoints.CurveTables}: flat arrays for the arrays of the C code, with its index arithmetic, in a structure, for the globals; arrays of Booleans for the arrays of \texttt{unsigned char} that hold \texttt{0} and \texttt{1}.)
\item \csource{j-points.c}{194}{194}: modeled by \lname{JPoints.CurveTables}, \lname{JPoints.CurveTables.initial}
  (deviation of \lname{JPoints.CurveTables}: flat arrays for the arrays of the C code, with its index arithmetic, in a structure, for the globals; arrays of Booleans for the arrays of \texttt{unsigned char} that hold \texttt{0} and \texttt{1}.)
\item \csource{j-points.c}{196}{196}: modeled by \lname{JPoints.CurveTables}, \lname{JPoints.CurveTables.initial}
  (deviation of \lname{JPoints.CurveTables}: flat arrays for the arrays of the C code, with its index arithmetic, in a structure, for the globals; arrays of Booleans for the arrays of \texttt{unsigned char} that hold \texttt{0} and \texttt{1}.)
\item \csource{j-points.c}{200}{200}: modeled by \lname{JPoints.SieveTables}, \lname{JPoints.SieveTables.initial}
  (deviation of \lname{JPoints.SieveTables}: an array of arrays for the array of pointers \texttt{sieve\_tab}, the empty array for a null pointer; a flat array for the array \texttt{help}, with the index arithmetic of the C code; a structure, for the globals.)
\item \csource{j-points.c}{208}{208}: modeled by \lname{JPoints.Cli.Input}, \lname{JPoints.SearchCurve}
  (deviation of \lname{JPoints.Cli.Input}: a structure for the globals; an array of seven \texttt{Int}s for the seven GMP integers of \texttt{coeffs}; \texttt{Int}s for the \texttt{long}s, whose ranges after \texttt{read\_input} later lemmas give; \texttt{Bool}s for the \texttt{int} flags; a string for \texttt{print\_format}, the empty string for the null pointer that it is before \texttt{read\_input} sets it; the flags \texttt{set} of the entries of \texttt{cost\_names} as an array of their own.)
  (deviation of \lname{JPoints.SearchCurve}: an array of integers for the array of GMP integers, a natural number for the \texttt{long} \texttt{degree} and Booleans for the \texttt{int} flags, in a structure, for the globals; the initializers \texttt{0} of the flags are the values \texttt{false} from which \texttt{JPoints.normalizeCurve} starts them.)
\item \csource{j-points.c}{211}{212}: excluded (memory): multi-precision temporaries, one set per thread (init\_thread\_mpz): the array bc of the coefficients times powers of b in the lifting test, and fff, tmp, tmp2, tmp3 and ddd, the scratch integers of the input, of the set-up, of the exact check and of the output; the model computes with Lean's integers in local values
\item \csource{j-points.c}{213}{214}: modeled by \lname{JPoints.KummerCoeffs}
  (deviation of \lname{JPoints.KummerCoeffs}: integers for the GMP integers, in a structure, for the fifteen globals.)
\item \csource{j-points.c}{216}{216}: excluded (memory): multi-precision temporaries, one set per thread: x12, x22 and x32, the squares of the coordinates in the lifting test; the model passes the values
\item \csource{j-points.c}{217}{217}: modeled by \lname{JPoints.KummerEqn}
  (deviation of \lname{JPoints.KummerEqn}: a structure of three integers, for the thread-local array of three GMP integers.)
\item \csource{j-points.c}{226}{228}: modeled by \lname{JPoints.TubeCoeffs}
  (deviation of \lname{JPoints.TubeCoeffs}: an array of seven doubles for the array \texttt{fd}; a structure, for the globals, without the zeros of static storage, which nothing reads: \texttt{tube\_init} writes every field first.)
\item \csource{j-points.c}{229}{230}: modeled by \lname{JPoints.Plan}
  (deviation of \lname{JPoints.Plan}: a structure for the globals; natural numbers for the \texttt{long} indices of \texttt{pnn} and for the \texttt{int} \texttt{sieve\_mode}; a Boolean for the \texttt{int} flag \texttt{tube\_cut}; natural numbers for the \texttt{long}s \texttt{sieve\_primes1}, \texttt{sieve\_primes2} and \texttt{sieve\_primes3} after the choice, which leaves them nonnegative (before it they hold what \texttt{read\_input} leaves, \texttt{JPoints.Cli.Input}).)
\item \csource{j-points.c}{235}{236}: modeled by \lname{JPoints.TubeStats}
  (deviation of \lname{JPoints.TubeStats}: \texttt{Float}s for the doubles, which only the cost model and message 4 read, in a structure, for the globals.)
\item \csource{j-points.c}{241}{242}: modeled by \lname{JPoints.NegativeIntervals}
  (deviation of \lname{JPoints.NegativeIntervals}: a natural number for the \texttt{long} \texttt{num\_neg}, which the C code keeps nonnegative; the arrays of eight doubles start at \texttt{0.0}, as the static storage does.)
\item \csource{j-points.c}{248}{251}: modeled by \lname{JPoints.TwoAdic}, \lname{JPoints.TwoAdic.initial}
  (deviation of \lname{JPoints.TwoAdic}: flat arrays for the arrays of the C code, at \texttt{a * 64 + b}, in a structure, for the globals; a Boolean for the \texttt{int} \texttt{use2}, and an array of Booleans for the array of \texttt{unsigned char} that holds \texttt{0} and \texttt{1}; \texttt{Float}s for the doubles, which only the cost model and message 4 read.)
\item \csource{j-points.c}{252}{252}: excluded (memory): the multi-precision temporaries cpf1, cpf2, cpf3, cpfa, cpfb and cpfc of check\_one\_point\_final, one set per thread; the model computes with Lean's integers in local values
\item \csource{j-points.c}{254}{254}: modeled by \lname{JPoints.Cli.Input}, \lname{JPoints.SearchCurve}
  (deviation of \lname{JPoints.Cli.Input}: a structure for the globals; an array of seven \texttt{Int}s for the seven GMP integers of \texttt{coeffs}; \texttt{Int}s for the \texttt{long}s, whose ranges after \texttt{read\_input} later lemmas give; \texttt{Bool}s for the \texttt{int} flags; a string for \texttt{print\_format}, the empty string for the null pointer that it is before \texttt{read\_input} sets it; the flags \texttt{set} of the entries of \texttt{cost\_names} as an array of their own.)
  (deviation of \lname{JPoints.SearchCurve}: an array of integers for the array of GMP integers, a natural number for the \texttt{long} \texttt{degree} and Booleans for the \texttt{int} flags, in a structure, for the globals; the initializers \texttt{0} of the flags are the values \texttt{false} from which \texttt{JPoints.normalizeCurve} starts them.)
\item \csource{j-points.c}{258}{258}: modeled by \lname{JPoints.SearchCurve}
  (deviation of \lname{JPoints.SearchCurve}: an array of integers for the array of GMP integers, a natural number for the \texttt{long} \texttt{degree} and Booleans for the \texttt{int} flags, in a structure, for the globals; the initializers \texttt{0} of the flags are the values \texttt{false} from which \texttt{JPoints.normalizeCurve} starts them.)
\item \csource{j-points.c}{259}{259}: modeled by \lname{JPoints.CurveTables}, \lname{JPoints.CurveTables.initial}
  (deviation of \lname{JPoints.CurveTables}: flat arrays for the arrays of the C code, with its index arithmetic, in a structure, for the globals; arrays of Booleans for the arrays of \texttt{unsigned char} that hold \texttt{0} and \texttt{1}.)
\item \csource{j-points.c}{262}{262}: modeled by \lname{JPoints.Rates}
  (deviation of \lname{JPoints.Rates}: the global and the local array (\texttt{j-{}points.c:809}) in a structure; an array of Booleans for the array of \texttt{int} flags.)
\item \csource{j-points.c}{265}{269}: modeled by \lname{JPoints.Cli.Input}
  (deviation of \lname{JPoints.Cli.Input}: a structure for the globals; an array of seven \texttt{Int}s for the seven GMP integers of \texttt{coeffs}; \texttt{Int}s for the \texttt{long}s, whose ranges after \texttt{read\_input} later lemmas give; \texttt{Bool}s for the \texttt{int} flags; a string for \texttt{print\_format}, the empty string for the null pointer that it is before \texttt{read\_input} sets it; the flags \texttt{set} of the entries of \texttt{cost\_names} as an array of their own.)
\item \csource{j-points.c}{270}{272}: modeled by \lname{JPoints.Cli.Input}, \lname{JPoints.Plan}
  (deviation of \lname{JPoints.Cli.Input}: a structure for the globals; an array of seven \texttt{Int}s for the seven GMP integers of \texttt{coeffs}; \texttt{Int}s for the \texttt{long}s, whose ranges after \texttt{read\_input} later lemmas give; \texttt{Bool}s for the \texttt{int} flags; a string for \texttt{print\_format}, the empty string for the null pointer that it is before \texttt{read\_input} sets it; the flags \texttt{set} of the entries of \texttt{cost\_names} as an array of their own.)
  (deviation of \lname{JPoints.Plan}: a structure for the globals; natural numbers for the \texttt{long} indices of \texttt{pnn} and for the \texttt{int} \texttt{sieve\_mode}; a Boolean for the \texttt{int} flag \texttt{tube\_cut}; natural numbers for the \texttt{long}s \texttt{sieve\_primes1}, \texttt{sieve\_primes2} and \texttt{sieve\_primes3} after the choice, which leaves them nonnegative (before it they hold what \texttt{read\_input} leaves, \texttt{JPoints.Cli.Input}).)
\item \csource{j-points.c}{273}{277}: modeled by \lname{JPoints.Cli.Input}
  (deviation of \lname{JPoints.Cli.Input}: a structure for the globals; an array of seven \texttt{Int}s for the seven GMP integers of \texttt{coeffs}; \texttt{Int}s for the \texttt{long}s, whose ranges after \texttt{read\_input} later lemmas give; \texttt{Bool}s for the \texttt{int} flags; a string for \texttt{print\_format}, the empty string for the null pointer that it is before \texttt{read\_input} sets it; the flags \texttt{set} of the entries of \texttt{cost\_names} as an array of their own.)
\item \csource{j-points.c}{279}{279}: modeled by \lname{JPoints.SiftState}, \lname{JPoints.SiftState.initial}
  (deviation of \lname{JPoints.SiftState}: an array for the bit array that \texttt{survivors} points to; indices into the sieve tables for the pointers \texttt{rowptr[n]} into \texttt{sieve\_tab[n]}; natural numbers for the counters, which start at \texttt{0} and count up, and for the index \texttt{cur\_chunk}; a structure, for the thread-local globals of the one thread.)
  (deviation of \lname{JPoints.SiftState.initial}: a fresh array of \texttt{arraySize + VW} zeros, \texttt{VW = 1}, for the bit array that \texttt{init\_thread\_sieve} allocates (memory management, not modeled), every word of which the sieve writes before it reads it; index \texttt{0} for the null pointers of \texttt{rowptr}, which the sieve reads only after \texttt{row\_setup} has set them.)
\item \csource{j-points.c}{282}{284}: modeled by \lname{JPoints.SiftState}, \lname{JPoints.SiftState.initial}
  (deviation of \lname{JPoints.SiftState}: an array for the bit array that \texttt{survivors} points to; indices into the sieve tables for the pointers \texttt{rowptr[n]} into \texttt{sieve\_tab[n]}; natural numbers for the counters, which start at \texttt{0} and count up, and for the index \texttt{cur\_chunk}; a structure, for the thread-local globals of the one thread.)
  (deviation of \lname{JPoints.SiftState.initial}: a fresh array of \texttt{arraySize + VW} zeros, \texttt{VW = 1}, for the bit array that \texttt{init\_thread\_sieve} allocates (memory management, not modeled), every word of which the sieve writes before it reads it; index \texttt{0} for the null pointers of \texttt{rowptr}, which the sieve reads only after \texttt{row\_setup} has set them.)
\item \csource{j-points.c}{285}{285}: modeled by \lname{JPoints.Output}
  (deviation of \lname{JPoints.Output}: the text collected, for the output on \texttt{stdout}, which the program writes; a natural number for the thread-local \texttt{long} \texttt{total} of the one thread, which starts at \texttt{0} and counts up.)
\item \csource{j-points.c}{286}{286}: modeled by \lname{JPoints.Totals}
  (deviation of \lname{JPoints.Totals}: a structure for the globals; natural numbers for the \texttt{long} sums of the counters, which count up from \texttt{0}.)
\item \csource{j-points.c}{292}{292}: modeled by \lname{JPoints.Cli.Input}
  (deviation of \lname{JPoints.Cli.Input}: a structure for the globals; an array of seven \texttt{Int}s for the seven GMP integers of \texttt{coeffs}; \texttt{Int}s for the \texttt{long}s, whose ranges after \texttt{read\_input} later lemmas give; \texttt{Bool}s for the \texttt{int} flags; a string for \texttt{print\_format}, the empty string for the null pointer that it is before \texttt{read\_input} sets it; the flags \texttt{set} of the entries of \texttt{cost\_names} as an array of their own.)
\item \csource{j-points.c}{293}{299}: excluded (threads): the state of the threads of -t: the lock of the output, the counters of the units of work handed out and printed, the list of the finished units waiting to be printed, and the buffer of the unit in work of each thread; the model runs the units in turn on one thread, and the threads print the points of the units in the order of the units, which is that output when the threads all start
\item \csource{j-points.c}{301}{301}: modeled by \lname{JPoints.SiftGlobals}
  (deviation of \lname{JPoints.SiftGlobals}: a structure, for the globals, of which only the declaration of \texttt{begmask} and \texttt{endmask} is recorded here, the others where the model defines them (\texttt{JPoints.Tables}, \texttt{JPoints.CurveTables}, \texttt{JPoints.SieveTables}, \texttt{JPoints.SearchCurve}, \texttt{JPoints.KummerCoeffs}, \texttt{JPoints.TwoAdic}, \texttt{JPoints.Plan}, \texttt{JPoints.Cli.Input}); \texttt{UInt64}s for the masks; natural numbers for the nonnegative \texttt{long}s \texttt{sieve\_primes1}, \texttt{sieve\_primes2}, \texttt{sieve\_primes3} and the entries of \texttt{pnn}, and for the \texttt{long} \texttt{array\_size}, which \texttt{main} converts from the kilobytes of \texttt{-{}s} into words.)
\item \csource{j-points.c}{308}{331}: excluded (other): the prototypes of the functions of j-points.c; the model has no forward declarations
\item \csource{j-points.c}{337}{338}: modeled by \lname{JPoints.Program.run}
  (deviation of \lname{JPoints.Program.run}: the text returned, which \texttt{JPoints.Program.main} writes at the end, where \texttt{j-{}points} prints as it goes; the exit through \texttt{error} as the text and the status returned; the words after the name of the program for \texttt{argv}, \texttt{args.length + 1} for \texttt{argc}, and the comparisons of \texttt{strcmp} as the equality of strings; the refusal of the formats of \texttt{-{}f} that the model does not admit, a decision of the project, after the command line is read and before any message (the formats that it admits are those that \texttt{printf} and \texttt{gmp\_printf} print as the model does); the words of the command line as the Lean runtime makes strings of them, which replaces the bytes that are not valid UTF-8 by the character U+FFFD, where \texttt{j-{}points} reads the bytes as they are.)
\item \csource{j-points.c}{339}{339}: modeled by \lname{JPoints.Program.mainArraySize}
  (deviation of \lname{JPoints.Program.mainArraySize}: integers for the \texttt{long}s; \texttt{JPoints.ceilDiv} for the macro \texttt{CEIL}; \texttt{>>>} for C's \texttt{>>} on the positive \texttt{s}, and the product with \texttt{2 \textasciicircum{} (13 -{} LONG\_SHIFT)} for the shift \texttt{<<} of the positive \texttt{array\_size}, which does not overflow when \texttt{array\_size ≤ s >> (13 -{} LONG\_SHIFT)}.)
\item \csource{j-points.c}{341}{341}: excluded (other): the test for a platform whose unsigned long has neither 16, 32 nor 64 bits (LONG\_SHIFT is 0 then), with error 1; on the LP64 platforms of the model LONG\_SHIFT is 6 and the test never fires
\item \csource{j-points.c}{342}{342}: modeled by \lname{JPoints.Program.run}
  (deviation of \lname{JPoints.Program.run}: the text returned, which \texttt{JPoints.Program.main} writes at the end, where \texttt{j-{}points} prints as it goes; the exit through \texttt{error} as the text and the status returned; the words after the name of the program for \texttt{argv}, \texttt{args.length + 1} for \texttt{argc}, and the comparisons of \texttt{strcmp} as the equality of strings; the refusal of the formats of \texttt{-{}f} that the model does not admit, a decision of the project, after the command line is read and before any message (the formats that it admits are those that \texttt{printf} and \texttt{gmp\_printf} print as the model does); the words of the command line as the Lean runtime makes strings of them, which replaces the bytes that are not valid UTF-8 by the character U+FFFD, where \texttt{j-{}points} reads the bytes as they are.)
\item \csource{j-points.c}{345}{345}: modeled by \lname{JPoints.Program.run}
  (deviation of \lname{JPoints.Program.run}: the text returned, which \texttt{JPoints.Program.main} writes at the end, where \texttt{j-{}points} prints as it goes; the exit through \texttt{error} as the text and the status returned; the words after the name of the program for \texttt{argv}, \texttt{args.length + 1} for \texttt{argc}, and the comparisons of \texttt{strcmp} as the equality of strings; the refusal of the formats of \texttt{-{}f} that the model does not admit, a decision of the project, after the command line is read and before any message (the formats that it admits are those that \texttt{printf} and \texttt{gmp\_printf} print as the model does); the words of the command line as the Lean runtime makes strings of them, which replaces the bytes that are not valid UTF-8 by the character U+FFFD, where \texttt{j-{}points} reads the bytes as they are.)
\item \csource{j-points.c}{346}{348}: modeled by \lname{JPoints.Program.mainCostList}
  (deviation of \lname{JPoints.Program.mainCostList}: the text returned, for the output on \texttt{stdout}; the constants in force passed, for the statics; a natural number for the \texttt{size\_t} index; the values by \texttt{JPoints.Double.ofFloat}, which takes a NaN to \texttt{0x7ff8000000000000} and would print \texttt{nan} for \texttt{-{}nan} (the defaults of the build are positive numbers).)
\item \csource{j-points.c}{349}{350}: modeled by \lname{JPoints.Program.run}
  (deviation of \lname{JPoints.Program.run}: the text returned, which \texttt{JPoints.Program.main} writes at the end, where \texttt{j-{}points} prints as it goes; the exit through \texttt{error} as the text and the status returned; the words after the name of the program for \texttt{argv}, \texttt{args.length + 1} for \texttt{argc}, and the comparisons of \texttt{strcmp} as the equality of strings; the refusal of the formats of \texttt{-{}f} that the model does not admit, a decision of the project, after the command line is read and before any message (the formats that it admits are those that \texttt{printf} and \texttt{gmp\_printf} print as the model does); the words of the command line as the Lean runtime makes strings of them, which replaces the bytes that are not valid UTF-8 by the character U+FFFD, where \texttt{j-{}points} reads the bytes as they are.)
\item \csource{j-points.c}{352}{353}: modeled by \lname{JPoints.Program.run}
  (deviation of \lname{JPoints.Program.run}: the text returned, which \texttt{JPoints.Program.main} writes at the end, where \texttt{j-{}points} prints as it goes; the exit through \texttt{error} as the text and the status returned; the words after the name of the program for \texttt{argv}, \texttt{args.length + 1} for \texttt{argc}, and the comparisons of \texttt{strcmp} as the equality of strings; the refusal of the formats of \texttt{-{}f} that the model does not admit, a decision of the project, after the command line is read and before any message (the formats that it admits are those that \texttt{printf} and \texttt{gmp\_printf} print as the model does); the words of the command line as the Lean runtime makes strings of them, which replaces the bytes that are not valid UTF-8 by the character U+FFFD, where \texttt{j-{}points} reads the bytes as they are.)
\item \csource{j-points.c}{354}{359}: modeled by \lname{JPoints.Program.mainHead}
  (deviation of \lname{JPoints.Program.mainHead}: the text returned, for the output on \texttt{stdout} that \texttt{fflush} flushes; the globals of \texttt{message} gathered from \texttt{input}, with the curve given as the curve of \texttt{JPoints.ChoiceGlobals} (the global \texttt{coeffs} before \texttt{normalise\_curve}, which message 5 prints), and the initial values of the statistics of the analysis, of \texttt{num\_neg} and of the condition at 2, which no message here reads.)
\item \csource{j-points.c}{360}{360}: modeled by \lname{JPoints.Program.mainSearch}
  (deviation of \lname{JPoints.Program.mainSearch}: the text returned, for the output on \texttt{stdout}; the globals computed in the order of the C code and passed, for the globals that the functions read and write, the plan, the static data of the sieve and the state of the analysis from their initial values; the words of the bit array in \texttt{JPoints.SearchGlobals.arraySize}, a natural number, for the global \texttt{array\_size} that \texttt{main} overwrites with them, while \texttt{JPoints.Cli.Input.arraySize} keeps the kilobytes, which nothing reads after the conversion; the globals of \texttt{message} for messages 12 and 2 with the totals, which they print, and the initial values of the plan and of the records, which they do not read.)
\item \csource{j-points.c}{361}{362}: modeled by \lname{JPoints.Program.mainMasks}
  (deviation of \lname{JPoints.Program.mainMasks}: \texttt{UInt64}s for the \texttt{unsigned long}s; the integers \texttt{(-{}H) \% 64} and \texttt{(-{}H -{} 1) \% 64}, in \texttt{[0, 64)}, for \texttt{(-{}height) \& LONG\_MASK} and \texttt{(\textasciitilde{}height) \& LONG\_MASK} on the \texttt{long} \texttt{height}, the same in two's complement.)
\item \csource{j-points.c}{363}{363}: modeled by \lname{JPoints.Program.mainArraySize}
  (deviation of \lname{JPoints.Program.mainArraySize}: integers for the \texttt{long}s; \texttt{JPoints.ceilDiv} for the macro \texttt{CEIL}; \texttt{>>>} for C's \texttt{>>} on the positive \texttt{s}, and the product with \texttt{2 \textasciicircum{} (13 -{} LONG\_SHIFT)} for the shift \texttt{<<} of the positive \texttt{array\_size}, which does not overflow when \texttt{array\_size ≤ s >> (13 -{} LONG\_SHIFT)}.)
\item \csource{j-points.c}{366}{366}: modeled by \lname{JPoints.Program.mainArraySize}
  (deviation of \lname{JPoints.Program.mainArraySize}: integers for the \texttt{long}s; \texttt{JPoints.ceilDiv} for the macro \texttt{CEIL}; \texttt{>>>} for C's \texttt{>>} on the positive \texttt{s}, and the product with \texttt{2 \textasciicircum{} (13 -{} LONG\_SHIFT)} for the shift \texttt{<<} of the positive \texttt{array\_size}, which does not overflow when \texttt{array\_size ≤ s >> (13 -{} LONG\_SHIFT)}.)
\item \csource{j-points.c}{368}{369}: modeled by \lname{JPoints.Program.mainSearch}
  (deviation of \lname{JPoints.Program.mainSearch}: the text returned, for the output on \texttt{stdout}; the globals computed in the order of the C code and passed, for the globals that the functions read and write, the plan, the static data of the sieve and the state of the analysis from their initial values; the words of the bit array in \texttt{JPoints.SearchGlobals.arraySize}, a natural number, for the global \texttt{array\_size} that \texttt{main} overwrites with them, while \texttt{JPoints.Cli.Input.arraySize} keeps the kilobytes, which nothing reads after the conversion; the globals of \texttt{message} for messages 12 and 2 with the totals, which they print, and the initial values of the plan and of the records, which they do not read.)
\item \csource{j-points.c}{371}{371}: modeled by \lname{JPoints.Program.mainSearch}
  (deviation of \lname{JPoints.Program.mainSearch}: the text returned, for the output on \texttt{stdout}; the globals computed in the order of the C code and passed, for the globals that the functions read and write, the plan, the static data of the sieve and the state of the analysis from their initial values; the words of the bit array in \texttt{JPoints.SearchGlobals.arraySize}, a natural number, for the global \texttt{array\_size} that \texttt{main} overwrites with them, while \texttt{JPoints.Cli.Input.arraySize} keeps the kilobytes, which nothing reads after the conversion; the globals of \texttt{message} for messages 12 and 2 with the totals, which they print, and the initial values of the plan and of the records, which they do not read.)
\item \csource{j-points.c}{373}{374}: modeled by \lname{JPoints.Program.mainSearch}
  (deviation of \lname{JPoints.Program.mainSearch}: the text returned, for the output on \texttt{stdout}; the globals computed in the order of the C code and passed, for the globals that the functions read and write, the plan, the static data of the sieve and the state of the analysis from their initial values; the words of the bit array in \texttt{JPoints.SearchGlobals.arraySize}, a natural number, for the global \texttt{array\_size} that \texttt{main} overwrites with them, while \texttt{JPoints.Cli.Input.arraySize} keeps the kilobytes, which nothing reads after the conversion; the globals of \texttt{message} for messages 12 and 2 with the totals, which they print, and the initial values of the plan and of the records, which they do not read.)
\item \csource{j-points.c}{375}{376}: modeled by \lname{JPoints.Program.run}
  (deviation of \lname{JPoints.Program.run}: the text returned, which \texttt{JPoints.Program.main} writes at the end, where \texttt{j-{}points} prints as it goes; the exit through \texttt{error} as the text and the status returned; the words after the name of the program for \texttt{argv}, \texttt{args.length + 1} for \texttt{argc}, and the comparisons of \texttt{strcmp} as the equality of strings; the refusal of the formats of \texttt{-{}f} that the model does not admit, a decision of the project, after the command line is read and before any message (the formats that it admits are those that \texttt{printf} and \texttt{gmp\_printf} print as the model does); the words of the command line as the Lean runtime makes strings of them, which replaces the bytes that are not valid UTF-8 by the character U+FFFD, where \texttt{j-{}points} reads the bytes as they are.)
\item \csource{j-points.c}{391}{397}: modeled by \lname{JPoints.Cli.setCosts}
  (deviation of \lname{JPoints.Cli.setCosts}: some globals for the return value 1, and none for 0, on which the caller exits, so that the values set by the items before a refused one are never read; the rest of the string from the pointer for the pointer \texttt{s}, the characters before the \texttt{=} for \texttt{eq}, the number of characters that \texttt{strtod} reads for \texttt{end}, and the flag of \texttt{JPoints.Cli.strtod} for \texttt{errno}, which \texttt{set\_costs} clears before the call; the return of 1 at the end of \texttt{s} as the exit from the loop, after which the globals are returned; the double stored as a \texttt{Float} (\texttt{JPoints.Double.toFloat}).)
\item \csource{j-points.c}{398}{400}: modeled by \lname{JPoints.Cli.costIndex}
  (deviation of \lname{JPoints.Cli.costIndex}: the lists of characters compared for \texttt{strlen} and \texttt{strncmp} on the bytes, which agree as the names are ASCII; the index returned, a natural number, for the \texttt{size\_t} \texttt{k} that the loop leaves.)
\item \csource{j-points.c}{401}{413}: modeled by \lname{JPoints.Cli.setCosts}
  (deviation of \lname{JPoints.Cli.setCosts}: some globals for the return value 1, and none for 0, on which the caller exits, so that the values set by the items before a refused one are never read; the rest of the string from the pointer for the pointer \texttt{s}, the characters before the \texttt{=} for \texttt{eq}, the number of characters that \texttt{strtod} reads for \texttt{end}, and the flag of \texttt{JPoints.Cli.strtod} for \texttt{errno}, which \texttt{set\_costs} clears before the call; the return of 1 at the end of \texttt{s} as the exit from the loop, after which the globals are returned; the double stored as a \texttt{Float} (\texttt{JPoints.Double.toFloat}).)
\item \csource{j-points.c}{415}{416}: modeled by \lname{JPoints.Cli.readInput}
  (deviation of \lname{JPoints.Cli.readInput}: the globals returned, or the error, for the exit through \texttt{error}; the words after the name of the program for \texttt{argv}, \texttt{args[0]} for \texttt{argv[1]} and \texttt{args[1]} for \texttt{argv[2]}, which \texttt{main} makes sure exist.)
\item \csource{j-points.c}{417}{424}: modeled by \lname{JPoints.Cli.readCoefficients}
  (deviation of \lname{JPoints.Cli.readCoefficients}: the global \texttt{coeffs} as an argument, returned with \texttt{degree}, or the error, for the exit through \texttt{error(3)}; the rest of the string from the pointer for the pointer \texttt{s}, with none for \texttt{NULL}; the eighth integer, which \texttt{scan\_mpz} writes into the scratch \texttt{fff}, not kept; an \texttt{Int} for the \texttt{long} \texttt{degree}.)
\item \csource{j-points.c}{425}{438}: modeled by \lname{JPoints.Cli.readInput}
  (deviation of \lname{JPoints.Cli.readInput}: the globals returned, or the error, for the exit through \texttt{error}; the words after the name of the program for \texttt{argv}, \texttt{args[0]} for \texttt{argv[1]} and \texttt{args[1]} for \texttt{argv[2]}, which \texttt{main} makes sure exist.)
\item \csource{j-points.c}{441}{443}: modeled by \lname{JPoints.Cli.readOptions}
  (deviation of \lname{JPoints.Cli.readOptions}: the globals returned, or the error, for the exit through \texttt{error}; the words after the name of the program for \texttt{argv} and \texttt{args.size} for \texttt{argc}; a natural number for the \texttt{long} \texttt{i}.)
\item \csource{j-points.c}{444}{448}: modeled by \lname{JPoints.Cli.readOption}
  (deviation of \lname{JPoints.Cli.readOption}: the globals and the index returned, or the error, for the exit through \texttt{error}; the words after the name of the program for \texttt{argv}, \texttt{args[i -{} 1]} for \texttt{argv[i]} and \texttt{args.size} for \texttt{argc}; the number of characters of the word for \texttt{strlen}, its number of bytes, which differ only for a word with a character that is not ASCII, refused either way, as neither the minus sign nor an option is such a character; the character \texttt{argv[i][1]} of the \texttt{switch} as \texttt{w[1]?}, whose \texttt{none}, which the test of the length excludes, is error 6 too; the test for the missing value, which each case of an option with a value makes, made once before them.)
\item \csource{j-points.c}{449}{450}: modeled by \lname{JPoints.Cli.setOption}
  (deviation of \lname{JPoints.Cli.setOption}: the globals returned, or the error, for the exit through \texttt{error}, so that a value refused after \texttt{scan\_long} has written it is not kept; \texttt{scan\_long} called once before the cases of the options that read a \texttt{long}, where each of these cases calls it first; the long that \texttt{scan\_long} writes when it fails, on which the caller exits, not kept.)
\item \csource{j-points.c}{451}{455}: modeled by \lname{JPoints.Cli.readOption}
  (deviation of \lname{JPoints.Cli.readOption}: the globals and the index returned, or the error, for the exit through \texttt{error}; the words after the name of the program for \texttt{argv}, \texttt{args[i -{} 1]} for \texttt{argv[i]} and \texttt{args.size} for \texttt{argc}; the number of characters of the word for \texttt{strlen}, its number of bytes, which differ only for a word with a character that is not ASCII, refused either way, as neither the minus sign nor an option is such a character; the character \texttt{argv[i][1]} of the \texttt{switch} as \texttt{w[1]?}, whose \texttt{none}, which the test of the length excludes, is error 6 too; the test for the missing value, which each case of an option with a value makes, made once before them.)
\item \csource{j-points.c}{456}{457}: modeled by \lname{JPoints.Cli.setOption}
  (deviation of \lname{JPoints.Cli.setOption}: the globals returned, or the error, for the exit through \texttt{error}, so that a value refused after \texttt{scan\_long} has written it is not kept; \texttt{scan\_long} called once before the cases of the options that read a \texttt{long}, where each of these cases calls it first; the long that \texttt{scan\_long} writes when it fails, on which the caller exits, not kept.)
\item \csource{j-points.c}{458}{462}: modeled by \lname{JPoints.Cli.readOption}
  (deviation of \lname{JPoints.Cli.readOption}: the globals and the index returned, or the error, for the exit through \texttt{error}; the words after the name of the program for \texttt{argv}, \texttt{args[i -{} 1]} for \texttt{argv[i]} and \texttt{args.size} for \texttt{argc}; the number of characters of the word for \texttt{strlen}, its number of bytes, which differ only for a word with a character that is not ASCII, refused either way, as neither the minus sign nor an option is such a character; the character \texttt{argv[i][1]} of the \texttt{switch} as \texttt{w[1]?}, whose \texttt{none}, which the test of the length excludes, is error 6 too; the test for the missing value, which each case of an option with a value makes, made once before them.)
\item \csource{j-points.c}{463}{464}: modeled by \lname{JPoints.Cli.setOption}
  (deviation of \lname{JPoints.Cli.setOption}: the globals returned, or the error, for the exit through \texttt{error}, so that a value refused after \texttt{scan\_long} has written it is not kept; \texttt{scan\_long} called once before the cases of the options that read a \texttt{long}, where each of these cases calls it first; the long that \texttt{scan\_long} writes when it fails, on which the caller exits, not kept.)
\item \csource{j-points.c}{465}{469}: modeled by \lname{JPoints.Cli.readOption}
  (deviation of \lname{JPoints.Cli.readOption}: the globals and the index returned, or the error, for the exit through \texttt{error}; the words after the name of the program for \texttt{argv}, \texttt{args[i -{} 1]} for \texttt{argv[i]} and \texttt{args.size} for \texttt{argc}; the number of characters of the word for \texttt{strlen}, its number of bytes, which differ only for a word with a character that is not ASCII, refused either way, as neither the minus sign nor an option is such a character; the character \texttt{argv[i][1]} of the \texttt{switch} as \texttt{w[1]?}, whose \texttt{none}, which the test of the length excludes, is error 6 too; the test for the missing value, which each case of an option with a value makes, made once before them.)
\item \csource{j-points.c}{470}{471}: modeled by \lname{JPoints.Cli.setOption}
  (deviation of \lname{JPoints.Cli.setOption}: the globals returned, or the error, for the exit through \texttt{error}, so that a value refused after \texttt{scan\_long} has written it is not kept; \texttt{scan\_long} called once before the cases of the options that read a \texttt{long}, where each of these cases calls it first; the long that \texttt{scan\_long} writes when it fails, on which the caller exits, not kept.)
\item \csource{j-points.c}{472}{476}: modeled by \lname{JPoints.Cli.readOption}
  (deviation of \lname{JPoints.Cli.readOption}: the globals and the index returned, or the error, for the exit through \texttt{error}; the words after the name of the program for \texttt{argv}, \texttt{args[i -{} 1]} for \texttt{argv[i]} and \texttt{args.size} for \texttt{argc}; the number of characters of the word for \texttt{strlen}, its number of bytes, which differ only for a word with a character that is not ASCII, refused either way, as neither the minus sign nor an option is such a character; the character \texttt{argv[i][1]} of the \texttt{switch} as \texttt{w[1]?}, whose \texttt{none}, which the test of the length excludes, is error 6 too; the test for the missing value, which each case of an option with a value makes, made once before them.)
\item \csource{j-points.c}{477}{477}: modeled by \lname{JPoints.Cli.setOption}
  (deviation of \lname{JPoints.Cli.setOption}: the globals returned, or the error, for the exit through \texttt{error}, so that a value refused after \texttt{scan\_long} has written it is not kept; \texttt{scan\_long} called once before the cases of the options that read a \texttt{long}, where each of these cases calls it first; the long that \texttt{scan\_long} writes when it fails, on which the caller exits, not kept.)
\item \csource{j-points.c}{478}{482}: modeled by \lname{JPoints.Cli.readOption}
  (deviation of \lname{JPoints.Cli.readOption}: the globals and the index returned, or the error, for the exit through \texttt{error}; the words after the name of the program for \texttt{argv}, \texttt{args[i -{} 1]} for \texttt{argv[i]} and \texttt{args.size} for \texttt{argc}; the number of characters of the word for \texttt{strlen}, its number of bytes, which differ only for a word with a character that is not ASCII, refused either way, as neither the minus sign nor an option is such a character; the character \texttt{argv[i][1]} of the \texttt{switch} as \texttt{w[1]?}, whose \texttt{none}, which the test of the length excludes, is error 6 too; the test for the missing value, which each case of an option with a value makes, made once before them.)
\item \csource{j-points.c}{483}{484}: modeled by \lname{JPoints.Cli.setOption}
  (deviation of \lname{JPoints.Cli.setOption}: the globals returned, or the error, for the exit through \texttt{error}, so that a value refused after \texttt{scan\_long} has written it is not kept; \texttt{scan\_long} called once before the cases of the options that read a \texttt{long}, where each of these cases calls it first; the long that \texttt{scan\_long} writes when it fails, on which the caller exits, not kept.)
\item \csource{j-points.c}{485}{501}: modeled by \lname{JPoints.Cli.readOption}
  (deviation of \lname{JPoints.Cli.readOption}: the globals and the index returned, or the error, for the exit through \texttt{error}; the words after the name of the program for \texttt{argv}, \texttt{args[i -{} 1]} for \texttt{argv[i]} and \texttt{args.size} for \texttt{argc}; the number of characters of the word for \texttt{strlen}, its number of bytes, which differ only for a word with a character that is not ASCII, refused either way, as neither the minus sign nor an option is such a character; the character \texttt{argv[i][1]} of the \texttt{switch} as \texttt{w[1]?}, whose \texttt{none}, which the test of the length excludes, is error 6 too; the test for the missing value, which each case of an option with a value makes, made once before them.)
\item \csource{j-points.c}{502}{503}: modeled by \lname{JPoints.Cli.setOption}
  (deviation of \lname{JPoints.Cli.setOption}: the globals returned, or the error, for the exit through \texttt{error}, so that a value refused after \texttt{scan\_long} has written it is not kept; \texttt{scan\_long} called once before the cases of the options that read a \texttt{long}, where each of these cases calls it first; the long that \texttt{scan\_long} writes when it fails, on which the caller exits, not kept.)
\item \csource{j-points.c}{504}{509}: modeled by \lname{JPoints.Cli.readOption}
  (deviation of \lname{JPoints.Cli.readOption}: the globals and the index returned, or the error, for the exit through \texttt{error}; the words after the name of the program for \texttt{argv}, \texttt{args[i -{} 1]} for \texttt{argv[i]} and \texttt{args.size} for \texttt{argc}; the number of characters of the word for \texttt{strlen}, its number of bytes, which differ only for a word with a character that is not ASCII, refused either way, as neither the minus sign nor an option is such a character; the character \texttt{argv[i][1]} of the \texttt{switch} as \texttt{w[1]?}, whose \texttt{none}, which the test of the length excludes, is error 6 too; the test for the missing value, which each case of an option with a value makes, made once before them.)
\item \csource{j-points.c}{510}{510}: modeled by \lname{JPoints.Cli.setOption}
  (deviation of \lname{JPoints.Cli.setOption}: the globals returned, or the error, for the exit through \texttt{error}, so that a value refused after \texttt{scan\_long} has written it is not kept; \texttt{scan\_long} called once before the cases of the options that read a \texttt{long}, where each of these cases calls it first; the long that \texttt{scan\_long} writes when it fails, on which the caller exits, not kept.)
\item \csource{j-points.c}{511}{515}: modeled by \lname{JPoints.Cli.readOption}
  (deviation of \lname{JPoints.Cli.readOption}: the globals and the index returned, or the error, for the exit through \texttt{error}; the words after the name of the program for \texttt{argv}, \texttt{args[i -{} 1]} for \texttt{argv[i]} and \texttt{args.size} for \texttt{argc}; the number of characters of the word for \texttt{strlen}, its number of bytes, which differ only for a word with a character that is not ASCII, refused either way, as neither the minus sign nor an option is such a character; the character \texttt{argv[i][1]} of the \texttt{switch} as \texttt{w[1]?}, whose \texttt{none}, which the test of the length excludes, is error 6 too; the test for the missing value, which each case of an option with a value makes, made once before them.)
\item \csource{j-points.c}{516}{518}: modeled by \lname{JPoints.Cli.setOption}
  (deviation of \lname{JPoints.Cli.setOption}: the globals returned, or the error, for the exit through \texttt{error}, so that a value refused after \texttt{scan\_long} has written it is not kept; \texttt{scan\_long} called once before the cases of the options that read a \texttt{long}, where each of these cases calls it first; the long that \texttt{scan\_long} writes when it fails, on which the caller exits, not kept.)
\item \csource{j-points.c}{519}{521}: modeled by \lname{JPoints.Cli.readOption}
  (deviation of \lname{JPoints.Cli.readOption}: the globals and the index returned, or the error, for the exit through \texttt{error}; the words after the name of the program for \texttt{argv}, \texttt{args[i -{} 1]} for \texttt{argv[i]} and \texttt{args.size} for \texttt{argc}; the number of characters of the word for \texttt{strlen}, its number of bytes, which differ only for a word with a character that is not ASCII, refused either way, as neither the minus sign nor an option is such a character; the character \texttt{argv[i][1]} of the \texttt{switch} as \texttt{w[1]?}, whose \texttt{none}, which the test of the length excludes, is error 6 too; the test for the missing value, which each case of an option with a value makes, made once before them.)
\item \csource{j-points.c}{522}{522}: modeled by \lname{JPoints.Cli.readOptions}
  (deviation of \lname{JPoints.Cli.readOptions}: the globals returned, or the error, for the exit through \texttt{error}; the words after the name of the program for \texttt{argv} and \texttt{args.size} for \texttt{argc}; a natural number for the \texttt{long} \texttt{i}.)
\item \csource{j-points.c}{525}{525}: modeled by \lname{JPoints.Cli.settleOptions}
  (deviation of \lname{JPoints.Cli.settleOptions}: the globals taken and returned.)
\item \csource{j-points.c}{529}{537}: modeled by \lname{JPoints.Cli.settleOptions}
  (deviation of \lname{JPoints.Cli.settleOptions}: the globals taken and returned.)
\item \csource{j-points.c}{540}{541}: modeled by \lname{JPoints.Cli.settleOptions}
  (deviation of \lname{JPoints.Cli.settleOptions}: the globals taken and returned.)
\item \csource{j-points.c}{542}{542}: modeled by \lname{JPoints.Cli.readInput}
  (deviation of \lname{JPoints.Cli.readInput}: the globals returned, or the error, for the exit through \texttt{error}; the words after the name of the program for \texttt{argv}, \texttt{args[0]} for \texttt{argv[1]} and \texttt{args[1]} for \texttt{argv[2]}, which \texttt{main} makes sure exist.)
\item \csource{j-points.c}{548}{554}: modeled by \lname{JPoints.Cli.scanLong}
  (deviation of \lname{JPoints.Cli.scanLong}: \texttt{some v} for the return value \texttt{1} with the \texttt{long} \texttt{v} written through the pointer, and \texttt{none} for the return value \texttt{0}, on which every caller exits, so that the \texttt{long} written then is never read; the flag of \texttt{JPoints.Cli.strtol} for \texttt{errno}, which \texttt{scan\_long} clears before the call; an \texttt{Int} for the \texttt{long}.)
\item \csource{j-points.c}{560}{578}: modeled by \lname{JPoints.Cli.scanMpz}
  (deviation of \lname{JPoints.Cli.scanMpz}: the rest of the string from the pointer for the pointer, with none for \texttt{NULL}; the integer returned with the pointer, for the \texttt{mpz\_t} that \texttt{scan\_mpz} writes through \texttt{x} when it does not return \texttt{NULL}; the error returned, for the exit through \texttt{error(3)}; a local integer for \texttt{tmp2}, a global GMP integer of the C (one per thread) that \texttt{scan\_mpz} sets before it reads it and that no function reads before setting it; a \texttt{Bool} for the \texttt{long} flag \texttt{neg}.)
\item \csource{j-points.c}{593}{604}: modeled by \lname{JPoints.normalizeCurve}
  (deviation of \lname{JPoints.normalizeCurve}: the curve and the flags returned, for the globals that the C code writes, the flags starting from \texttt{false}, the values \texttt{0} of their initializers; a natural number for the \texttt{long} \texttt{degree}; a Boolean for the \texttt{int} \texttt{quiet}, and an argument for this global; the text of the message returned, which the caller writes to \texttt{stdout}.)
\item \csource{j-points.c}{605}{611}: modeled by \lname{JPoints.normalizeCurve}
  (deviation of \lname{JPoints.normalizeCurve}: the curve and the flags returned, for the globals that the C code writes, the flags starting from \texttt{false}, the values \texttt{0} of their initializers; a natural number for the \texttt{long} \texttt{degree}; a Boolean for the \texttt{int} \texttt{quiet}, and an argument for this global; the text of the message returned, which the caller writes to \texttt{stdout}.)
\item \csource{j-points.c}{612}{613}: modeled by \lname{JPoints.normalizeCurve}
  (deviation of \lname{JPoints.normalizeCurve}: the curve and the flags returned, for the globals that the C code writes, the flags starting from \texttt{false}, the values \texttt{0} of their initializers; a natural number for the \texttt{long} \texttt{degree}; a Boolean for the \texttt{int} \texttt{quiet}, and an argument for this global; the text of the message returned, which the caller writes to \texttt{stdout}.)
\item \csource{j-points.c}{618}{622}: modeled by \lname{JPoints.Cli.squarefree}
  (deviation of \lname{JPoints.Cli.squarefree}: the globals \texttt{coeffs} and \texttt{degree} as arguments; arrays of seven \texttt{Rat}s for the arrays \texttt{u} and \texttt{v} of \texttt{mpq\_t}s, zero at the start as \texttt{mpq\_init} sets them (lines 623-624, memory), and exchanged where the C exchanges the pointers \texttt{pu} and \texttt{pv} to them; an \texttt{Int} for the \texttt{long}s \texttt{du} and \texttt{dv}, which reach \texttt{-{}1}; a \texttt{Bool} for the \texttt{int} result.)
\item \csource{j-points.c}{623}{624}: excluded (memory): the initialization of the rationals of squarefree: the coefficients of the polynomials u and v, seven each, and q and t; every value that squarefree reads is set before (u up to the degree, v up to the degree less one, q and t by the operations)
\item \csource{j-points.c}{625}{629}: modeled by \lname{JPoints.Cli.squarefree}
  (deviation of \lname{JPoints.Cli.squarefree}: the globals \texttt{coeffs} and \texttt{degree} as arguments; arrays of seven \texttt{Rat}s for the arrays \texttt{u} and \texttt{v} of \texttt{mpq\_t}s, zero at the start as \texttt{mpq\_init} sets them (lines 623-624, memory), and exchanged where the C exchanges the pointers \texttt{pu} and \texttt{pv} to them; an \texttt{Int} for the \texttt{long}s \texttt{du} and \texttt{dv}, which reach \texttt{-{}1}; a \texttt{Bool} for the \texttt{int} result.)
\item \csource{j-points.c}{632}{632}: modeled by \lname{JPoints.Cli.squarefree}
  (deviation of \lname{JPoints.Cli.squarefree}: the globals \texttt{coeffs} and \texttt{degree} as arguments; arrays of seven \texttt{Rat}s for the arrays \texttt{u} and \texttt{v} of \texttt{mpq\_t}s, zero at the start as \texttt{mpq\_init} sets them (lines 623-624, memory), and exchanged where the C exchanges the pointers \texttt{pu} and \texttt{pv} to them; an \texttt{Int} for the \texttt{long}s \texttt{du} and \texttt{dv}, which reach \texttt{-{}1}; a \texttt{Bool} for the \texttt{int} result.)
\item \csource{j-points.c}{633}{640}: modeled by \lname{JPoints.Cli.remainderModulo}
  (deviation of \lname{JPoints.Cli.remainderModulo}: the coefficients that \texttt{pu} points to as an array, returned; \texttt{Rat} for the \texttt{mpq\_t}s; an \texttt{Int} for the \texttt{long}s \texttt{du} and \texttt{dv}.)
\item \csource{j-points.c}{641}{644}: modeled by \lname{JPoints.Cli.squarefree}
  (deviation of \lname{JPoints.Cli.squarefree}: the globals \texttt{coeffs} and \texttt{degree} as arguments; arrays of seven \texttt{Rat}s for the arrays \texttt{u} and \texttt{v} of \texttt{mpq\_t}s, zero at the start as \texttt{mpq\_init} sets them (lines 623-624, memory), and exchanged where the C exchanges the pointers \texttt{pu} and \texttt{pv} to them; an \texttt{Int} for the \texttt{long}s \texttt{du} and \texttt{dv}, which reach \texttt{-{}1}; a \texttt{Bool} for the \texttt{int} result.)
\item \csource{j-points.c}{645}{646}: excluded (memory): the release of the rationals of squarefree
\item \csource{j-points.c}{647}{648}: modeled by \lname{JPoints.Cli.squarefree}
  (deviation of \lname{JPoints.Cli.squarefree}: the globals \texttt{coeffs} and \texttt{degree} as arguments; arrays of seven \texttt{Rat}s for the arrays \texttt{u} and \texttt{v} of \texttt{mpq\_t}s, zero at the start as \texttt{mpq\_init} sets them (lines 623-624, memory), and exchanged where the C exchanges the pointers \texttt{pu} and \texttt{pv} to them; an \texttt{Int} for the \texttt{long}s \texttt{du} and \texttt{dv}, which reach \texttt{-{}1}; a \texttt{Bool} for the \texttt{int} result.)
\item \csource{j-points.c}{652}{653}: modeled by \lname{JPoints.Cli.unescape}
  (deviation of \lname{JPoints.Cli.unescape}: the string returned, for the copy that \texttt{unescape} allocates (lines 654-655, memory); the rest of the string from the pointer for the pointer \texttt{s}.)
\item \csource{j-points.c}{654}{655}: excluded (memory): the buffer of the unescaped format, as long as the argument, with the write position at its start, and error 7 when it cannot be had; the model builds the string from the empty string
\item \csource{j-points.c}{656}{663}: modeled by \lname{JPoints.Cli.unescape}
  (deviation of \lname{JPoints.Cli.unescape}: the string returned, for the copy that \texttt{unescape} allocates (lines 654-655, memory); the rest of the string from the pointer for the pointer \texttt{s}.)
\item \csource{j-points.c}{669}{681}: excluded (memory): init\_thread\_mpz, which initializes the multi-precision temporaries of the calling thread (lines 211-212, 216-217, 252) and which every further thread calls too (line 1373); the model has no temporaries
\item \csource{j-points.c}{683}{686}: modeled by \lname{JPoints.initMain}
  (deviation of \lname{JPoints.initMain}: the tables taken and returned, for the globals; a natural number for the index \texttt{n}. The initializations of the GMP integers (\texttt{j-{}points.c:687-{}693}) are memory management, not modeled.)
\item \csource{j-points.c}{687}{693}: excluded (memory): the initialization of the integers of the coefficients of f and of the Kummer surface (to 0, a value that is never read: read\_input sets every coefficient, kummer\_init every coefficient of the surface), and of the temporaries (init\_thread\_mpz)
\item \csource{j-points.c}{695}{702}: modeled by \lname{JPoints.initMain}
  (deviation of \lname{JPoints.initMain}: the tables taken and returned, for the globals; a natural number for the index \texttt{n}. The initializations of the GMP integers (\texttt{j-{}points.c:687-{}693}) are memory management, not modeled.)
\item \csource{j-points.c}{708}{716}: modeled by \lname{JPoints.invert}
  (deviation of \lname{JPoints.invert}: integers for the \texttt{long}s; \texttt{n \% 2 == 1} for \texttt{n \& 1} and \texttt{n >>> 1} for \texttt{n >>= 1}, which agree on the positive \texttt{n} of the loop; \texttt{Int.tmod} for C's \texttt{\%}.)
\item \csource{j-points.c}{718}{724}: modeled by \lname{JPoints.initInverses}
  (deviation of \lname{JPoints.initInverses}: the table taken and returned, for the global; natural numbers for the nonnegative \texttt{long}s \texttt{pn}, \texttt{p} and \texttt{a}.)
\item \csource{j-points.c}{726}{738}: modeled by \lname{JPoints.initSquares}
  (deviation of \lname{JPoints.initSquares}: the table taken and returned, for the global; natural numbers for the nonnegative \texttt{long}s \texttt{pn}, \texttt{p} and \texttt{a}, so that \texttt{a² \% p} is the remainder of natural numbers.)
\item \csource{j-points.c}{742}{746}: modeled by \lname{JPoints.compareEntries}
  (deviation of \lname{JPoints.compareEntries}: the records passed, for the pointers.)
\item \csource{j-points.c}{751}{753}: modeled by \lname{JPoints.rateOf}
\item \csource{j-points.c}{761}{761}: modeled by \lname{JPoints.modp}
  (deviation of \lname{JPoints.modp}: a function of \texttt{p} and \texttt{recip}, which the macro takes from the scope where it is used; integers for the \texttt{long}s \texttt{x} and \texttt{p}, with the conversion of \texttt{x} to an \texttt{unsigned long} modulo \texttt{2 \textasciicircum{} 64} as C converts it; the quotient is below \texttt{2 \textasciicircum{} 32}, so that the conversion back to a \texttt{long} keeps it, and \texttt{x -{} p q} fits a \texttt{long} for the operands of the calls, below \texttt{2 \textasciicircum{} 24}.)
\item \csource{j-points.c}{762}{777}: modeled by \lname{JPoints.pointOnJ}
  (deviation of \lname{JPoints.pointOnJ}: natural numbers for the index \texttt{pn} and the \texttt{long} \texttt{degree}, and the steps of the loop counted up, for the index \texttt{n} that counts down from \texttt{degree -{} 2} to \texttt{0}; the tables and the degree passed, for the globals that the C code reads; the entries of \texttt{squares} read at the natural number of the index, which the reductions make nonnegative; a Boolean for the \texttt{int}.)
\item \csource{j-points.c}{781}{795}: modeled by \lname{JPoints.initTable}
  (deviation of \lname{JPoints.initTable}: the tables and the degree taken and the tables of the curve returned, for the globals (the degree for \texttt{point\_on\_j}, which reads it); natural numbers for the indices \texttt{pn}, \texttt{b} and \texttt{c}, for the count, which the C code keeps nonnegative, and for the degree; \texttt{p \& 2} by the AND of natural numbers on the positive \texttt{p}; \texttt{Int.tdiv} and \texttt{Int.tmod} for C's \texttt{/} and \texttt{\%}.)
\item \csource{j-points.c}{804}{807}: modeled by \lname{JPoints.initFmodpsquare}
  (deviation of \lname{JPoints.initFmodpsquare}: the tables, the records, the plan and the globals of the condition at 2 taken and returned, for the globals, and the text returned, for the output on \texttt{stdout}; \texttt{num\_primes}, which \texttt{read\_input} leaves in \texttt{[1, 53]}, as a natural number; the messages as their texts, where the C code calls \texttt{message}, which prints them and flushes \texttt{stdout}.)
\item \csource{j-points.c}{808}{812}: modeled by \lname{JPoints.initFmodpsquareRates}
  (deviation of \lname{JPoints.initFmodpsquareRates}: the tables and the records taken and returned, for the globals; a natural number for the index \texttt{pn}; the local array \texttt{tabled}, which the C code leaves uninitialized, starts as \texttt{false}, and only the entries that the loop writes are read; the sort of glibc as \texttt{JPoints.qsort}, its merge sort, as glibc sorts when it can allocate its buffer.)
\item \csource{j-points.c}{813}{817}: modeled by \lname{JPoints.initFmodpsquarePrime}
  (deviation of \lname{JPoints.initFmodpsquarePrime}: the tables taken and the tables of the curve and the records returned, for the globals and the local arrays; natural numbers for the index \texttt{pn}, the \texttt{long} \texttt{degree}, the prime, its square, the stride and the counts, which the C code keeps nonnegative; Booleans for the \texttt{int} flags.)
\item \csource{j-points.c}{819}{820}: modeled by \lname{JPoints.initFmodpsquareCoeffs}
  (deviation of \lname{JPoints.initFmodpsquareCoeffs}: the table taken and returned, for the global; the remainder of \texttt{mpz\_fdiv\_r\_ui}, which rounds the quotient down, as \texttt{Int.emod} (\texttt{\%} of integers) for \texttt{p > 0}, and the GMP temporary \texttt{tmp}, into which the C code also writes it, left out (no reader of \texttt{tmp} reads it before writing it).)
\item \csource{j-points.c}{822}{828}: modeled by \lname{JPoints.initFmodpsquareValues}
  (deviation of \lname{JPoints.initFmodpsquareValues}: the tables taken and returned, for the globals; natural numbers for the index \texttt{pn}, the \texttt{long} \texttt{degree}, the prime and the count \texttt{np0}, and the entries of \texttt{squares} read at the natural number of the index, which the reductions make nonnegative; Booleans for the \texttt{unsigned char}s, each tested as the value of its assignment is in the C code.)
\item \csource{j-points.c}{829}{834}: modeled by \lname{JPoints.initFmodpsquareEval}
  (deviation of \lname{JPoints.initFmodpsquareEval}: natural numbers for the index \texttt{pn} and the \texttt{long} \texttt{degree}, and the steps of the loop counted up, for the index \texttt{n} that counts down from \texttt{degree -{} 1} to \texttt{0}; \texttt{Int.tmod} for C's \texttt{\%}.)
\item \csource{j-points.c}{835}{836}: modeled by \lname{JPoints.initFmodpsquareValues}
  (deviation of \lname{JPoints.initFmodpsquareValues}: the tables taken and returned, for the globals; natural numbers for the index \texttt{pn}, the \texttt{long} \texttt{degree}, the prime and the count \texttt{np0}, and the entries of \texttt{squares} read at the natural number of the index, which the reductions make nonnegative; Booleans for the \texttt{unsigned char}s, each tested as the value of its assignment is in the C code.)
\item \csource{j-points.c}{837}{842}: modeled by \lname{JPoints.initFmodpsquarePrime}
  (deviation of \lname{JPoints.initFmodpsquarePrime}: the tables taken and the tables of the curve and the records returned, for the globals and the local arrays; natural numbers for the index \texttt{pn}, the \texttt{long} \texttt{degree}, the prime, its square, the stride and the counts, which the C code keeps nonnegative; Booleans for the \texttt{int} flags.)
\item \csource{j-points.c}{843}{850}: modeled by \lname{JPoints.initFmodpsquareSample}
  (deviation of \lname{JPoints.initFmodpsquareSample}: the tables and the degree passed, for the globals that \texttt{point\_on\_j} reads; natural numbers for the index \texttt{pn}, the \texttt{long} \texttt{degree}, the prime, its square, the stride, the residues \texttt{i}, \texttt{b} and \texttt{c} and the counts, which the C code keeps nonnegative, so that \texttt{/} is the division of natural numbers; \texttt{Int.tdiv} and \texttt{Int.tmod} for C's \texttt{/} and \texttt{\%} in \texttt{b² / 4}; the loop \texttt{for(i = 0; i < pp; i += stride)} over \texttt{i = t stride} for \texttt{t < (pp + stride -{} 1) / stride = ⌈pp / stride⌉}, the values of \texttt{i} of the C code in their order, as \texttt{stride ≥ 1}.)
\item \csource{j-points.c}{851}{854}: modeled by \lname{JPoints.initFmodpsquarePrime}
  (deviation of \lname{JPoints.initFmodpsquarePrime}: the tables taken and the tables of the curve and the records returned, for the globals and the local arrays; natural numbers for the index \texttt{pn}, the \texttt{long} \texttt{degree}, the prime, its square, the stride and the counts, which the C code keeps nonnegative; Booleans for the \texttt{int} flags.)
\item \csource{j-points.c}{855}{855}: modeled by \lname{JPoints.initFmodpsquareRates}
  (deviation of \lname{JPoints.initFmodpsquareRates}: the tables and the records taken and returned, for the globals; a natural number for the index \texttt{pn}; the local array \texttt{tabled}, which the C code leaves uninitialized, starts as \texttt{false}, and only the entries that the loop writes are read; the sort of glibc as \texttt{JPoints.qsort}, its merge sort, as glibc sorts when it can allocate its buffer.)
\item \csource{j-points.c}{857}{857}: modeled by \lname{JPoints.initFmodpsquareRates}
  (deviation of \lname{JPoints.initFmodpsquareRates}: the tables and the records taken and returned, for the globals; a natural number for the index \texttt{pn}; the local array \texttt{tabled}, which the C code leaves uninitialized, starts as \texttt{false}, and only the entries that the loop writes are read; the sort of glibc as \texttt{JPoints.qsort}, its merge sort, as glibc sorts when it can allocate its buffer.)
\item \csource{j-points.c}{858}{858}: modeled by \lname{JPoints.initFmodpsquare}
  (deviation of \lname{JPoints.initFmodpsquare}: the tables, the records, the plan and the globals of the condition at 2 taken and returned, for the globals, and the text returned, for the output on \texttt{stdout}; \texttt{num\_primes}, which \texttt{read\_input} leaves in \texttt{[1, 53]}, as a natural number; the messages as their texts, where the C code calls \texttt{message}, which prints them and flushes \texttt{stdout}.)
\item \csource{j-points.c}{861}{861}: modeled by \lname{JPoints.initFmodpsquare}
  (deviation of \lname{JPoints.initFmodpsquare}: the tables, the records, the plan and the globals of the condition at 2 taken and returned, for the globals, and the text returned, for the output on \texttt{stdout}; \texttt{num\_primes}, which \texttt{read\_input} leaves in \texttt{[1, 53]}, as a natural number; the messages as their texts, where the C code calls \texttt{message}, which prints them and flushes \texttt{stdout}.)
\item \csource{j-points.c}{863}{864}: modeled by \lname{JPoints.initFmodpsquareTables}
  (deviation of \lname{JPoints.initFmodpsquareTables}: the tables of the curve and the flags taken and returned, for the globals and the local array; a natural number for the index \texttt{n}; the count that \texttt{init\_table} returns dropped, as the C code drops it.)
\item \csource{j-points.c}{865}{871}: modeled by \lname{JPoints.initFmodpsquare}
  (deviation of \lname{JPoints.initFmodpsquare}: the tables, the records, the plan and the globals of the condition at 2 taken and returned, for the globals, and the text returned, for the output on \texttt{stdout}; \texttt{num\_primes}, which \texttt{read\_input} leaves in \texttt{[1, 53]}, as a natural number; the messages as their texts, where the C code calls \texttt{message}, which prints them and flushes \texttt{stdout}.)
\item \csource{j-points.c}{876}{876}: modeled by \lname{JPoints.Cand}
  (deviation of \lname{JPoints.Cand}: a natural number for the \texttt{long} index; a \texttt{Float} for the double.)
\item \csource{j-points.c}{877}{880}: modeled by \lname{JPoints.compareCands}
  (deviation of \lname{JPoints.compareCands}: the records passed, for the pointers.)
\item \csource{j-points.c}{900}{900}: excluded (other): the forward declaration of choose\_primes\_mode; the model has no forward declarations
\item \csource{j-points.c}{907}{926}: modeled by \lname{JPoints.choosePrimes}
  (deviation of \lname{JPoints.choosePrimes}: the plan taken and returned, for the globals; the pins \texttt{pin1}, \texttt{pin2} and \texttt{pin3}, to which the C code sets \texttt{sieve\_primes1}, \texttt{sieve\_primes2} and \texttt{sieve\_primes3} back before each mode, read by \texttt{JPoints.choosePrimesMode} from \texttt{JPoints.Cli.Input}, which holds them; natural numbers for the \texttt{long} numbers and indices and the \texttt{int} modes, Booleans for the \texttt{int} cuts; \texttt{Float}s for the doubles \texttt{best} and \texttt{cost}; an array of \texttt{JPoints.numPrimes} entries for the local array \texttt{pn}, which starts as zeros (the C code leaves it uninitialized, and reads only the entries it writes).)
\item \csource{j-points.c}{928}{929}: modeled by \lname{JPoints.choosePrimesMode}
  (deviation of \lname{JPoints.choosePrimesMode}: the plan taken and returned, for the globals that the C code writes, its mode and cut unchanged; the pins \texttt{sieve\_primes1}, \texttt{sieve\_primes2} and \texttt{sieve\_primes3} read from \texttt{JPoints.Cli.Input}, which holds the values that \texttt{choose\_primes} sets them to before the call; a natural number for the \texttt{int} \texttt{mode} and a Boolean for the \texttt{int} flag \texttt{cut}.)
\item \csource{j-points.c}{930}{930}: modeled by \lname{JPoints.choosePrimesModeShape}
  (deviation of \lname{JPoints.choosePrimesModeShape}: the globals passed, for the globals that the C code reads, and the locals returned in a structure (\texttt{JPoints.ModeShape}); a natural number for the \texttt{int} \texttt{mode}, Booleans for the \texttt{int} flags \texttt{cut} and \texttt{tube}; \texttt{Float}s for the doubles, the integers converted to \texttt{Float} exactly, \texttt{H >>> 6} for \texttt{height>>LONG\_SHIFT} on the positive \texttt{H}; the test \texttt{mpz\_cmp\_si} of \texttt{f₅} with \texttt{1} as the equality of integers; \texttt{sqrt} and \texttt{floor} of \texttt{Float}.)
\item \csource{j-points.c}{931}{931}: modeled by \lname{JPoints.choosePrimesMode}, \lname{JPoints.choosePrimesModeList1}
  (deviation of \lname{JPoints.choosePrimesMode}: the plan taken and returned, for the globals that the C code writes, its mode and cut unchanged; the pins \texttt{sieve\_primes1}, \texttt{sieve\_primes2} and \texttt{sieve\_primes3} read from \texttt{JPoints.Cli.Input}, which holds the values that \texttt{choose\_primes} sets them to before the call; a natural number for the \texttt{int} \texttt{mode} and a Boolean for the \texttt{int} flag \texttt{cut}.)
  (deviation of \lname{JPoints.choosePrimesModeList1}: an array of the candidates for the array \texttt{list1}, whose size is their count \texttt{ne}; the sort of glibc as \texttt{JPoints.qsort}.)
\item \csource{j-points.c}{932}{932}: modeled by \lname{JPoints.choosePrimesMode}
  (deviation of \lname{JPoints.choosePrimesMode}: the plan taken and returned, for the globals that the C code writes, its mode and cut unchanged; the pins \texttt{sieve\_primes1}, \texttt{sieve\_primes2} and \texttt{sieve\_primes3} read from \texttt{JPoints.Cli.Input}, which holds the values that \texttt{choose\_primes} sets them to before the call; a natural number for the \texttt{int} \texttt{mode} and a Boolean for the \texttt{int} flag \texttt{cut}.)
\item \csource{j-points.c}{933}{933}: modeled by \lname{JPoints.ModeBest}
  (deviation of \lname{JPoints.ModeBest}: a structure for the locals; a \texttt{Float} for the double \texttt{best}; natural numbers for the \texttt{long} numbers and indices; arrays of \texttt{JPoints.numPrimes} entries for the local arrays.)
\item \csource{j-points.c}{934}{938}: modeled by \lname{JPoints.choosePrimesModeShape}
  (deviation of \lname{JPoints.choosePrimesModeShape}: the globals passed, for the globals that the C code reads, and the locals returned in a structure (\texttt{JPoints.ModeShape}); a natural number for the \texttt{int} \texttt{mode}, Booleans for the \texttt{int} flags \texttt{cut} and \texttt{tube}; \texttt{Float}s for the doubles, the integers converted to \texttt{Float} exactly, \texttt{H >>> 6} for \texttt{height>>LONG\_SHIFT} on the positive \texttt{H}; the test \texttt{mpz\_cmp\_si} of \texttt{f₅} with \texttt{1} as the equality of integers; \texttt{sqrt} and \texttt{floor} of \texttt{Float}.)
\item \csource{j-points.c}{939}{939}: modeled by \lname{JPoints.ModeBest}, \lname{JPoints.choosePrimesModeShape}
  (deviation of \lname{JPoints.ModeBest}: a structure for the locals; a \texttt{Float} for the double \texttt{best}; natural numbers for the \texttt{long} numbers and indices; arrays of \texttt{JPoints.numPrimes} entries for the local arrays.)
  (deviation of \lname{JPoints.choosePrimesModeShape}: the globals passed, for the globals that the C code reads, and the locals returned in a structure (\texttt{JPoints.ModeShape}); a natural number for the \texttt{int} \texttt{mode}, Booleans for the \texttt{int} flags \texttt{cut} and \texttt{tube}; \texttt{Float}s for the doubles, the integers converted to \texttt{Float} exactly, \texttt{H >>> 6} for \texttt{height>>LONG\_SHIFT} on the positive \texttt{H}; the test \texttt{mpz\_cmp\_si} of \texttt{f₅} with \texttt{1} as the equality of integers; \texttt{sqrt} and \texttt{floor} of \texttt{Float}.)
\item \csource{j-points.c}{940}{940}: modeled by \lname{JPoints.choosePrimesModeShape}
  (deviation of \lname{JPoints.choosePrimesModeShape}: the globals passed, for the globals that the C code reads, and the locals returned in a structure (\texttt{JPoints.ModeShape}); a natural number for the \texttt{int} \texttt{mode}, Booleans for the \texttt{int} flags \texttt{cut} and \texttt{tube}; \texttt{Float}s for the doubles, the integers converted to \texttt{Float} exactly, \texttt{H >>> 6} for \texttt{height>>LONG\_SHIFT} on the positive \texttt{H}; the test \texttt{mpz\_cmp\_si} of \texttt{f₅} with \texttt{1} as the equality of integers; \texttt{sqrt} and \texttt{floor} of \texttt{Float}.)
\item \csource{j-points.c}{941}{941}: modeled by \lname{JPoints.ModeBest}
  (deviation of \lname{JPoints.ModeBest}: a structure for the locals; a \texttt{Float} for the double \texttt{best}; natural numbers for the \texttt{long} numbers and indices; arrays of \texttt{JPoints.numPrimes} entries for the local arrays.)
\item \csource{j-points.c}{942}{942}: modeled by \lname{JPoints.choosePrimesModeShape}
  (deviation of \lname{JPoints.choosePrimesModeShape}: the globals passed, for the globals that the C code reads, and the locals returned in a structure (\texttt{JPoints.ModeShape}); a natural number for the \texttt{int} \texttt{mode}, Booleans for the \texttt{int} flags \texttt{cut} and \texttt{tube}; \texttt{Float}s for the doubles, the integers converted to \texttt{Float} exactly, \texttt{H >>> 6} for \texttt{height>>LONG\_SHIFT} on the positive \texttt{H}; the test \texttt{mpz\_cmp\_si} of \texttt{f₅} with \texttt{1} as the equality of integers; \texttt{sqrt} and \texttt{floor} of \texttt{Float}.)
\item \csource{j-points.c}{943}{943}: modeled by \lname{JPoints.choosePrimesModeTerms}
  (deviation of \lname{JPoints.choosePrimesModeTerms}: arrays of \texttt{JPoints.numPrimes} \texttt{Float}s for the local arrays, which start as zeros (the C code leaves them uninitialized, and reads only the entries it writes); \texttt{Float.log} for \texttt{log}, the \texttt{Float} of the bits of the rate.)
\item \csource{j-points.c}{944}{944}: modeled by \lname{JPoints.choosePrimesModeLoop1}
  (deviation of \lname{JPoints.choosePrimesModeLoop1}: an array of \texttt{JPoints.numPrimes} Booleans for the array of \texttt{int} flags \texttt{used}, which the C code leaves uninitialized (the first stage clears the entries it reads); \texttt{Float.pow} for \texttt{pow}; the pins \texttt{sieve\_primes2} and \texttt{sieve\_primes3} read from \texttt{JPoints.Cli.Input}.)
\item \csource{j-points.c}{951}{965}: modeled by \lname{JPoints.choosePrimesModeShape}
  (deviation of \lname{JPoints.choosePrimesModeShape}: the globals passed, for the globals that the C code reads, and the locals returned in a structure (\texttt{JPoints.ModeShape}); a natural number for the \texttt{int} \texttt{mode}, Booleans for the \texttt{int} flags \texttt{cut} and \texttt{tube}; \texttt{Float}s for the doubles, the integers converted to \texttt{Float} exactly, \texttt{H >>> 6} for \texttt{height>>LONG\_SHIFT} on the positive \texttt{H}; the test \texttt{mpz\_cmp\_si} of \texttt{f₅} with \texttt{1} as the equality of integers; \texttt{sqrt} and \texttt{floor} of \texttt{Float}.)
\item \csource{j-points.c}{968}{968}: modeled by \lname{JPoints.choosePrimesModeShape}
  (deviation of \lname{JPoints.choosePrimesModeShape}: the globals passed, for the globals that the C code reads, and the locals returned in a structure (\texttt{JPoints.ModeShape}); a natural number for the \texttt{int} \texttt{mode}, Booleans for the \texttt{int} flags \texttt{cut} and \texttt{tube}; \texttt{Float}s for the doubles, the integers converted to \texttt{Float} exactly, \texttt{H >>> 6} for \texttt{height>>LONG\_SHIFT} on the positive \texttt{H}; the test \texttt{mpz\_cmp\_si} of \texttt{f₅} with \texttt{1} as the equality of integers; \texttt{sqrt} and \texttt{floor} of \texttt{Float}.)
\item \csource{j-points.c}{969}{974}: modeled by \lname{JPoints.choosePrimesModeTerms}
  (deviation of \lname{JPoints.choosePrimesModeTerms}: arrays of \texttt{JPoints.numPrimes} \texttt{Float}s for the local arrays, which start as zeros (the C code leaves them uninitialized, and reads only the entries it writes); \texttt{Float.log} for \texttt{log}, the \texttt{Float} of the bits of the rate.)
\item \csource{j-points.c}{976}{984}: modeled by \lname{JPoints.choosePrimesModeList1}
  (deviation of \lname{JPoints.choosePrimesModeList1}: an array of the candidates for the array \texttt{list1}, whose size is their count \texttt{ne}; the sort of glibc as \texttt{JPoints.qsort}.)
\item \csource{j-points.c}{986}{990}: modeled by \lname{JPoints.choosePrimesModeRange1}
  (deviation of \lname{JPoints.choosePrimesModeRange1}: integers for the \texttt{long}s, the ends returned as natural numbers, which they are.)
\item \csource{j-points.c}{991}{991}: modeled by \lname{JPoints.choosePrimesModeLoop1}
  (deviation of \lname{JPoints.choosePrimesModeLoop1}: an array of \texttt{JPoints.numPrimes} Booleans for the array of \texttt{int} flags \texttt{used}, which the C code leaves uninitialized (the first stage clears the entries it reads); \texttt{Float.pow} for \texttt{pow}; the pins \texttt{sieve\_primes2} and \texttt{sieve\_primes3} read from \texttt{JPoints.Cli.Input}.)
\item \csource{j-points.c}{992}{992}: modeled by \lname{JPoints.choosePrimesModeFirst}
  (deviation of \lname{JPoints.choosePrimesModeFirst}: the flags taken and returned, for the local array that the loop over \texttt{n1} reuses; an array of Booleans for the array of \texttt{int} flags; the \texttt{Float} of the bits of the rate.)
\item \csource{j-points.c}{993}{993}: modeled by \lname{JPoints.choosePrimesModeList2}, \lname{JPoints.choosePrimesModeLoop1}
  (deviation of \lname{JPoints.choosePrimesModeList2}: an array of the candidates for the array \texttt{list2}, whose size is their count \texttt{ne2}; the sort of glibc as \texttt{JPoints.qsort}.)
  (deviation of \lname{JPoints.choosePrimesModeLoop1}: an array of \texttt{JPoints.numPrimes} Booleans for the array of \texttt{int} flags \texttt{used}, which the C code leaves uninitialized (the first stage clears the entries it reads); \texttt{Float.pow} for \texttt{pow}; the pins \texttt{sieve\_primes2} and \texttt{sieve\_primes3} read from \texttt{JPoints.Cli.Input}.)
\item \csource{j-points.c}{994}{1001}: modeled by \lname{JPoints.choosePrimesModeFirst}
  (deviation of \lname{JPoints.choosePrimesModeFirst}: the flags taken and returned, for the local array that the loop over \texttt{n1} reuses; an array of Booleans for the array of \texttt{int} flags; the \texttt{Float} of the bits of the rate.)
\item \csource{j-points.c}{1002}{1002}: modeled by \lname{JPoints.choosePrimesModeLoop1}
  (deviation of \lname{JPoints.choosePrimesModeLoop1}: an array of \texttt{JPoints.numPrimes} Booleans for the array of \texttt{int} flags \texttt{used}, which the C code leaves uninitialized (the first stage clears the entries it reads); \texttt{Float.pow} for \texttt{pow}; the pins \texttt{sieve\_primes2} and \texttt{sieve\_primes3} read from \texttt{JPoints.Cli.Input}.)
\item \csource{j-points.c}{1004}{1013}: modeled by \lname{JPoints.choosePrimesModeList2}
  (deviation of \lname{JPoints.choosePrimesModeList2}: an array of the candidates for the array \texttt{list2}, whose size is their count \texttt{ne2}; the sort of glibc as \texttt{JPoints.qsort}.)
\item \csource{j-points.c}{1014}{1018}: modeled by \lname{JPoints.choosePrimesModeRange2}
  (deviation of \lname{JPoints.choosePrimesModeRange2}: integers for the \texttt{long}s, the ends returned as natural numbers, which they are.)
\item \csource{j-points.c}{1019}{1019}: modeled by \lname{JPoints.choosePrimesModeLoop2}
  (deviation of \lname{JPoints.choosePrimesModeLoop2}: the pin \texttt{sieve\_primes3} read from \texttt{JPoints.Cli.Input}.)
\item \csource{j-points.c}{1020}{1020}: modeled by \lname{JPoints.choosePrimesModeSecond}
  (deviation of \lname{JPoints.choosePrimesModeSecond}: \texttt{Float.pow} for \texttt{pow}; the \texttt{Float} of the bits of the rate.)
\item \csource{j-points.c}{1021}{1021}: modeled by \lname{JPoints.choosePrimesModeList3}, \lname{JPoints.choosePrimesModeLoop2}
  (deviation of \lname{JPoints.choosePrimesModeList3}: an array of the candidates for the array \texttt{list3}, whose size is their count \texttt{ne3}; a Boolean for the \texttt{int} flag \texttt{in2}; the sort of glibc as \texttt{JPoints.qsort}.)
  (deviation of \lname{JPoints.choosePrimesModeLoop2}: the pin \texttt{sieve\_primes3} read from \texttt{JPoints.Cli.Input}.)
\item \csource{j-points.c}{1023}{1031}: modeled by \lname{JPoints.choosePrimesModeSecond}
  (deviation of \lname{JPoints.choosePrimesModeSecond}: \texttt{Float.pow} for \texttt{pow}; the \texttt{Float} of the bits of the rate.)
\item \csource{j-points.c}{1033}{1044}: modeled by \lname{JPoints.choosePrimesModeList3}
  (deviation of \lname{JPoints.choosePrimesModeList3}: an array of the candidates for the array \texttt{list3}, whose size is their count \texttt{ne3}; a Boolean for the \texttt{int} flag \texttt{in2}; the sort of glibc as \texttt{JPoints.qsort}.)
\item \csource{j-points.c}{1045}{1048}: modeled by \lname{JPoints.choosePrimesModeRange3}
  (deviation of \lname{JPoints.choosePrimesModeRange3}: integers for the \texttt{long}s, the ends returned as natural numbers, which they are.)
\item \csource{j-points.c}{1049}{1056}: modeled by \lname{JPoints.choosePrimesModeLoop3}
  (deviation of \lname{JPoints.choosePrimesModeLoop3}: the \texttt{Float} of the bits of the rate.)
\item \csource{j-points.c}{1057}{1060}: modeled by \lname{JPoints.choosePrimesModeKeep}
  (deviation of \lname{JPoints.choosePrimesModeKeep}: the cheapest choice taken and returned, for the locals of \texttt{choose\_primes\_mode} that the C code writes; natural numbers for the \texttt{long} numbers and indices; a \texttt{Float} for the double \texttt{total}.)
\item \csource{j-points.c}{1061}{1070}: modeled by \lname{JPoints.choosePrimesModeLoop3}
  (deviation of \lname{JPoints.choosePrimesModeLoop3}: the \texttt{Float} of the bits of the rate.)
\item \csource{j-points.c}{1071}{1071}: modeled by \lname{JPoints.choosePrimesModeLoop2}
  (deviation of \lname{JPoints.choosePrimesModeLoop2}: the pin \texttt{sieve\_primes3} read from \texttt{JPoints.Cli.Input}.)
\item \csource{j-points.c}{1072}{1072}: modeled by \lname{JPoints.choosePrimesModeLoop1}
  (deviation of \lname{JPoints.choosePrimesModeLoop1}: an array of \texttt{JPoints.numPrimes} Booleans for the array of \texttt{int} flags \texttt{used}, which the C code leaves uninitialized (the first stage clears the entries it reads); \texttt{Float.pow} for \texttt{pow}; the pins \texttt{sieve\_primes2} and \texttt{sieve\_primes3} read from \texttt{JPoints.Cli.Input}.)
\item \csource{j-points.c}{1073}{1078}: modeled by \lname{JPoints.choosePrimesMode}
  (deviation of \lname{JPoints.choosePrimesMode}: the plan taken and returned, for the globals that the C code writes, its mode and cut unchanged; the pins \texttt{sieve\_primes1}, \texttt{sieve\_primes2} and \texttt{sieve\_primes3} read from \texttt{JPoints.Cli.Input}, which holds the values that \texttt{choose\_primes} sets them to before the call; a natural number for the \texttt{int} \texttt{mode} and a Boolean for the \texttt{int} flag \texttt{cut}.)
\item \csource{j-points.c}{1085}{1085}: modeled by \lname{JPoints.SieveTables}, \lname{JPoints.SieveTables.initial}, \lname{JPoints.helpWidth}
  (deviation of \lname{JPoints.SieveTables}: an array of arrays for the array of pointers \texttt{sieve\_tab}, the empty array for a null pointer; a flat array for the array \texttt{help}, with the index arithmetic of the C code; a structure, for the globals.)
\item \csource{j-points.c}{1088}{1090}: modeled by \lname{JPoints.initSieve}
  (deviation of \lname{JPoints.initSieve}: the sieve tables passed and returned, for the globals \texttt{sieve\_tab} and \texttt{help}; the indices of the sieving primes in the table and their number in the first two stages passed, for the globals \texttt{pnn} and \texttt{sieve\_primes2}, and the tables passed, for the globals that the C code reads; each table stored once it is filled, for the pointer that the C code stores before it fills the table through it; natural numbers for the nonnegative \texttt{long}s \texttt{pn} and \texttt{n}.)
\item \csource{j-points.c}{1092}{1094}: excluded (instrumentation): the heading of the output of the sieve tables under DEBUG >= 2
\item \csource{j-points.c}{1095}{1097}: modeled by \lname{JPoints.initSieve}
  (deviation of \lname{JPoints.initSieve}: the sieve tables passed and returned, for the globals \texttt{sieve\_tab} and \texttt{help}; the indices of the sieving primes in the table and their number in the first two stages passed, for the globals \texttt{pnn} and \texttt{sieve\_primes2}, and the tables passed, for the globals that the C code reads; each table stored once it is filled, for the pointer that the C code stores before it fills the table through it; natural numbers for the nonnegative \texttt{long}s \texttt{pn} and \texttt{n}.)
\item \csource{j-points.c}{1098}{1100}: modeled by \lname{JPoints.initSievePrime}
  (deviation of \lname{JPoints.initSievePrime}: the table returned, for the table that the C code allocates and stores in \texttt{sieve\_tab} (\texttt{j-{}points.c:1101-{}1102}, memory management, not modeled), a fresh array of zeros, every word of which the rows write; \texttt{help} passed and returned, for the global; the tables passed, for the globals that the C code reads; natural numbers for the nonnegative \texttt{long}s, so that \texttt{/} is the division of natural numbers, as for the unsigned operand \texttt{LONG\_LENGTH + p -{} 1}.)
\item \csource{j-points.c}{1101}{1102}: excluded (memory): the allocation of the sieve table of the prime p on the heap, p p sieve\_rowlen(p) words, and error 7 when it fails; the model makes the table as a fresh array, every word of which the copy of the pattern writes (lines 1148-1169)
\item \csource{j-points.c}{1104}{1105}: modeled by \lname{JPoints.initSievePrime}
  (deviation of \lname{JPoints.initSievePrime}: the table returned, for the table that the C code allocates and stores in \texttt{sieve\_tab} (\texttt{j-{}points.c:1101-{}1102}, memory management, not modeled), a fresh array of zeros, every word of which the rows write; \texttt{help} passed and returned, for the global; the tables passed, for the globals that the C code reads; natural numbers for the nonnegative \texttt{long}s, so that \texttt{/} is the division of natural numbers, as for the unsigned operand \texttt{LONG\_LENGTH + p -{} 1}.)
\item \csource{j-points.c}{1107}{1109}: excluded (instrumentation): the prime p and aa, the index of the last word of its rows of help, printed under DEBUG >= 2
\item \csource{j-points.c}{1111}{1115}: modeled by \lname{JPoints.initSieveClear}
  (deviation of \lname{JPoints.initSieveClear}: the array passed and returned, for the global; natural numbers for the nonnegative \texttt{long}s; \texttt{c} for the index \texttt{(0 * maxTablePrime + 0) * helpWidth + c} of \texttt{help[0][0][c]}.)
\item \csource{j-points.c}{1117}{1122}: modeled by \lname{JPoints.initSieveInfinity}
  (deviation of \lname{JPoints.initSieveInfinity}: the array passed and returned, for the global; the tables of the curve and the table \texttt{bits} passed, for the globals that the C code reads; natural numbers for the nonnegative \texttt{long}s, so that \texttt{(b*c)\%p} is the remainder of natural numbers; the loop \texttt{for(i = (b*c)\%p; i < kp; i += p)} over \texttt{i = i₀ + t p}, \texttt{i₀ = (b*c)\%p}, for \texttt{t < (kp -{} i₀ + p -{} 1) / p = ⌈(kp -{} i₀) / p⌉}, the values of \texttt{i} of the C code in their order, as \texttt{p ≥ 3} and \texttt{i₀ < p ≤ kp}; \texttt{b * helpWidth + w} for the index \texttt{(0 * maxTablePrime + b) * helpWidth + w} of \texttt{help[0][b][w]}.)
\item \csource{j-points.c}{1124}{1131}: modeled by \lname{JPoints.initSieveOnJ}
  (deviation of \lname{JPoints.initSieveOnJ}: the array passed and returned, for the global; the tables of the curve and the table \texttt{bits} passed, for the globals that the C code reads; natural numbers for the nonnegative \texttt{long}s, so that \texttt{(a*b)\%p} and \texttt{(a*c)\%p} are remainders of natural numbers; the loop \texttt{for(i = (a*c)\%p; i < kp; i += p)} over \texttt{i = i₀ + t p}, \texttt{i₀ = (a*c)\%p}, for \texttt{t < (kp -{} i₀ + p -{} 1) / p = ⌈(kp -{} i₀) / p⌉}, the values of \texttt{i} of the C code in their order, as \texttt{p ≥ 3} and \texttt{i₀ < p ≤ kp}.)
\item \csource{j-points.c}{1133}{1143}: excluded (instrumentation): the rows of help of the prime printed word by word under DEBUG >= 3
\item \csource{j-points.c}{1148}{1149}: modeled by \lname{JPoints.initSieveRows}
  (deviation of \lname{JPoints.initSieveRows}: the table and \texttt{help} passed and returned, for the table that \texttt{sieve\_tab[pn]} points to and the global; natural numbers for the nonnegative \texttt{long}s \texttt{a} and \texttt{b}.)
\item \csource{j-points.c}{1150}{1159}: modeled by \lname{JPoints.initSieveRow}
  (deviation of \lname{JPoints.initSieveRow}: indices into the arrays for the pointers \texttt{si} and \texttt{he}, and the table and \texttt{help} passed and returned, for the table that \texttt{si} points into and the global; \texttt{sieve\_rowlen(p)} computed once, for its calls at \texttt{j-{}points.c:1150} and in the test of the last loop, which give the same value; natural numbers for the nonnegative \texttt{long}s, so that \texttt{LONG\_LENGTH/p} is the division of natural numbers, as for its unsigned operand, and the counter \texttt{c} of the last loop from \texttt{p}, where the C code goes on with \texttt{c} from the rotation, the same since \texttt{wp ≤ p}.)
\item \csource{j-points.c}{1161}{1167}: modeled by \lname{JPoints.initSieveRotate}
  (deviation of \lname{JPoints.initSieveRotate}: indices into the arrays for the pointers \texttt{si} and \texttt{he}, and the table and \texttt{help} passed and returned, for the table that \texttt{si} points into and the global; the values of the variables of \texttt{init\_sieve} that the loop reads passed; natural numbers for the nonnegative \texttt{long}s, and the counter \texttt{c} of the loop from \texttt{wp}, the value that the copy of the words \texttt{c < wp} of the row in \texttt{init\_sieve} leaves in it.)
\item \csource{j-points.c}{1169}{1170}: modeled by \lname{JPoints.initSieveRow}
  (deviation of \lname{JPoints.initSieveRow}: indices into the arrays for the pointers \texttt{si} and \texttt{he}, and the table and \texttt{help} passed and returned, for the table that \texttt{si} points into and the global; \texttt{sieve\_rowlen(p)} computed once, for its calls at \texttt{j-{}points.c:1150} and in the test of the last loop, which give the same value; natural numbers for the nonnegative \texttt{long}s, so that \texttt{LONG\_LENGTH/p} is the division of natural numbers, as for its unsigned operand, and the counter \texttt{c} of the last loop from \texttt{p}, where the C code goes on with \texttt{c} from the rotation, the same since \texttt{wp ≤ p}.)
\item \csource{j-points.c}{1172}{1182}: excluded (instrumentation): the sieve table of the prime printed word by word under DEBUG >= 3
\item \csource{j-points.c}{1184}{1186}: modeled by \lname{JPoints.initSieve}
  (deviation of \lname{JPoints.initSieve}: the sieve tables passed and returned, for the globals \texttt{sieve\_tab} and \texttt{help}; the indices of the sieving primes in the table and their number in the first two stages passed, for the globals \texttt{pnn} and \texttt{sieve\_primes2}, and the tables passed, for the globals that the C code reads; each table stored once it is filled, for the pointer that the C code stores before it fills the table through it; natural numbers for the nonnegative \texttt{long}s \texttt{pn} and \texttt{n}.)
\item \csource{j-points.c}{1193}{1198}: modeled by \lname{JPoints.findDoublePoints}
  (deviation of \lname{JPoints.findDoublePoints}: integers for the \texttt{long}s; \texttt{wheight} by the doubles of \texttt{JPoints.Model.DoubleArith}, the exact conversion of \texttt{height}, \texttt{sqrt} and \texttt{floor}, and the conversion back to a \texttt{long}, which truncates (\texttt{JPoints.Double.toInt64}); the local array \texttt{ssp0} of \texttt{NUM\_PRIMES} records that start as \texttt{JPoints.Stype}'s defaults, for the uninitialized array; the test \texttt{mpz\_cmp\_si} of \texttt{f₅} with \texttt{1} as the equality of integers; a Boolean for the \texttt{int} returned; the output taken and returned, for the thread-local counter \texttt{total} and \texttt{stdout}.)
\item \csource{j-points.c}{1199}{1199}: modeled by \lname{JPoints.Stype}, \lname{JPoints.findDoublePoints}
  (deviation of \lname{JPoints.Stype}: integers for the \texttt{long}s \texttt{p} and \texttt{ap}; a \texttt{JPoints.TablePtr} for the pointer \texttt{ptr}.)
  (deviation of \lname{JPoints.findDoublePoints}: integers for the \texttt{long}s; \texttt{wheight} by the doubles of \texttt{JPoints.Model.DoubleArith}, the exact conversion of \texttt{height}, \texttt{sqrt} and \texttt{floor}, and the conversion back to a \texttt{long}, which truncates (\texttt{JPoints.Double.toInt64}); the local array \texttt{ssp0} of \texttt{NUM\_PRIMES} records that start as \texttt{JPoints.Stype}'s defaults, for the uninitialized array; the test \texttt{mpz\_cmp\_si} of \texttt{f₅} with \texttt{1} as the equality of integers; a Boolean for the \texttt{int} returned; the output taken and returned, for the thread-local counter \texttt{total} and \texttt{stdout}.)
\item \csource{j-points.c}{1200}{1204}: modeled by \lname{JPoints.findDoublePoints}
  (deviation of \lname{JPoints.findDoublePoints}: integers for the \texttt{long}s; \texttt{wheight} by the doubles of \texttt{JPoints.Model.DoubleArith}, the exact conversion of \texttt{height}, \texttt{sqrt} and \texttt{floor}, and the conversion back to a \texttt{long}, which truncates (\texttt{JPoints.Double.toInt64}); the local array \texttt{ssp0} of \texttt{NUM\_PRIMES} records that start as \texttt{JPoints.Stype}'s defaults, for the uninitialized array; the test \texttt{mpz\_cmp\_si} of \texttt{f₅} with \texttt{1} as the equality of integers; a Boolean for the \texttt{int} returned; the output taken and returned, for the thread-local counter \texttt{total} and \texttt{stdout}.)
\item \csource{j-points.c}{1205}{1207}: modeled by \lname{JPoints.findDoublePointsSquares}
  (deviation of \lname{JPoints.findDoublePointsSquares}: integers for the \texttt{long}s; a \texttt{while} loop over \texttt{aa}, which ends on the data, for the \texttt{for} loop whose test sets \texttt{a}; the records taken, for the local array \texttt{ssp0}, and threaded through the loop; \texttt{true} for the \texttt{return(1)} of \texttt{find\_double\_points}; the output taken and returned, for the thread-local counter \texttt{total} and \texttt{stdout}.)
\item \csource{j-points.c}{1208}{1217}: modeled by \lname{JPoints.findDoublePointsSetup}
  (deviation of \lname{JPoints.findDoublePointsSetup}: one function for the loops of the two branches of \texttt{find\_double\_points}, which differ only in the test that skips \texttt{a} (line 1240), made under \texttt{skip}, and returning the records written so far with \texttt{false} for the \texttt{goto}; the records taken and returned, for the local array \texttt{ssp0}; the index \texttt{n} for the pointer \texttt{ssp}, which the C code moves along with \texttt{n}; integers for the \texttt{long}s, \texttt{Int.tmod} for C's \texttt{\%}.)
\item \csource{j-points.c}{1218}{1221}: modeled by \lname{JPoints.findDoublePointsRows}
  (deviation of \lname{JPoints.findDoublePointsRows}: one function for the loops of the two branches of \texttt{find\_double\_points}, which are the same code; integers for the \texttt{long}s, \texttt{/} for C's division of the positive \texttt{height} by \texttt{2 a}; \texttt{true} for the \texttt{return(1)} of \texttt{find\_double\_points}; the output taken and returned, for the thread-local counter \texttt{total} and \texttt{stdout}.)
\item \csource{j-points.c}{1222}{1226}: modeled by \lname{JPoints.findDoublePointsSieve}
  (deviation of \lname{JPoints.findDoublePointsSieve}: one function for the loops of the two branches of \texttt{find\_double\_points}, which are the same code; the index \texttt{k} of the record for the pointer \texttt{ssp}, which the C code moves on while it counts \texttt{n} down from \texttt{sieve\_primes3}; integers for the \texttt{long}s, \texttt{Int.tmod} for C's \texttt{\%}; \texttt{false} for the \texttt{goto} to the next \texttt{b} (\texttt{nextb1}, \texttt{nextb2}).)
\item \csource{j-points.c}{1227}{1231}: modeled by \lname{JPoints.findDoublePointsRows}
  (deviation of \lname{JPoints.findDoublePointsRows}: one function for the loops of the two branches of \texttt{find\_double\_points}, which are the same code; integers for the \texttt{long}s, \texttt{/} for C's division of the positive \texttt{height} by \texttt{2 a}; \texttt{true} for the \texttt{return(1)} of \texttt{find\_double\_points}; the output taken and returned, for the thread-local counter \texttt{total} and \texttt{stdout}.)
\item \csource{j-points.c}{1232}{1232}: modeled by \lname{JPoints.findDoublePointsSquares}
  (deviation of \lname{JPoints.findDoublePointsSquares}: integers for the \texttt{long}s; a \texttt{while} loop over \texttt{aa}, which ends on the data, for the \texttt{for} loop whose test sets \texttt{a}; the records taken, for the local array \texttt{ssp0}, and threaded through the loop; \texttt{true} for the \texttt{return(1)} of \texttt{find\_double\_points}; the output taken and returned, for the thread-local counter \texttt{total} and \texttt{stdout}.)
\item \csource{j-points.c}{1233}{1234}: modeled by \lname{JPoints.findDoublePoints}
  (deviation of \lname{JPoints.findDoublePoints}: integers for the \texttt{long}s; \texttt{wheight} by the doubles of \texttt{JPoints.Model.DoubleArith}, the exact conversion of \texttt{height}, \texttt{sqrt} and \texttt{floor}, and the conversion back to a \texttt{long}, which truncates (\texttt{JPoints.Double.toInt64}); the local array \texttt{ssp0} of \texttt{NUM\_PRIMES} records that start as \texttt{JPoints.Stype}'s defaults, for the uninitialized array; the test \texttt{mpz\_cmp\_si} of \texttt{f₅} with \texttt{1} as the equality of integers; a Boolean for the \texttt{int} returned; the output taken and returned, for the thread-local counter \texttt{total} and \texttt{stdout}.)
\item \csource{j-points.c}{1235}{1236}: modeled by \lname{JPoints.findDoublePointsAll}
  (deviation of \lname{JPoints.findDoublePointsAll}: integers for the \texttt{long}s; the records taken, for the local array \texttt{ssp0}, and threaded through the loop; the skip of \texttt{a} (\texttt{goto nexta2}) by the flag of \texttt{JPoints.findDoublePointsSetup}; \texttt{true} for the \texttt{return(1)} of \texttt{find\_double\_points}; the output taken and returned, for the thread-local counter \texttt{total} and \texttt{stdout}.)
\item \csource{j-points.c}{1237}{1247}: modeled by \lname{JPoints.findDoublePointsSetup}
  (deviation of \lname{JPoints.findDoublePointsSetup}: one function for the loops of the two branches of \texttt{find\_double\_points}, which differ only in the test that skips \texttt{a} (line 1240), made under \texttt{skip}, and returning the records written so far with \texttt{false} for the \texttt{goto}; the records taken and returned, for the local array \texttt{ssp0}; the index \texttt{n} for the pointer \texttt{ssp}, which the C code moves along with \texttt{n}; integers for the \texttt{long}s, \texttt{Int.tmod} for C's \texttt{\%}.)
\item \csource{j-points.c}{1248}{1251}: modeled by \lname{JPoints.findDoublePointsRows}
  (deviation of \lname{JPoints.findDoublePointsRows}: one function for the loops of the two branches of \texttt{find\_double\_points}, which are the same code; integers for the \texttt{long}s, \texttt{/} for C's division of the positive \texttt{height} by \texttt{2 a}; \texttt{true} for the \texttt{return(1)} of \texttt{find\_double\_points}; the output taken and returned, for the thread-local counter \texttt{total} and \texttt{stdout}.)
\item \csource{j-points.c}{1252}{1256}: modeled by \lname{JPoints.findDoublePointsSieve}
  (deviation of \lname{JPoints.findDoublePointsSieve}: one function for the loops of the two branches of \texttt{find\_double\_points}, which are the same code; the index \texttt{k} of the record for the pointer \texttt{ssp}, which the C code moves on while it counts \texttt{n} down from \texttt{sieve\_primes3}; integers for the \texttt{long}s, \texttt{Int.tmod} for C's \texttt{\%}; \texttt{false} for the \texttt{goto} to the next \texttt{b} (\texttt{nextb1}, \texttt{nextb2}).)
\item \csource{j-points.c}{1257}{1261}: modeled by \lname{JPoints.findDoublePointsRows}
  (deviation of \lname{JPoints.findDoublePointsRows}: one function for the loops of the two branches of \texttt{find\_double\_points}, which are the same code; integers for the \texttt{long}s, \texttt{/} for C's division of the positive \texttt{height} by \texttt{2 a}; \texttt{true} for the \texttt{return(1)} of \texttt{find\_double\_points}; the output taken and returned, for the thread-local counter \texttt{total} and \texttt{stdout}.)
\item \csource{j-points.c}{1262}{1263}: modeled by \lname{JPoints.findDoublePointsAll}
  (deviation of \lname{JPoints.findDoublePointsAll}: integers for the \texttt{long}s; the records taken, for the local array \texttt{ssp0}, and threaded through the loop; the skip of \texttt{a} (\texttt{goto nexta2}) by the flag of \texttt{JPoints.findDoublePointsSetup}; \texttt{true} for the \texttt{return(1)} of \texttt{find\_double\_points}; the output taken and returned, for the thread-local counter \texttt{total} and \texttt{stdout}.)
\item \csource{j-points.c}{1264}{1265}: modeled by \lname{JPoints.findDoublePoints}
  (deviation of \lname{JPoints.findDoublePoints}: integers for the \texttt{long}s; \texttt{wheight} by the doubles of \texttt{JPoints.Model.DoubleArith}, the exact conversion of \texttt{height}, \texttt{sqrt} and \texttt{floor}, and the conversion back to a \texttt{long}, which truncates (\texttt{JPoints.Double.toInt64}); the local array \texttt{ssp0} of \texttt{NUM\_PRIMES} records that start as \texttt{JPoints.Stype}'s defaults, for the uninitialized array; the test \texttt{mpz\_cmp\_si} of \texttt{f₅} with \texttt{1} as the equality of integers; a Boolean for the \texttt{int} returned; the output taken and returned, for the thread-local counter \texttt{total} and \texttt{stdout}.)
\item \csource{j-points.c}{1275}{1286}: modeled by \lname{JPoints.rowStep}
  (deviation of \lname{JPoints.rowStep}: integers for the \texttt{long}s, \texttt{Int.tmod} for C's \texttt{\%}; a natural number for the index \texttt{n}.)
\item \csource{j-points.c}{1288}{1290}: excluded (other): the forward declarations of unit\_a, run\_units and emit; the model has no forward declarations
\item \csource{j-points.c}{1292}{1293}: modeled by \lname{JPoints.findPoints}
  (deviation of \lname{JPoints.findPoints}: the totals returned, for the globals that \texttt{run\_units} sets; the globals, the static data and the state passed; the output taken and returned, for the thread-local counter \texttt{total} and \texttt{stdout}; the state of the thread from \texttt{JPoints.SiftState.initial} with the state of the analysis \texttt{t} and the output, for the thread-local globals, whose bit array and bands \texttt{init\_sift} allocates anew, and which nothing reads before it writes them.)
\item \csource{j-points.c}{1296}{1296}: modeled by \lname{JPoints.findPointsSetup}
  (deviation of \lname{JPoints.findPointsSetup}: the globals that it writes returned, the tables of the curve and the records of the primes from their initial values, which nothing writes before \texttt{init\_fmodpsquare}; the output taken and returned, for \texttt{stdout}.)
\item \csource{j-points.c}{1298}{1298}: modeled by \lname{JPoints.findPointsSetup}
  (deviation of \lname{JPoints.findPointsSetup}: the globals that it writes returned, the tables of the curve and the records of the primes from their initial values, which nothing writes before \texttt{init\_fmodpsquare}; the output taken and returned, for \texttt{stdout}.)
\item \csource{j-points.c}{1300}{1300}: modeled by \lname{JPoints.findPointsSetup}
  (deviation of \lname{JPoints.findPointsSetup}: the globals that it writes returned, the tables of the curve and the records of the primes from their initial values, which nothing writes before \texttt{init\_fmodpsquare}; the output taken and returned, for \texttt{stdout}.)
\item \csource{j-points.c}{1302}{1307}: modeled by \lname{JPoints.findPointsAtInfinity}
  (deviation of \lname{JPoints.findPointsAtInfinity}: integers for the GMP integers \texttt{fff}, \texttt{tmp} and \texttt{tmp2}, locals for the thread-local scratch integers; a Boolean for the test after which \texttt{find\_points} returns (\texttt{JPoints.findPoints} returns); the output taken and returned, for the thread-local counter \texttt{total} and \texttt{stdout}.)
\item \csource{j-points.c}{1308}{1308}: modeled by \lname{JPoints.findPoints}
  (deviation of \lname{JPoints.findPoints}: the totals returned, for the globals that \texttt{run\_units} sets; the globals, the static data and the state passed; the output taken and returned, for the thread-local counter \texttt{total} and \texttt{stdout}; the state of the thread from \texttt{JPoints.SiftState.initial} with the state of the analysis \texttt{t} and the output, for the thread-local globals, whose bit array and bands \texttt{init\_sift} allocates anew, and which nothing reads before it writes them.)
\item \csource{j-points.c}{1309}{1309}: modeled by \lname{JPoints.findPointsAtInfinity}
  (deviation of \lname{JPoints.findPointsAtInfinity}: integers for the GMP integers \texttt{fff}, \texttt{tmp} and \texttt{tmp2}, locals for the thread-local scratch integers; a Boolean for the test after which \texttt{find\_points} returns (\texttt{JPoints.findPoints} returns); the output taken and returned, for the thread-local counter \texttt{total} and \texttt{stdout}.)
\item \csource{j-points.c}{1311}{1311}: modeled by \lname{JPoints.findPoints}
  (deviation of \lname{JPoints.findPoints}: the totals returned, for the globals that \texttt{run\_units} sets; the globals, the static data and the state passed; the output taken and returned, for the thread-local counter \texttt{total} and \texttt{stdout}; the state of the thread from \texttt{JPoints.SiftState.initial} with the state of the analysis \texttt{t} and the output, for the thread-local globals, whose bit array and bands \texttt{init\_sift} allocates anew, and which nothing reads before it writes them.)
\item \csource{j-points.c}{1315}{1316}: modeled by \lname{JPoints.findPointsUnits}
  (deviation of \lname{JPoints.findPointsUnits}: an integer for the \texttt{long} \texttt{units}; \texttt{⌊√H⌋} by the doubles of \texttt{JPoints.Model.DoubleArith}, the exact conversion of \texttt{height}, \texttt{sqrt} and \texttt{floor}, and the conversion back to a \texttt{long}, which truncates (\texttt{JPoints.Double.toInt64}); the test \texttt{mpz\_cmp\_si} of \texttt{f₅} with \texttt{1} as the equality of integers.)
\item \csource{j-points.c}{1317}{1320}: modeled by \lname{JPoints.findPoints}
  (deviation of \lname{JPoints.findPoints}: the totals returned, for the globals that \texttt{run\_units} sets; the globals, the static data and the state passed; the output taken and returned, for the thread-local counter \texttt{total} and \texttt{stdout}; the state of the thread from \texttt{JPoints.SiftState.initial} with the state of the analysis \texttt{t} and the output, for the thread-local globals, whose bit array and bands \texttt{init\_sift} allocates anew, and which nothing reads before it writes them.)
\item \csource{j-points.c}{1324}{1325}: modeled by \lname{JPoints.unitA}
  (deviation of \lname{JPoints.unitA}: integers for the \texttt{long}s; the test \texttt{mpz\_cmp\_si} of \texttt{f₅} with \texttt{1} as the equality of integers.)
\item \csource{j-points.c}{1330}{1335}: modeled by \lname{JPoints.runA}
  (deviation of \lname{JPoints.runA}: the state taken and returned, for the thread-local globals; the globals and the static data that it reads passed; integers for the \texttt{long}s, with \texttt{Int.tdiv} for C's \texttt{/} on the positive \texttt{height} and \texttt{m}; the loop \texttt{for(b = b0; b <= height; b += m)} over \texttt{b = b₀ + k m} for \texttt{k < (H -{} b₀) / m + 1}, with Lean's \texttt{/}, which rounds down, none when \texttt{b₀ > H}; \texttt{a \% 64} and \texttt{b \% 64}, in \texttt{[0, 64)}, for \texttt{a \& 63} and \texttt{b \& 63}, the same in two's complement, and the entry \texttt{(a \& 63) * 64 + (b \& 63)} of the flat array \texttt{alive2}; a Boolean for the \texttt{int} returned.)
\item \csource{j-points.c}{1336}{1338}: excluded (instrumentation): the row (a, b) printed before it is sieved, under VERBOSE
\item \csource{j-points.c}{1339}{1343}: modeled by \lname{JPoints.runA}
  (deviation of \lname{JPoints.runA}: the state taken and returned, for the thread-local globals; the globals and the static data that it reads passed; integers for the \texttt{long}s, with \texttt{Int.tdiv} for C's \texttt{/} on the positive \texttt{height} and \texttt{m}; the loop \texttt{for(b = b0; b <= height; b += m)} over \texttt{b = b₀ + k m} for \texttt{k < (H -{} b₀) / m + 1}, with Lean's \texttt{/}, which rounds down, none when \texttt{b₀ > H}; \texttt{a \% 64} and \texttt{b \% 64}, in \texttt{[0, 64)}, for \texttt{a \& 63} and \texttt{b \& 63}, the same in two's complement, and the entry \texttt{(a \& 63) * 64 + (b \& 63)} of the flat array \texttt{alive2}; a Boolean for the \texttt{int} returned.)
\item \csource{j-points.c}{1345}{1383}: excluded (threads): finish\_unit and worker: the loop of a thread over the units of work it takes from the shared counter, the hand-over of the points of each unit to the printer, which prints the finished units in their order (with error 7 when the node of a unit cannot be had), and the sums of the counters of the threads; the model runs the units in turn on one thread (lines 1396-1403)
\item \csource{j-points.c}{1385}{1386}: modeled by \lname{JPoints.runUnits}
  (deviation of \lname{JPoints.runUnits}: the arm of one thread for every number of threads of \texttt{-{}t} (the threads are not modeled, and print the points in the order of the units, as this arm does, when the system lets them all start: when it refuses a thread or the memory for them, \texttt{j-{}points} exits with error 8 or 7 after part of the points, which depends on the limits of the machine; \texttt{j-{}points.c:1387-{}1395}, the bound of the number of threads by the number of units and the points collected for the threads, and \texttt{j-{}points.c:1404-{}1413} are excluded), so that the test \texttt{num\_threads == 1} is left out; the state taken and returned, for the thread-local globals, and the totals returned, for the globals; the globals and the static data that it reads passed; the loop over \texttt{k < units} as a loop over the natural numbers below \texttt{units.toNat}, \texttt{units} being nonnegative; natural numbers for the counters and the totals.)
\item \csource{j-points.c}{1387}{1395}: excluded (threads): the number of units for the workers, the bound of the number of threads by the number of units, and the printing of the points found before the units, which emit collects in the buffer of the calling thread when there are several threads; with one thread the bound leaves the number at 1 and the buffer is empty, since emit\_point and printf\_mpz print at once; with several it leaves more than one, since there are at least two units
\item \csource{j-points.c}{1396}{1403}: modeled by \lname{JPoints.runUnits}
  (deviation of \lname{JPoints.runUnits}: the arm of one thread for every number of threads of \texttt{-{}t} (the threads are not modeled, and print the points in the order of the units, as this arm does, when the system lets them all start: when it refuses a thread or the memory for them, \texttt{j-{}points} exits with error 8 or 7 after part of the points, which depends on the limits of the machine; \texttt{j-{}points.c:1387-{}1395}, the bound of the number of threads by the number of units and the points collected for the threads, and \texttt{j-{}points.c:1404-{}1413} are excluded), so that the test \texttt{num\_threads == 1} is left out; the state taken and returned, for the thread-local globals, and the totals returned, for the globals; the globals and the static data that it reads passed; the loop over \texttt{k < units} as a loop over the natural numbers below \texttt{units.toNat}, \texttt{units} being nonnegative; natural numbers for the counters and the totals.)
\item \csource{j-points.c}{1404}{1413}: excluded (threads): the run with several threads: the threads started (error 7 when their array cannot be had, error 8 when one cannot be started), the calling thread working as one of them (worker), the threads joined, and the return after them; the model runs the units in turn (lines 1396-1403), whose output the threads reproduce when they all start (errors 7 and 8 depend on the limits of the machine and are not modeled)
\item \csource{j-points.c}{1414}{1414}: modeled by \lname{JPoints.runUnits}
  (deviation of \lname{JPoints.runUnits}: the arm of one thread for every number of threads of \texttt{-{}t} (the threads are not modeled, and print the points in the order of the units, as this arm does, when the system lets them all start: when it refuses a thread or the memory for them, \texttt{j-{}points} exits with error 8 or 7 after part of the points, which depends on the limits of the machine; \texttt{j-{}points.c:1387-{}1395}, the bound of the number of threads by the number of units and the points collected for the threads, and \texttt{j-{}points.c:1404-{}1413} are excluded), so that the test \texttt{num\_threads == 1} is left out; the state taken and returned, for the thread-local globals, and the totals returned, for the globals; the globals and the static data that it reads passed; the loop over \texttt{k < units} as a loop over the natural numbers below \texttt{units.toNat}, \texttt{units} being nonnegative; natural numbers for the counters and the totals.)
\item \csource{j-points.c}{1421}{1463}: modeled by \lname{JPoints.kummerInit}
  (deviation of \lname{JPoints.kummerInit}: the coefficients returned in a structure, for the fifteen globals that the C code sets; integers for the GMP integers, without the scratch integer \texttt{tmp}.)
\item \csource{j-points.c}{1469}{1471}: modeled by \lname{JPoints.tubeInit}
  (deviation of \lname{JPoints.tubeInit}: the globals that it writes returned, the plan, the static data and the state of the analysis taken and returned; the globals that it reads passed.)
\item \csource{j-points.c}{1472}{1480}: modeled by \lname{JPoints.tubeInitCoeffs}
  (deviation of \lname{JPoints.tubeInitCoeffs}: the curve and the coefficients \texttt{k} passed, for the globals that it reads; the coefficients returned, for the globals \texttt{kd400, …, kd004} and \texttt{fd} that the C code writes; the array \texttt{fd} filled by the loop over \texttt{n ≤ 6} from seven zeros, for the global array, every entry of which the loop writes.)
\item \csource{j-points.c}{1481}{1485}: modeled by \lname{JPoints.tubeInit}
  (deviation of \lname{JPoints.tubeInit}: the globals that it writes returned, the plan, the static data and the state of the analysis taken and returned; the globals that it reads passed.)
\item \csource{j-points.c}{1497}{1498}: modeled by \lname{JPoints.SturmChain}
  (deviation of \lname{JPoints.SturmChain}: one flat array for the array \texttt{sturm[7][7]} of \texttt{mpq\_t}s, rationals, which start at \texttt{0} as \texttt{mpq\_init} sets them; natural numbers for the \texttt{long}s \texttt{sturm\_deg[i]} and \texttt{sturm\_len}, which the C code keeps nonnegative, starting at \texttt{0} as the static storage does.)
\item \csource{j-points.c}{1499}{1500}: modeled by \lname{JPoints.RootIntervals}
  (deviation of \lname{JPoints.RootIntervals}: arrays of six rationals for the arrays of \texttt{mpq\_t}s, which start at \texttt{0} as \texttt{mpq\_init} sets them; a natural number for the \texttt{long} \texttt{num\_roots}, which the C code keeps nonnegative.)
\item \csource{j-points.c}{1501}{1504}: modeled by \lname{JPoints.rootBits}
\item \csource{j-points.c}{1506}{1510}: modeled by \lname{JPoints.sturmChanges}
  (deviation of \lname{JPoints.sturmChanges}: the statics of the chain as an argument; rationals for the \texttt{mpq\_t}s; a natural number for the \texttt{long} count, an integer for the \texttt{int}s \texttt{s} and \texttt{last}; the steps of the loop of Horner's rule counted up, for the index \texttt{k} that counts down from \texttt{deg[i] -{} 1}.)
\item \csource{j-points.c}{1511}{1511}: excluded (memory): the initialization of the rational v of sturm\_changes, which is set before it is read
\item \csource{j-points.c}{1512}{1522}: modeled by \lname{JPoints.sturmChanges}
  (deviation of \lname{JPoints.sturmChanges}: the statics of the chain as an argument; rationals for the \texttt{mpq\_t}s; a natural number for the \texttt{long} count, an integer for the \texttt{int}s \texttt{s} and \texttt{last}; the steps of the loop of Horner's rule counted up, for the index \texttt{k} that counts down from \texttt{deg[i] -{} 1}.)
\item \csource{j-points.c}{1523}{1523}: excluded (memory): the release of v of sturm\_changes
\item \csource{j-points.c}{1524}{1525}: modeled by \lname{JPoints.sturmChanges}
  (deviation of \lname{JPoints.sturmChanges}: the statics of the chain as an argument; rationals for the \texttt{mpq\_t}s; a natural number for the \texttt{long} count, an integer for the \texttt{int}s \texttt{s} and \texttt{last}; the steps of the loop of Horner's rule counted up, for the index \texttt{k} that counts down from \texttt{deg[i] -{} 1}.)
\item \csource{j-points.c}{1528}{1532}: modeled by \lname{JPoints.sturmIsolate}
  (deviation of \lname{JPoints.sturmIsolate}: the fuel, the depth that the recursion may still descend, of which \texttt{JPoints.sturmIsolateFuel} suffices whenever \texttt{f} is squarefree, for a recursion that the C code ends by Sturm's theorem; none when it runs out, for which the caller records no interval; the statics as arguments and the isolating intervals returned; the ends by value, for the rationals \texttt{l2} and \texttt{h2} that the C code sets to the halves; natural numbers for the \texttt{long} counts, and \texttt{nlo -{} nhi == 1} as \texttt{nlo = nhi + 1}, which is the same for integers.)
\item \csource{j-points.c}{1533}{1533}: excluded (memory): the initialization of the rational mid of sturm\_isolate, which is set before it is read
\item \csource{j-points.c}{1534}{1535}: modeled by \lname{JPoints.sturmIsolate}
  (deviation of \lname{JPoints.sturmIsolate}: the fuel, the depth that the recursion may still descend, of which \texttt{JPoints.sturmIsolateFuel} suffices whenever \texttt{f} is squarefree, for a recursion that the C code ends by Sturm's theorem; none when it runs out, for which the caller records no interval; the statics as arguments and the isolating intervals returned; the ends by value, for the rationals \texttt{l2} and \texttt{h2} that the C code sets to the halves; natural numbers for the \texttt{long} counts, and \texttt{nlo -{} nhi == 1} as \texttt{nlo = nhi + 1}, which is the same for integers.)
\item \csource{j-points.c}{1536}{1543}: modeled by \lname{JPoints.sturmRefine}
  (deviation of \lname{JPoints.sturmRefine}: the ends by value and returned, for the rationals that the C code refines in place; the doubles by their bits (\texttt{JPoints.Model.DoubleArith}); natural numbers for the \texttt{long} counts, and \texttt{nlo -{} nmid == 1} as \texttt{nlo = nmid + 1}, which is the same for integers.)
\item \csource{j-points.c}{1544}{1545}: modeled by \lname{JPoints.sturmIsolate}
  (deviation of \lname{JPoints.sturmIsolate}: the fuel, the depth that the recursion may still descend, of which \texttt{JPoints.sturmIsolateFuel} suffices whenever \texttt{f} is squarefree, for a recursion that the C code ends by Sturm's theorem; none when it runs out, for which the caller records no interval; the statics as arguments and the isolating intervals returned; the ends by value, for the rationals \texttt{l2} and \texttt{h2} that the C code sets to the halves; natural numbers for the \texttt{long} counts, and \texttt{nlo -{} nhi == 1} as \texttt{nlo = nhi + 1}, which is the same for integers.)
\item \csource{j-points.c}{1546}{1546}: excluded (memory): the release of mid before the return after the refinement of one root
\item \csource{j-points.c}{1547}{1551}: modeled by \lname{JPoints.sturmIsolate}
  (deviation of \lname{JPoints.sturmIsolate}: the fuel, the depth that the recursion may still descend, of which \texttt{JPoints.sturmIsolateFuel} suffices whenever \texttt{f} is squarefree, for a recursion that the C code ends by Sturm's theorem; none when it runs out, for which the caller records no interval; the statics as arguments and the isolating intervals returned; the ends by value, for the rationals \texttt{l2} and \texttt{h2} that the C code sets to the halves; natural numbers for the \texttt{long} counts, and \texttt{nlo -{} nhi == 1} as \texttt{nlo = nhi + 1}, which is the same for integers.)
\item \csource{j-points.c}{1552}{1552}: excluded (memory): the initialization of l2 and h2, the copies of the ends of each half that the recursive calls refine in place
\item \csource{j-points.c}{1553}{1556}: modeled by \lname{JPoints.sturmIsolate}
  (deviation of \lname{JPoints.sturmIsolate}: the fuel, the depth that the recursion may still descend, of which \texttt{JPoints.sturmIsolateFuel} suffices whenever \texttt{f} is squarefree, for a recursion that the C code ends by Sturm's theorem; none when it runs out, for which the caller records no interval; the statics as arguments and the isolating intervals returned; the ends by value, for the rationals \texttt{l2} and \texttt{h2} that the C code sets to the halves; natural numbers for the \texttt{long} counts, and \texttt{nlo -{} nhi == 1} as \texttt{nlo = nhi + 1}, which is the same for integers.)
\item \csource{j-points.c}{1557}{1557}: excluded (memory): the release of l2 and h2
\item \csource{j-points.c}{1558}{1558}: modeled by \lname{JPoints.sturmIsolate}
  (deviation of \lname{JPoints.sturmIsolate}: the fuel, the depth that the recursion may still descend, of which \texttt{JPoints.sturmIsolateFuel} suffices whenever \texttt{f} is squarefree, for a recursion that the C code ends by Sturm's theorem; none when it runs out, for which the caller records no interval; the statics as arguments and the isolating intervals returned; the ends by value, for the rationals \texttt{l2} and \texttt{h2} that the C code sets to the halves; natural numbers for the \texttt{long} counts, and \texttt{nlo -{} nhi == 1} as \texttt{nlo = nhi + 1}, which is the same for integers.)
\item \csource{j-points.c}{1559}{1559}: excluded (memory): the release of mid
\item \csource{j-points.c}{1560}{1561}: modeled by \lname{JPoints.sturmIsolate}
  (deviation of \lname{JPoints.sturmIsolate}: the fuel, the depth that the recursion may still descend, of which \texttt{JPoints.sturmIsolateFuel} suffices whenever \texttt{f} is squarefree, for a recursion that the C code ends by Sturm's theorem; none when it runs out, for which the caller records no interval; the statics as arguments and the isolating intervals returned; the ends by value, for the rationals \texttt{l2} and \texttt{h2} that the C code sets to the halves; natural numbers for the \texttt{long} counts, and \texttt{nlo -{} nhi == 1} as \texttt{nlo = nhi + 1}, which is the same for integers.)
\item \csource{j-points.c}{1563}{1567}: modeled by \lname{JPoints.realRootsInit}
  (deviation of \lname{JPoints.realRootsInit}: the globals \texttt{coeffs} and \texttt{degree} as arguments, and the globals of the intervals returned, from their initial values, as \texttt{real\_roots\_init} runs once; a natural number for the \texttt{long} \texttt{degree}, which is \texttt{5} or \texttt{6}; \texttt{mpz\_sizeinbase(B, 2)} as the number of bits of \texttt{B}; the fuel \texttt{JPoints.sturmIsolateFuel} for the recursion of \texttt{sturm\_isolate}, and no interval when it runs out, which it does not for a squarefree \texttt{f}; the statics of the chain and of the isolating intervals passed on.)
\item \csource{j-points.c}{1568}{1570}: excluded (memory): the initialization of the rationals of the Sturm chain, of the isolating intervals and of q, lo and hi of real\_roots\_init (to 0, a value that is never read: the coefficients of the chain beyond the degree of a polynomial, which line 1583 copies along, are never read, and every other value is set before it is read)
\item \csource{j-points.c}{1572}{1579}: modeled by \lname{JPoints.realRootsStart}
  (deviation of \lname{JPoints.realRootsStart}: the globals \texttt{coeffs} and \texttt{degree} as arguments, and the chain returned, from the statics as they start (\texttt{0}); a natural number for the \texttt{long} \texttt{degree}; the product \texttt{k fₖ} by the rational \texttt{k}, as the C code makes it with the scratch \texttt{q}.)
\item \csource{j-points.c}{1581}{1584}: modeled by \lname{JPoints.realRootsChain}
  (deviation of \lname{JPoints.realRootsChain}: the chain by value and returned.)
\item \csource{j-points.c}{1585}{1588}: modeled by \lname{JPoints.realRootsReduce}
  (deviation of \lname{JPoints.realRootsReduce}: the chain by value and returned; an integer for the \texttt{long} \texttt{d}, which reaches \texttt{-{}1}; the scratch rational \texttt{t} of each step a local.)
\item \csource{j-points.c}{1589}{1589}: excluded (memory): the initialization of the scratch rational t of a step of the division in the Sturm chain, which line 1588 declares for each coefficient with the brace of the body of the loop (that line stays with the model); the model subtracts the product directly
\item \csource{j-points.c}{1590}{1591}: modeled by \lname{JPoints.realRootsReduce}
  (deviation of \lname{JPoints.realRootsReduce}: the chain by value and returned; an integer for the \texttt{long} \texttt{d}, which reaches \texttt{-{}1}; the scratch rational \texttt{t} of each step a local.)
\item \csource{j-points.c}{1592}{1592}: excluded (memory): the release of the scratch rational t
\item \csource{j-points.c}{1593}{1596}: modeled by \lname{JPoints.realRootsReduce}
  (deviation of \lname{JPoints.realRootsReduce}: the chain by value and returned; an integer for the \texttt{long} \texttt{d}, which reaches \texttt{-{}1}; the scratch rational \texttt{t} of each step a local.)
\item \csource{j-points.c}{1597}{1601}: modeled by \lname{JPoints.realRootsChain}
  (deviation of \lname{JPoints.realRootsChain}: the chain by value and returned.)
\item \csource{j-points.c}{1604}{1610}: modeled by \lname{JPoints.realRootsBound}
  (deviation of \lname{JPoints.realRootsBound}: the globals \texttt{coeffs} and \texttt{degree} as arguments, and the bound returned; integers for the GMP integers \texttt{B}, \texttt{t} and \texttt{fd}, \texttt{B} from \texttt{0} as \texttt{mpz\_init} sets it; the quotient rounded up, as \texttt{mpz\_cdiv\_q} does, as \texttt{-{}(-{}t / fd)} with Lean's \texttt{/}, which rounds down for the positive \texttt{fd}.)
\item \csource{j-points.c}{1611}{1618}: modeled by \lname{JPoints.realRootsInit}
  (deviation of \lname{JPoints.realRootsInit}: the globals \texttt{coeffs} and \texttt{degree} as arguments, and the globals of the intervals returned, from their initial values, as \texttt{real\_roots\_init} runs once; a natural number for the \texttt{long} \texttt{degree}, which is \texttt{5} or \texttt{6}; \texttt{mpz\_sizeinbase(B, 2)} as the number of bits of \texttt{B}; the fuel \texttt{JPoints.sturmIsolateFuel} for the recursion of \texttt{sturm\_isolate}, and no interval when it runs out, which it does not for a squarefree \texttt{f}; the statics of the chain and of the isolating intervals passed on.)
\item \csource{j-points.c}{1621}{1633}: modeled by \lname{JPoints.realRootsRecord}
  (deviation of \lname{JPoints.realRootsRecord}: the globals \texttt{coeffs} and \texttt{degree} and the statics of the isolating intervals as arguments, and the intervals returned; the steps of the loop counted up, for the index \texttt{i} that counts down from \texttt{num\_roots}; an integer for the \texttt{long}s \texttt{sgn\_inf} and \texttt{s}.)
\item \csource{j-points.c}{1634}{1634}: modeled by \lname{JPoints.realRootsInit}
  (deviation of \lname{JPoints.realRootsInit}: the globals \texttt{coeffs} and \texttt{degree} as arguments, and the globals of the intervals returned, from their initial values, as \texttt{real\_roots\_init} runs once; a natural number for the \texttt{long} \texttt{degree}, which is \texttt{5} or \texttt{6}; \texttt{mpz\_sizeinbase(B, 2)} as the number of bits of \texttt{B}; the fuel \texttt{JPoints.sturmIsolateFuel} for the recursion of \texttt{sturm\_isolate}, and no interval when it runs out, which it does not for a squarefree \texttt{f}; the statics of the chain and of the isolating intervals passed on.)
\item \csource{j-points.c}{1635}{1638}: excluded (memory): the release of the integers B, t and fd, of the rationals q, lo and hi, of the Sturm chain and of the isolating intervals at the end of real\_roots\_init (line 1604, whose mpz\_init gives B its start 0 for the maximum that follows, stays with the model)
\item \csource{j-points.c}{1639}{1640}: modeled by \lname{JPoints.realRootsInit}
  (deviation of \lname{JPoints.realRootsInit}: the globals \texttt{coeffs} and \texttt{degree} as arguments, and the globals of the intervals returned, from their initial values, as \texttt{real\_roots\_init} runs once; a natural number for the \texttt{long} \texttt{degree}, which is \texttt{5} or \texttt{6}; \texttt{mpz\_sizeinbase(B, 2)} as the number of bits of \texttt{B}; the fuel \texttt{JPoints.sturmIsolateFuel} for the recursion of \texttt{sturm\_isolate}, and no interval when it runs out, which it does not for a squarefree \texttt{f}; the statics of the chain and of the isolating intervals passed on.)
\item \csource{j-points.c}{1660}{1661}: modeled by \lname{JPoints.TwoadicTable}, \lname{JPoints.twoadicTable}
  (deviation of \lname{JPoints.TwoadicTable}: a structure, for the arrays and the integer that the arguments point to and for the result; natural numbers for the counts, which the C code keeps nonnegative.)
  (deviation of \lname{JPoints.twoadicTable}: the arrays taken and returned, for those that the arguments point to; the scratch array \texttt{sol} a value, its allocation and release (\texttt{j-{}points.c:1663}, \texttt{j-{}points.c:1668}, \texttt{j-{}points.c:1737}, with the exit of a failed allocation) memory management, not modeled.)
\item \csource{j-points.c}{1662}{1662}: modeled by \lname{JPoints.twoadicTable}
  (deviation of \lname{JPoints.twoadicTable}: the arrays taken and returned, for those that the arguments point to; the scratch array \texttt{sol} a value, its allocation and release (\texttt{j-{}points.c:1663}, \texttt{j-{}points.c:1668}, \texttt{j-{}points.c:1737}, with the exit of a failed allocation) memory management, not modeled.)
\item \csource{j-points.c}{1663}{1663}: excluded (memory): the allocation of the table sol of 64\textasciicircum{}3 words of twoadic\_table on the heap; the model makes the table as a fresh array, which line 1672 clears
\item \csource{j-points.c}{1664}{1664}: modeled by \lname{JPoints.twoadicTable}
  (deviation of \lname{JPoints.twoadicTable}: the arrays taken and returned, for those that the arguments point to; the scratch array \texttt{sol} a value, its allocation and release (\texttt{j-{}points.c:1663}, \texttt{j-{}points.c:1668}, \texttt{j-{}points.c:1737}, with the exit of a failed allocation) memory management, not modeled.)
\item \csource{j-points.c}{1665}{1665}: modeled by \lname{JPoints.twoadicClasses}
  (deviation of \lname{JPoints.twoadicClasses}: the arrays taken and returned with the counts, for the arrays and the integer that the arguments of the C code point to and for its result; natural numbers for the residues and the counts; the count of each row added after its loop.)
\item \csource{j-points.c}{1666}{1667}: modeled by \lname{JPoints.twoadicTable}
  (deviation of \lname{JPoints.twoadicTable}: the arrays taken and returned, for those that the arguments point to; the scratch array \texttt{sol} a value, its allocation and release (\texttt{j-{}points.c:1663}, \texttt{j-{}points.c:1668}, \texttt{j-{}points.c:1737}, with the exit of a failed allocation) memory management, not modeled.)
\item \csource{j-points.c}{1668}{1668}: excluded (memory): error 7 when the table sol cannot be had
\item \csource{j-points.c}{1669}{1669}: modeled by \lname{JPoints.twoadicTable}
  (deviation of \lname{JPoints.twoadicTable}: the arrays taken and returned, for those that the arguments point to; the scratch array \texttt{sol} a value, its allocation and release (\texttt{j-{}points.c:1663}, \texttt{j-{}points.c:1668}, \texttt{j-{}points.c:1737}, with the exit of a failed allocation) memory management, not modeled.)
\item \csource{j-points.c}{1670}{1671}: modeled by \lname{JPoints.twoadicSquares}
  (deviation of \lname{JPoints.twoadicSquares}: Booleans for the \texttt{int}s \texttt{0} and \texttt{1}; the local array of \texttt{64} entries, which the C code leaves uninitialized and the first loop clears, starting as \texttt{false}; a natural number for \texttt{i}, and \texttt{i * i \% 64} for \texttt{(i*i) \& 63}.)
\item \csource{j-points.c}{1672}{1678}: modeled by \lname{JPoints.twoadicSol}
  (deviation of \lname{JPoints.twoadicSol}: a fresh array of zeros for the array that the C code allocates and the loop clears (\texttt{malloc}, memory management, not modeled); natural numbers for the nonnegative \texttt{long}s; \texttt{x \% 64} for \texttt{x \& 63}, equal for the nonnegative \texttt{x}.)
\item \csource{j-points.c}{1679}{1688}: modeled by \lname{JPoints.twoadicCoeffs}
  (deviation of \lname{JPoints.twoadicCoeffs}: integers for the \texttt{long}s, \texttt{x \% 64}, the Euclidean remainder, for the remainder of the floor division of \texttt{mpz\_fdiv\_ui}; the arrays returned, for the local arrays; a fresh array of zeros for the array \texttt{km}, which the loop at \texttt{j-{}points.c:1680} clears.)
\item \csource{j-points.c}{1689}{1691}: modeled by \lname{JPoints.twoadicClasses}
  (deviation of \lname{JPoints.twoadicClasses}: the arrays taken and returned with the counts, for the arrays and the integer that the arguments of the C code point to and for its result; natural numbers for the residues and the counts; the count of each row added after its loop.)
\item \csource{j-points.c}{1692}{1692}: modeled by \lname{JPoints.twoadicPowers}
  (deviation of \lname{JPoints.twoadicPowers}: one function for the three loops of the C code, which differ only in the array and the residue; the array returned, for the local array of five \texttt{long}s, which starts as zeros; a natural number for the residue \texttt{x}; \texttt{y \% 64} for \texttt{y \& 63}, equal for the nonnegative \texttt{y}.)
\item \csource{j-points.c}{1693}{1693}: modeled by \lname{JPoints.twoadicClasses}
  (deviation of \lname{JPoints.twoadicClasses}: the arrays taken and returned with the counts, for the arrays and the integer that the arguments of the C code point to and for its result; natural numbers for the residues and the counts; the count of each row added after its loop.)
\item \csource{j-points.c}{1694}{1694}: modeled by \lname{JPoints.twoadicRow}
  (deviation of \lname{JPoints.twoadicRow}: the number of the admitted \texttt{c2} returned, which the caller adds to the count that the C code increments for each; natural numbers for the residues and the count.)
\item \csource{j-points.c}{1695}{1695}: modeled by \lname{JPoints.twoadicPowers}
  (deviation of \lname{JPoints.twoadicPowers}: one function for the three loops of the C code, which differ only in the array and the residue; the array returned, for the local array of five \texttt{long}s, which starts as zeros; a natural number for the residue \texttt{x}; \texttt{y \% 64} for \texttt{y \& 63}, equal for the nonnegative \texttt{y}.)
\item \csource{j-points.c}{1696}{1696}: modeled by \lname{JPoints.twoadicRow}
  (deviation of \lname{JPoints.twoadicRow}: the number of the admitted \texttt{c2} returned, which the caller adds to the count that the C code increments for each; natural numbers for the residues and the count.)
\item \csource{j-points.c}{1697}{1697}: modeled by \lname{JPoints.twoadicCell}, \lname{JPoints.twoadicK0}
  (deviation of \lname{JPoints.twoadicCell}: natural numbers for the residues \texttt{a2}, \texttt{b2} and \texttt{c2}; integers for the other \texttt{long}s; \texttt{x \% 64} for \texttt{x \& 63}, equal also for the negative \texttt{x} of \texttt{k₂} and \texttt{kk} (two's complement).)
  (deviation of \lname{JPoints.twoadicK0}: natural numbers for the indices \texttt{i}, \texttt{j} and \texttt{k}, the bound \texttt{i + j < 5} written as \texttt{j < 5 -{} i}; \texttt{x \% 64} for \texttt{x \& 63}.)
\item \csource{j-points.c}{1698}{1699}: modeled by \lname{JPoints.twoadicAdmits}
  (deviation of \lname{JPoints.twoadicAdmits}: a Boolean for the \texttt{int} \texttt{ok}; natural numbers for the indices \texttt{d} and \texttt{e}; the condition \texttt{!ok} of the loops a \texttt{break} at the start of each step; the entries of \texttt{sol} and \texttt{sq} read at the natural number of the index, which the reductions make nonnegative; \texttt{x \% 64} for \texttt{x \& 63}.)
\item \csource{j-points.c}{1700}{1700}: modeled by \lname{JPoints.twoadicRow}
  (deviation of \lname{JPoints.twoadicRow}: the number of the admitted \texttt{c2} returned, which the caller adds to the count that the C code increments for each; natural numbers for the residues and the count.)
\item \csource{j-points.c}{1701}{1701}: modeled by \lname{JPoints.twoadicPowers}
  (deviation of \lname{JPoints.twoadicPowers}: one function for the three loops of the C code, which differ only in the array and the residue; the array returned, for the local array of five \texttt{long}s, which starts as zeros; a natural number for the residue \texttt{x}; \texttt{y \% 64} for \texttt{y \& 63}, equal for the nonnegative \texttt{y}.)
\item \csource{j-points.c}{1702}{1706}: modeled by \lname{JPoints.twoadicK0}
  (deviation of \lname{JPoints.twoadicK0}: natural numbers for the indices \texttt{i}, \texttt{j} and \texttt{k}, the bound \texttt{i + j < 5} written as \texttt{j < 5 -{} i}; \texttt{x \% 64} for \texttt{x \& 63}.)
\item \csource{j-points.c}{1707}{1714}: modeled by \lname{JPoints.twoadicCell}
  (deviation of \lname{JPoints.twoadicCell}: natural numbers for the residues \texttt{a2}, \texttt{b2} and \texttt{c2}; integers for the other \texttt{long}s; \texttt{x \% 64} for \texttt{x \& 63}, equal also for the negative \texttt{x} of \texttt{k₂} and \texttt{kk} (two's complement).)
\item \csource{j-points.c}{1716}{1729}: modeled by \lname{JPoints.twoadicAdmits}
  (deviation of \lname{JPoints.twoadicAdmits}: a Boolean for the \texttt{int} \texttt{ok}; natural numbers for the indices \texttt{d} and \texttt{e}; the condition \texttt{!ok} of the loops a \texttt{break} at the start of each step; the entries of \texttt{sol} and \texttt{sq} read at the natural number of the index, which the reductions make nonnegative; \texttt{x \% 64} for \texttt{x \& 63}.)
\item \csource{j-points.c}{1730}{1731}: modeled by \lname{JPoints.twoadicRow}
  (deviation of \lname{JPoints.twoadicRow}: the number of the admitted \texttt{c2} returned, which the caller adds to the count that the C code increments for each; natural numbers for the residues and the count.)
\item \csource{j-points.c}{1732}{1736}: modeled by \lname{JPoints.twoadicClasses}
  (deviation of \lname{JPoints.twoadicClasses}: the arrays taken and returned with the counts, for the arrays and the integer that the arguments of the C code point to and for its result; natural numbers for the residues and the counts; the count of each row added after its loop.)
\item \csource{j-points.c}{1737}{1737}: excluded (memory): the release of the table sol
\item \csource{j-points.c}{1738}{1739}: modeled by \lname{JPoints.twoadicTable}
  (deviation of \lname{JPoints.twoadicTable}: the arrays taken and returned, for those that the arguments point to; the scratch array \texttt{sol} a value, its allocation and release (\texttt{j-{}points.c:1663}, \texttt{j-{}points.c:1668}, \texttt{j-{}points.c:1737}, with the exit of a failed allocation) memory management, not modeled.)
\item \csource{j-points.c}{1741}{1753}: modeled by \lname{JPoints.twoadicInit}
  (deviation of \lname{JPoints.twoadicInit}: the globals taken and returned; \texttt{Float}s for the doubles, the integers converted exactly.)
\item \csource{j-points.c}{1759}{1815}: modeled by \lname{JPoints.kummerEqn}
  (deviation of \lname{JPoints.kummerEqn}: integers for the \texttt{long}s: the products \texttt{aa, …, bc} and the factors \texttt{-{}4a}, \texttt{-{}2b}, \texttt{-{}4c} are \texttt{long}s in the C code, which hold them for \texttt{|a|, |b|, |c| ≤ JPoints.maxHeight}, as the comment of the function requires and its callers ensure; the equation returned in a structure, for the thread-local array \texttt{kummer} that the C code fills; no scratch integer \texttt{tmp}.)
\item \csource{j-points.c}{1821}{1834}: modeled by \lname{JPoints.relprime}
  (deviation of \lname{JPoints.relprime}: integers for the \texttt{long}s, so that the negation of \texttt{m = -{}2 \textasciicircum{} 63}, which overflows in the C code, is exact (the callers pass \texttt{|m| ≤ ⌊√H⌋}); \texttt{x \% 2 == 0} for \texttt{!(x \& 1)}, which agree on every \texttt{long} in two's complement; \texttt{x >>> 1}, the floor of \texttt{x / 2}, for \texttt{x >>= 1}, which gcc computes so also for a negative \texttt{long}; a Boolean for the \texttt{int}.)
\item \csource{j-points.c}{1835}{1840}: modeled by \lname{JPoints.relprimeLoop}
  (deviation of \lname{JPoints.relprimeLoop}: integers for the \texttt{long}s; \texttt{x \% 2 == 0} for \texttt{!(x \& 1)} and \texttt{x >>> 1}, the floor of \texttt{x / 2}, for \texttt{x >>= 1}, which gcc computes so also for a negative \texttt{long}; the variables \texttt{n} and \texttt{m} passed in and returned.)
\item \csource{j-points.c}{1841}{1842}: modeled by \lname{JPoints.relprime}
  (deviation of \lname{JPoints.relprime}: integers for the \texttt{long}s, so that the negation of \texttt{m = -{}2 \textasciicircum{} 63}, which overflows in the C code, is exact (the callers pass \texttt{|m| ≤ ⌊√H⌋}); \texttt{x \% 2 == 0} for \texttt{!(x \& 1)}, which agree on every \texttt{long} in two's complement; \texttt{x >>> 1}, the floor of \texttt{x / 2}, for \texttt{x >>= 1}, which gcc computes so also for a negative \texttt{long}; a Boolean for the \texttt{int}.)
\item \csource{j-points.c}{1849}{1865}: modeled by \lname{JPoints.checkLifts}
  (deviation of \lname{JPoints.checkLifts}: integers for the \texttt{long}s: the squares \texttt{a²}, \texttt{b²} and \texttt{c²} are \texttt{long}s in the C code, which hold them for \texttt{|a|, |b|, |c| ≤ JPoints.maxHeight}, as \texttt{check\_one\_point\_final} ensures; \texttt{x12}, \texttt{x22} and \texttt{x32} locals, for the thread-local scratch integers; the computations of \texttt{b⁶ f(c / b)}, \texttt{A²} and \texttt{B²} functions of their own; a Boolean for the \texttt{int}.)
\item \csource{j-points.c}{1867}{1878}: modeled by \lname{JPoints.checkLiftsForm}
  (deviation of \lname{JPoints.checkLiftsForm}: integers for the \texttt{long}s \texttt{b} and \texttt{c}; a fresh array for the thread-local array \texttt{bc}, whose entries \texttt{0, …, 5} the first loop writes before the second reads them; the steps of the loops counted up, for the index \texttt{k} that counts down from \texttt{5}; the value returned, for the scratch integer \texttt{fff}.)
\item \csource{j-points.c}{1879}{1889}: modeled by \lname{JPoints.checkLifts}
  (deviation of \lname{JPoints.checkLifts}: integers for the \texttt{long}s: the squares \texttt{a²}, \texttt{b²} and \texttt{c²} are \texttt{long}s in the C code, which hold them for \texttt{|a|, |b|, |c| ≤ JPoints.maxHeight}, as \texttt{check\_one\_point\_final} ensures; \texttt{x12}, \texttt{x22} and \texttt{x32} locals, for the thread-local scratch integers; the computations of \texttt{b⁶ f(c / b)}, \texttt{A²} and \texttt{B²} functions of their own; a Boolean for the \texttt{int}.)
\item \csource{j-points.c}{1891}{1910}: modeled by \lname{JPoints.checkLiftsA2}
  (deviation of \lname{JPoints.checkLiftsA2}: integers for the \texttt{long}s: the products \texttt{a d} and \texttt{a b} are \texttt{long}s in the C code, which hold them for \texttt{|a|, |b|, |d| ≤ JPoints.maxHeight}, as the callers ensure (\texttt{b² -{} ac}, which can exceed a \texttt{long}, is a GMP integer there); the value returned, for the scratch integer \texttt{tmp}, without the scratch integers \texttt{tmp2} and \texttt{tmp3}.)
\item \csource{j-points.c}{1911}{1913}: excluded (instrumentation): the value of A\textasciicircum{}2 in check\_lifts printed under VERBOSE
\item \csource{j-points.c}{1915}{1916}: modeled by \lname{JPoints.checkLifts}
  (deviation of \lname{JPoints.checkLifts}: integers for the \texttt{long}s: the squares \texttt{a²}, \texttt{b²} and \texttt{c²} are \texttt{long}s in the C code, which hold them for \texttt{|a|, |b|, |c| ≤ JPoints.maxHeight}, as \texttt{check\_one\_point\_final} ensures; \texttt{x12}, \texttt{x22} and \texttt{x32} locals, for the thread-local scratch integers; the computations of \texttt{b⁶ f(c / b)}, \texttt{A²} and \texttt{B²} functions of their own; a Boolean for the \texttt{int}.)
\item \csource{j-points.c}{1918}{1930}: modeled by \lname{JPoints.checkLiftsB2}
  (deviation of \lname{JPoints.checkLiftsB2}: integers for the \texttt{long}s: the products \texttt{c d} and \texttt{a b} are \texttt{long}s in the C code, which hold them for \texttt{|a|, |b|, |c|, |d| ≤ JPoints.maxHeight}, as the callers ensure; the value returned, for the scratch integer \texttt{tmp}, without the scratch integer \texttt{tmp2}.)
\item \csource{j-points.c}{1931}{1933}: excluded (instrumentation): the value of B\textasciicircum{}2 in check\_lifts printed under VERBOSE
\item \csource{j-points.c}{1934}{1936}: modeled by \lname{JPoints.checkLifts}
  (deviation of \lname{JPoints.checkLifts}: integers for the \texttt{long}s: the squares \texttt{a²}, \texttt{b²} and \texttt{c²} are \texttt{long}s in the C code, which hold them for \texttt{|a|, |b|, |c| ≤ JPoints.maxHeight}, as \texttt{check\_one\_point\_final} ensures; \texttt{x12}, \texttt{x22} and \texttt{x32} locals, for the thread-local scratch integers; the computations of \texttt{b⁶ f(c / b)}, \texttt{A²} and \texttt{B²} functions of their own; a Boolean for the \texttt{int}.)
\item \csource{j-points.c}{1938}{1954}: modeled by \lname{JPoints.checkLiftsMpz}
  (deviation of \lname{JPoints.checkLiftsMpz}: integers for the GMP integers, the arguments passed by value, for the pointers; \texttt{x12}, \texttt{x22} and \texttt{x32} locals, for the thread-local scratch integers; the computations of \texttt{b⁶ f(c / b)}, \texttt{A²} and \texttt{B²} functions of their own; a Boolean for the \texttt{int}.)
\item \csource{j-points.c}{1956}{1967}: modeled by \lname{JPoints.checkLiftsMpzForm}
  (deviation of \lname{JPoints.checkLiftsMpzForm}: a fresh array for the thread-local array \texttt{bc}, whose entries \texttt{0, …, 5} the first loop writes before the second reads them; the steps of the loops counted up, for the index \texttt{k} that counts down from \texttt{5}; the value returned, for the scratch integer \texttt{fff}.)
\item \csource{j-points.c}{1968}{1978}: modeled by \lname{JPoints.checkLiftsMpz}
  (deviation of \lname{JPoints.checkLiftsMpz}: integers for the GMP integers, the arguments passed by value, for the pointers; \texttt{x12}, \texttt{x22} and \texttt{x32} locals, for the thread-local scratch integers; the computations of \texttt{b⁶ f(c / b)}, \texttt{A²} and \texttt{B²} functions of their own; a Boolean for the \texttt{int}.)
\item \csource{j-points.c}{1980}{1996}: modeled by \lname{JPoints.checkLiftsMpzA2}
  (deviation of \lname{JPoints.checkLiftsMpzA2}: the value returned, for the scratch integer \texttt{tmp}, without the scratch integers \texttt{tmp2} and \texttt{tmp3}.)
\item \csource{j-points.c}{1997}{1999}: excluded (instrumentation): the value of A\textasciicircum{}2 in check\_lifts\_mpz printed under VERBOSE
\item \csource{j-points.c}{2001}{2002}: modeled by \lname{JPoints.checkLiftsMpz}
  (deviation of \lname{JPoints.checkLiftsMpz}: integers for the GMP integers, the arguments passed by value, for the pointers; \texttt{x12}, \texttt{x22} and \texttt{x32} locals, for the thread-local scratch integers; the computations of \texttt{b⁶ f(c / b)}, \texttt{A²} and \texttt{B²} functions of their own; a Boolean for the \texttt{int}.)
\item \csource{j-points.c}{2004}{2016}: modeled by \lname{JPoints.checkLiftsMpzB2}
  (deviation of \lname{JPoints.checkLiftsMpzB2}: the value returned, for the scratch integer \texttt{tmp}, without the scratch integer \texttt{tmp2}.)
\item \csource{j-points.c}{2017}{2019}: excluded (instrumentation): the value of B\textasciicircum{}2 in check\_lifts\_mpz printed under VERBOSE
\item \csource{j-points.c}{2020}{2022}: modeled by \lname{JPoints.checkLiftsMpz}
  (deviation of \lname{JPoints.checkLiftsMpz}: integers for the GMP integers, the arguments passed by value, for the pointers; \texttt{x12}, \texttt{x22} and \texttt{x32} locals, for the thread-local scratch integers; the computations of \texttt{b⁶ f(c / b)}, \texttt{A²} and \texttt{B²} functions of their own; a Boolean for the \texttt{int}.)
\item \csource{j-points.c}{2031}{2052}: modeled by \lname{JPoints.Printf.gmpSpec}
  (deviation of \lname{JPoints.Printf.gmpSpec}: lists of characters for the characters from \texttt{s} to \texttt{e} and from \texttt{d}, and the list returned for the characters written with their number; the loops over the flags and the width as \texttt{List.takeWhile} and \texttt{List.dropWhile}, the flags \texttt{minus} and \texttt{plus} as the membership of \texttt{-{}} and \texttt{+} in the flags, and the loop over the digits of the precision with the test \texttt{p != e} as the test that what follows the width is empty or a \texttt{.} followed by digits only; the test \texttt{w < e} as a precision that is not empty, and \texttt{e -{} w == 1} as a precision that is a \texttt{.} alone.)
\item \csource{j-points.c}{2054}{2055}: modeled by \lname{JPoints.Printf.printfMpz}
  (deviation of \lname{JPoints.Printf.printfMpz}: nothing printed for a rewritten format that \texttt{gmp\_printf} does not read as an admitted one (the C code passes it to \texttt{gmp\_printf} whatever it is), which no admitted format gives; \texttt{Bool}s for the flags \texttt{reversed} and \texttt{negated}, and arguments for these globals; the coordinates by value, which the C code negates in place in the variables of its caller, which sets them anew before each call; the text, which the caller writes to \texttt{stdout}; the path of one thread.)
\item \csource{j-points.c}{2059}{2059}: modeled by \lname{JPoints.Printf.rewriteFormat}
  (deviation of \lname{JPoints.Printf.rewriteFormat}: the characters of a string for the bytes of the C string, which the scan treats alike, since the characters it looks for are ASCII and no byte of a character beyond ASCII in UTF-8 is; the end of the string for its NUL, the string having none inside (an argument of the command line has none), so that the length \texttt{l} is the number of its characters; the string of the characters built for the buffer \texttt{fmt}, which its NUL ends.)
\item \csource{j-points.c}{2060}{2061}: excluded (memory): the buffer of 2 l + 1 characters for the format that printf\_mpz builds from a format of l characters for gmp\_printf, and error 7 when it cannot be had; the model builds the format as a fresh array (lines 2062-2075)
\item \csource{j-points.c}{2062}{2065}: modeled by \lname{JPoints.Printf.rewriteFormatLoop}
  (deviation of \lname{JPoints.Printf.rewriteFormatLoop}: an array of characters for the C string \texttt{format}, its size for the length \texttt{l} and its end for the NUL that ends the C string, so that \texttt{format[i+1]}, \texttt{format[k]} and \texttt{format[k+1]} past the end are the NUL and the test \texttt{format[k] != 0} is \texttt{k < s.size}; an array of characters for the buffer \texttt{fmt}, empty at the start, its size for the index \texttt{j}, and the result of \texttt{JPoints.Printf.gmpSpec} appended to it, for the characters that \texttt{gmp\_spec} writes and their number.)
\item \csource{j-points.c}{2066}{2066}: modeled by \lname{JPoints.Printf.rewriteFormatScan}
  (deviation of \lname{JPoints.Printf.rewriteFormatScan}: the end of \texttt{s} for the NUL that ends the C string, so that the test \texttt{format[k] != 0} is \texttt{k < s.size}; the membership in \texttt{JPoints.Printf.scanChars} for the test of \texttt{strchr}, which the C code makes only on a character that is not the NUL.)
\item \csource{j-points.c}{2067}{2074}: modeled by \lname{JPoints.Printf.rewriteFormatLoop}
  (deviation of \lname{JPoints.Printf.rewriteFormatLoop}: an array of characters for the C string \texttt{format}, its size for the length \texttt{l} and its end for the NUL that ends the C string, so that \texttt{format[i+1]}, \texttt{format[k]} and \texttt{format[k+1]} past the end are the NUL and the test \texttt{format[k] != 0} is \texttt{k < s.size}; an array of characters for the buffer \texttt{fmt}, empty at the start, its size for the index \texttt{j}, and the result of \texttt{JPoints.Printf.gmpSpec} appended to it, for the characters that \texttt{gmp\_spec} writes and their number.)
\item \csource{j-points.c}{2075}{2075}: modeled by \lname{JPoints.Printf.rewriteFormat}
  (deviation of \lname{JPoints.Printf.rewriteFormat}: the characters of a string for the bytes of the C string, which the scan treats alike, since the characters it looks for are ASCII and no byte of a character beyond ASCII in UTF-8 is; the end of the string for its NUL, the string having none inside (an argument of the command line has none), so that the length \texttt{l} is the number of its characters; the string of the characters built for the buffer \texttt{fmt}, which its NUL ends.)
\item \csource{j-points.c}{2076}{2082}: modeled by \lname{JPoints.Printf.printfMpz}
  (deviation of \lname{JPoints.Printf.printfMpz}: nothing printed for a rewritten format that \texttt{gmp\_printf} does not read as an admitted one (the C code passes it to \texttt{gmp\_printf} whatever it is), which no admitted format gives; \texttt{Bool}s for the flags \texttt{reversed} and \texttt{negated}, and arguments for these globals; the coordinates by value, which the C code negates in place in the variables of its caller, which sets them anew before each call; the text, which the caller writes to \texttt{stdout}; the path of one thread.)
\item \csource{j-points.c}{2083}{2088}: excluded (threads): the output of a point with several threads in printf\_mpz: printed by gmp\_asprintf into a string (error 7 when that fails) and handed to emit; the model prints at once, as one thread does (line 2082)
\item \csource{j-points.c}{2089}{2089}: excluded (memory): the release of the format built for gmp\_printf
\item \csource{j-points.c}{2090}{2090}: modeled by \lname{JPoints.Printf.printfMpz}
  (deviation of \lname{JPoints.Printf.printfMpz}: nothing printed for a rewritten format that \texttt{gmp\_printf} does not read as an admitted one (the C code passes it to \texttt{gmp\_printf} whatever it is), which no admitted format gives; \texttt{Bool}s for the flags \texttt{reversed} and \texttt{negated}, and arguments for these globals; the coordinates by value, which the C code negates in place in the variables of its caller, which sets them anew before each call; the text, which the caller writes to \texttt{stdout}; the path of one thread.)
\item \csource{j-points.c}{2092}{2106}: excluded (threads): emit, which appends the text of a point to the buffer of the unit of work in hand of the thread, enlarging the buffer (error 7 when it cannot), called only by the arms of printf\_mpz and emit\_point for several threads
\item \csource{j-points.c}{2108}{2109}: modeled by \lname{JPoints.Printf.emitPoint}
  (deviation of \lname{JPoints.Printf.emitPoint}: nothing printed for a format that is not admitted (the C code passes any format to \texttt{printf}), which the program refuses before the search; \texttt{Bool}s for the flags \texttt{reversed} and \texttt{negated}, and arguments for these globals and for \texttt{print\_format}; integers for the \texttt{long}s, which are at most \texttt{MAX\_HEIGHT} in absolute value, so that the negations stay in range; the text, which the caller writes to \texttt{stdout}; the path of one thread.)
\item \csource{j-points.c}{2110}{2110}: excluded (threads): the string s of the output of emit\_point with several threads
\item \csource{j-points.c}{2111}{2116}: modeled by \lname{JPoints.Printf.emitPoint}
  (deviation of \lname{JPoints.Printf.emitPoint}: nothing printed for a format that is not admitted (the C code passes any format to \texttt{printf}), which the program refuses before the search; \texttt{Bool}s for the flags \texttt{reversed} and \texttt{negated}, and arguments for these globals and for \texttt{print\_format}; integers for the \texttt{long}s, which are at most \texttt{MAX\_HEIGHT} in absolute value, so that the negations stay in range; the text, which the caller writes to \texttt{stdout}; the path of one thread.)
\item \csource{j-points.c}{2117}{2120}: excluded (threads): the output of a point with several threads in emit\_point: printed by asprintf into a string (error 7 when that fails) and handed to emit; the model prints at once, as one thread does (line 2116)
\item \csource{j-points.c}{2121}{2121}: modeled by \lname{JPoints.Printf.emitPoint}
  (deviation of \lname{JPoints.Printf.emitPoint}: nothing printed for a format that is not admitted (the C code passes any format to \texttt{printf}), which the program refuses before the search; \texttt{Bool}s for the flags \texttt{reversed} and \texttt{negated}, and arguments for these globals and for \texttt{print\_format}; integers for the \texttt{long}s, which are at most \texttt{MAX\_HEIGHT} in absolute value, so that the negations stay in range; the text, which the caller writes to \texttt{stdout}; the path of one thread.)
\item \csource{j-points.c}{2123}{2143}: modeled by \lname{JPoints.checkOnePointFinal}
  (deviation of \lname{JPoints.checkOnePointFinal}: integers for the \texttt{long}s \texttt{a}, \texttt{b}, \texttt{c}, \texttt{m}, \texttt{h} and \texttt{dmax}, and for the GMP integers \texttt{d1} and \texttt{d2}, these by value, for the pointers of the C code; the division \texttt{H / m} and the comparison of \texttt{dbound} with \texttt{MAX\_HEIGHT}, which the C code makes in \texttt{unsigned long}s, made in the integers, where they agree for the \texttt{1 ≤ m ≤ height ≤ MAX\_HEIGHT} and \texttt{dbound ≥ 1} of the calls; \texttt{/} for \texttt{mpz\_divexact}, which divides exactly; locals for the thread-local scratch integers \texttt{cpf1}, \texttt{cpf2} and \texttt{cpf3}; Booleans for the \texttt{int} \texttt{dok} and for the \texttt{int} returned; the output taken and returned, for the thread-local counter \texttt{total} and \texttt{stdout}; the two checks functions of their own.)
\item \csource{j-points.c}{2144}{2146}: modeled by \lname{JPoints.checkOnePointWords}
  (deviation of \lname{JPoints.checkOnePointWords}: integers for the \texttt{long}s \texttt{g}, \texttt{d}, \texttt{a}, \texttt{b} and \texttt{c}, which hold the values: \texttt{mpz\_get\_si} of \texttt{cpf2} and \texttt{cpf1}, within the bounds \texttt{h} and \texttt{dmax}, and the products \texttt{g a}, \texttt{g b} and \texttt{g c}, at most \texttt{m h ≤ MAX\_HEIGHT} in absolute value; \texttt{u} and \texttt{v} by value, for the scratch integers \texttt{cpf1} and \texttt{cpf2}; a Boolean for the \texttt{int} returned; the output taken and returned, for the thread-local counter \texttt{total} and \texttt{stdout}.)
\item \csource{j-points.c}{2147}{2149}: excluded (instrumentation): the coordinates of the point in check\_one\_point\_final printed under VERBOSE
\item \csource{j-points.c}{2150}{2151}: modeled by \lname{JPoints.checkOnePointWords}
  (deviation of \lname{JPoints.checkOnePointWords}: integers for the \texttt{long}s \texttt{g}, \texttt{d}, \texttt{a}, \texttt{b} and \texttt{c}, which hold the values: \texttt{mpz\_get\_si} of \texttt{cpf2} and \texttt{cpf1}, within the bounds \texttt{h} and \texttt{dmax}, and the products \texttt{g a}, \texttt{g b} and \texttt{g c}, at most \texttt{m h ≤ MAX\_HEIGHT} in absolute value; \texttt{u} and \texttt{v} by value, for the scratch integers \texttt{cpf1} and \texttt{cpf2}; a Boolean for the \texttt{int} returned; the output taken and returned, for the thread-local counter \texttt{total} and \texttt{stdout}.)
\item \csource{j-points.c}{2152}{2154}: excluded (instrumentation): the message that the point lifts, under VERBOSE
\item \csource{j-points.c}{2155}{2160}: modeled by \lname{JPoints.checkOnePointWords}
  (deviation of \lname{JPoints.checkOnePointWords}: integers for the \texttt{long}s \texttt{g}, \texttt{d}, \texttt{a}, \texttt{b} and \texttt{c}, which hold the values: \texttt{mpz\_get\_si} of \texttt{cpf2} and \texttt{cpf1}, within the bounds \texttt{h} and \texttt{dmax}, and the products \texttt{g a}, \texttt{g b} and \texttt{g c}, at most \texttt{m h ≤ MAX\_HEIGHT} in absolute value; \texttt{u} and \texttt{v} by value, for the scratch integers \texttt{cpf1} and \texttt{cpf2}; a Boolean for the \texttt{int} returned; the output taken and returned, for the thread-local counter \texttt{total} and \texttt{stdout}.)
\item \csource{j-points.c}{2161}{2163}: excluded (instrumentation): the message that the point does not lift, under VERBOSE
\item \csource{j-points.c}{2164}{2166}: modeled by \lname{JPoints.checkOnePointWords}
  (deviation of \lname{JPoints.checkOnePointWords}: integers for the \texttt{long}s \texttt{g}, \texttt{d}, \texttt{a}, \texttt{b} and \texttt{c}, which hold the values: \texttt{mpz\_get\_si} of \texttt{cpf2} and \texttt{cpf1}, within the bounds \texttt{h} and \texttt{dmax}, and the products \texttt{g a}, \texttt{g b} and \texttt{g c}, at most \texttt{m h ≤ MAX\_HEIGHT} in absolute value; \texttt{u} and \texttt{v} by value, for the scratch integers \texttt{cpf1} and \texttt{cpf2}; a Boolean for the \texttt{int} returned; the output taken and returned, for the thread-local counter \texttt{total} and \texttt{stdout}.)
\item \csource{j-points.c}{2167}{2168}: modeled by \lname{JPoints.checkOnePointFinal}
  (deviation of \lname{JPoints.checkOnePointFinal}: integers for the \texttt{long}s \texttt{a}, \texttt{b}, \texttt{c}, \texttt{m}, \texttt{h} and \texttt{dmax}, and for the GMP integers \texttt{d1} and \texttt{d2}, these by value, for the pointers of the C code; the division \texttt{H / m} and the comparison of \texttt{dbound} with \texttt{MAX\_HEIGHT}, which the C code makes in \texttt{unsigned long}s, made in the integers, where they agree for the \texttt{1 ≤ m ≤ height ≤ MAX\_HEIGHT} and \texttt{dbound ≥ 1} of the calls; \texttt{/} for \texttt{mpz\_divexact}, which divides exactly; locals for the thread-local scratch integers \texttt{cpf1}, \texttt{cpf2} and \texttt{cpf3}; Booleans for the \texttt{int} \texttt{dok} and for the \texttt{int} returned; the output taken and returned, for the thread-local counter \texttt{total} and \texttt{stdout}; the two checks functions of their own.)
\item \csource{j-points.c}{2169}{2172}: modeled by \lname{JPoints.checkOnePointMpz}
  (deviation of \lname{JPoints.checkOnePointMpz}: integers for the \texttt{long}s \texttt{a}, \texttt{b} and \texttt{c} and for the GMP integers; \texttt{u} and \texttt{v} by value, and locals \texttt{cpfa}, \texttt{cpfb} and \texttt{cpfc}, for the thread-local scratch integers; a Boolean for the \texttt{int} returned; the output taken and returned, for the thread-local counter \texttt{total} and \texttt{stdout}.)
\item \csource{j-points.c}{2173}{2175}: excluded (instrumentation): the message that the coordinates exceed a machine word, under VERBOSE
\item \csource{j-points.c}{2176}{2181}: modeled by \lname{JPoints.checkOnePointMpz}
  (deviation of \lname{JPoints.checkOnePointMpz}: integers for the \texttt{long}s \texttt{a}, \texttt{b} and \texttt{c} and for the GMP integers; \texttt{u} and \texttt{v} by value, and locals \texttt{cpfa}, \texttt{cpfb} and \texttt{cpfc}, for the thread-local scratch integers; a Boolean for the \texttt{int} returned; the output taken and returned, for the thread-local counter \texttt{total} and \texttt{stdout}.)
\item \csource{j-points.c}{2182}{2184}: excluded (instrumentation): the message that the point with multi-precision coordinates lifts, under VERBOSE
\item \csource{j-points.c}{2185}{2190}: modeled by \lname{JPoints.checkOnePointMpz}
  (deviation of \lname{JPoints.checkOnePointMpz}: integers for the \texttt{long}s \texttt{a}, \texttt{b} and \texttt{c} and for the GMP integers; \texttt{u} and \texttt{v} by value, and locals \texttt{cpfa}, \texttt{cpfb} and \texttt{cpfc}, for the thread-local scratch integers; a Boolean for the \texttt{int} returned; the output taken and returned, for the thread-local counter \texttt{total} and \texttt{stdout}.)
\item \csource{j-points.c}{2191}{2193}: excluded (instrumentation): the message that the point with multi-precision coordinates does not lift, under VERBOSE
\item \csource{j-points.c}{2194}{2196}: modeled by \lname{JPoints.checkOnePointMpz}
  (deviation of \lname{JPoints.checkOnePointMpz}: integers for the \texttt{long}s \texttt{a}, \texttt{b} and \texttt{c} and for the GMP integers; \texttt{u} and \texttt{v} by value, and locals \texttt{cpfa}, \texttt{cpfb} and \texttt{cpfc}, for the thread-local scratch integers; a Boolean for the \texttt{int} returned; the output taken and returned, for the thread-local counter \texttt{total} and \texttt{stdout}.)
\item \csource{j-points.c}{2197}{2198}: modeled by \lname{JPoints.checkOnePointFinal}
  (deviation of \lname{JPoints.checkOnePointFinal}: integers for the \texttt{long}s \texttt{a}, \texttt{b}, \texttt{c}, \texttt{m}, \texttt{h} and \texttt{dmax}, and for the GMP integers \texttt{d1} and \texttt{d2}, these by value, for the pointers of the C code; the division \texttt{H / m} and the comparison of \texttt{dbound} with \texttt{MAX\_HEIGHT}, which the C code makes in \texttt{unsigned long}s, made in the integers, where they agree for the \texttt{1 ≤ m ≤ height ≤ MAX\_HEIGHT} and \texttt{dbound ≥ 1} of the calls; \texttt{/} for \texttt{mpz\_divexact}, which divides exactly; locals for the thread-local scratch integers \texttt{cpf1}, \texttt{cpf2} and \texttt{cpf3}; Booleans for the \texttt{int} \texttt{dok} and for the \texttt{int} returned; the output taken and returned, for the thread-local counter \texttt{total} and \texttt{stdout}; the two checks functions of their own.)
\item \csource{j-points.c}{2199}{2201}: excluded (instrumentation): the message that the height is too big, under VERBOSE
\item \csource{j-points.c}{2202}{2204}: modeled by \lname{JPoints.checkOnePointFinal}
  (deviation of \lname{JPoints.checkOnePointFinal}: integers for the \texttt{long}s \texttt{a}, \texttt{b}, \texttt{c}, \texttt{m}, \texttt{h} and \texttt{dmax}, and for the GMP integers \texttt{d1} and \texttt{d2}, these by value, for the pointers of the C code; the division \texttt{H / m} and the comparison of \texttt{dbound} with \texttt{MAX\_HEIGHT}, which the C code makes in \texttt{unsigned long}s, made in the integers, where they agree for the \texttt{1 ≤ m ≤ height ≤ MAX\_HEIGHT} and \texttt{dbound ≥ 1} of the calls; \texttt{/} for \texttt{mpz\_divexact}, which divides exactly; locals for the thread-local scratch integers \texttt{cpf1}, \texttt{cpf2} and \texttt{cpf3}; Booleans for the \texttt{int} \texttt{dok} and for the \texttt{int} returned; the output taken and returned, for the thread-local counter \texttt{total} and \texttt{stdout}; the two checks functions of their own.)
\item \csource{j-points.c}{2206}{2210}: modeled by \lname{JPoints.checkOnePoint}
  (deviation of \lname{JPoints.checkOnePoint}: integers for the \texttt{long}s \texttt{a}, \texttt{b} and \texttt{c} and for the GMP integers; locals for the thread-local scratch integers \texttt{tmp}, \texttt{tmp2} and \texttt{ddd}, and \texttt{2 K₂} a local, for \texttt{kummer[2]}, which the C code doubles in place; a Boolean for the \texttt{int} returned; the output taken and returned, for the thread-local counter \texttt{total} and \texttt{stdout}.)
\item \csource{j-points.c}{2211}{2213}: excluded (instrumentation): the triple of check\_one\_point printed under VERBOSE
\item \csource{j-points.c}{2214}{2218}: modeled by \lname{JPoints.checkOnePoint}
  (deviation of \lname{JPoints.checkOnePoint}: integers for the \texttt{long}s \texttt{a}, \texttt{b} and \texttt{c} and for the GMP integers; locals for the thread-local scratch integers \texttt{tmp}, \texttt{tmp2} and \texttt{ddd}, and \texttt{2 K₂} a local, for \texttt{kummer[2]}, which the C code doubles in place; a Boolean for the \texttt{int} returned; the output taken and returned, for the thread-local counter \texttt{total} and \texttt{stdout}.)
\item \csource{j-points.c}{2219}{2223}: excluded (instrumentation): the Kummer equation and its discriminant printed under VERBOSE
\item \csource{j-points.c}{2224}{2225}: modeled by \lname{JPoints.checkOnePoint}
  (deviation of \lname{JPoints.checkOnePoint}: integers for the \texttt{long}s \texttt{a}, \texttt{b} and \texttt{c} and for the GMP integers; locals for the thread-local scratch integers \texttt{tmp}, \texttt{tmp2} and \texttt{ddd}, and \texttt{2 K₂} a local, for \texttt{kummer[2]}, which the C code doubles in place; a Boolean for the \texttt{int} returned; the output taken and returned, for the thread-local counter \texttt{total} and \texttt{stdout}.)
\item \csource{j-points.c}{2226}{2228}: excluded (instrumentation): the message that the discriminant is 0, under VERBOSE
\item \csource{j-points.c}{2229}{2230}: modeled by \lname{JPoints.checkOnePoint}
  (deviation of \lname{JPoints.checkOnePoint}: integers for the \texttt{long}s \texttt{a}, \texttt{b} and \texttt{c} and for the GMP integers; locals for the thread-local scratch integers \texttt{tmp}, \texttt{tmp2} and \texttt{ddd}, and \texttt{2 K₂} a local, for \texttt{kummer[2]}, which the C code doubles in place; a Boolean for the \texttt{int} returned; the output taken and returned, for the thread-local counter \texttt{total} and \texttt{stdout}.)
\item \csource{j-points.c}{2231}{2233}: excluded (instrumentation): the message that the Kummer equation is constant, under VERBOSE
\item \csource{j-points.c}{2234}{2241}: modeled by \lname{JPoints.checkOnePoint}
  (deviation of \lname{JPoints.checkOnePoint}: integers for the \texttt{long}s \texttt{a}, \texttt{b} and \texttt{c} and for the GMP integers; locals for the thread-local scratch integers \texttt{tmp}, \texttt{tmp2} and \texttt{ddd}, and \texttt{2 K₂} a local, for \texttt{kummer[2]}, which the C code doubles in place; a Boolean for the \texttt{int} returned; the output taken and returned, for the thread-local counter \texttt{total} and \texttt{stdout}.)
\item \csource{j-points.c}{2242}{2244}: excluded (instrumentation): the message that the discriminant is a square, under VERBOSE
\item \csource{j-points.c}{2245}{2246}: modeled by \lname{JPoints.checkOnePoint}
  (deviation of \lname{JPoints.checkOnePoint}: integers for the \texttt{long}s \texttt{a}, \texttt{b} and \texttt{c} and for the GMP integers; locals for the thread-local scratch integers \texttt{tmp}, \texttt{tmp2} and \texttt{ddd}, and \texttt{2 K₂} a local, for \texttt{kummer[2]}, which the C code doubles in place; a Boolean for the \texttt{int} returned; the output taken and returned, for the thread-local counter \texttt{total} and \texttt{stdout}.)
\item \csource{j-points.c}{2247}{2249}: excluded (instrumentation): the message that the leading coefficient is 0, under VERBOSE
\item \csource{j-points.c}{2250}{2257}: modeled by \lname{JPoints.checkOnePoint}
  (deviation of \lname{JPoints.checkOnePoint}: integers for the \texttt{long}s \texttt{a}, \texttt{b} and \texttt{c} and for the GMP integers; locals for the thread-local scratch integers \texttt{tmp}, \texttt{tmp2} and \texttt{ddd}, and \texttt{2 K₂} a local, for \texttt{kummer[2]}, which the C code doubles in place; a Boolean for the \texttt{int} returned; the output taken and returned, for the thread-local counter \texttt{total} and \texttt{stdout}.)
\item \csource{j-points.c}{2258}{2262}: excluded (instrumentation): the square root of the discriminant and the first fourth coordinate printed under VERBOSE
\item \csource{j-points.c}{2263}{2266}: modeled by \lname{JPoints.checkOnePoint}
  (deviation of \lname{JPoints.checkOnePoint}: integers for the \texttt{long}s \texttt{a}, \texttt{b} and \texttt{c} and for the GMP integers; locals for the thread-local scratch integers \texttt{tmp}, \texttt{tmp2} and \texttt{ddd}, and \texttt{2 K₂} a local, for \texttt{kummer[2]}, which the C code doubles in place; a Boolean for the \texttt{int} returned; the output taken and returned, for the thread-local counter \texttt{total} and \texttt{stdout}.)
\item \csource{j-points.c}{2267}{2271}: excluded (instrumentation): the square root of the discriminant and the second fourth coordinate printed under VERBOSE
\item \csource{j-points.c}{2272}{2276}: modeled by \lname{JPoints.checkOnePoint}
  (deviation of \lname{JPoints.checkOnePoint}: integers for the \texttt{long}s \texttt{a}, \texttt{b} and \texttt{c} and for the GMP integers; locals for the thread-local scratch integers \texttt{tmp}, \texttt{tmp2} and \texttt{ddd}, and \texttt{2 K₂} a local, for \texttt{kummer[2]}, which the C code doubles in place; a Boolean for the \texttt{int} returned; the output taken and returned, for the thread-local counter \texttt{total} and \texttt{stdout}.)
\item \csource{j-points.c}{2277}{2279}: excluded (instrumentation): the message that the discriminant is not a square, under VERBOSE
\item \csource{j-points.c}{2280}{2283}: modeled by \lname{JPoints.checkOnePoint}
  (deviation of \lname{JPoints.checkOnePoint}: integers for the \texttt{long}s \texttt{a}, \texttt{b} and \texttt{c} and for the GMP integers; locals for the thread-local scratch integers \texttt{tmp}, \texttt{tmp2} and \texttt{ddd}, and \texttt{2 K₂} a local, for \texttt{kummer[2]}, which the C code doubles in place; a Boolean for the \texttt{int} returned; the output taken and returned, for the thread-local counter \texttt{total} and \texttt{stdout}.)
\item \csource{j-points.c}{2290}{2319}: modeled by \lname{JPoints.Printf.printPoly}
  (deviation of \lname{JPoints.Printf.printPoly}: the text returned, which the caller writes to \texttt{stdout}, and of which the flush of \texttt{stdout} at the end changes nothing; natural numbers for the \texttt{long} \texttt{degree} and the \texttt{int} index \texttt{i}; the steps of the loop counted up, for the index \texttt{i} that counts down from \texttt{degree}; a local for the scratch integer \texttt{tmp3}.)
\item \csource{j-points.c}{2321}{2326}: modeled by \lname{JPoints.message}
  (deviation of \lname{JPoints.message}: the text returned, for the output on \texttt{stdout} that \texttt{fflush} flushes; the \texttt{long} \texttt{n} an integer; the \texttt{long} \texttt{total} an integer, which message \texttt{6} prints, and message \texttt{5} passes as a natural number to \texttt{print\_poly}, as the degree, at least \texttt{5}, that \texttt{main} passes; the global \texttt{coeffs} read from the curve of \texttt{g.choice}; \texttt{num\_primes}, which \texttt{read\_input} leaves in \texttt{[1, 53]}, as a natural number; the messages \texttt{4}, \texttt{7} and \texttt{8} by the functions of their cases.)
\item \csource{j-points.c}{2327}{2337}: modeled by \lname{JPoints.message4Mode}
  (deviation of \lname{JPoints.message4Mode}: the text returned, for the output on \texttt{stdout}; the doubles by \texttt{JPoints.Double.ofFloat}, which takes a NaN to \texttt{0x7ff8000000000000} and would print \texttt{nan} for \texttt{-{}nan} (the statistics are quotients of sums of counts by a positive number of rows, or \texttt{1}, never NaNs).)
\item \csource{j-points.c}{2338}{2347}: modeled by \lname{JPoints.message4}
  (deviation of \lname{JPoints.message4}: the text returned, for the output on \texttt{stdout}; the fractions by \texttt{JPoints.Double.ofFloat}, which takes a NaN to \texttt{0x7ff8000000000000} and would print \texttt{nan} for \texttt{-{}nan} (they are quotients by positive counts, never NaNs).)
\item \csource{j-points.c}{2348}{2356}: modeled by \lname{JPoints.message4Costs}
  (deviation of \lname{JPoints.message4Costs}: the text returned, for the output on \texttt{stdout}; a natural number for the \texttt{size\_t} index; the values by \texttt{JPoints.Double.ofFloat}, which takes a NaN to \texttt{0x7ff8000000000000} and would print \texttt{nan} for \texttt{-{}nan} (\texttt{-{}c} sets only positive normal values, never NaNs).)
\item \csource{j-points.c}{2357}{2357}: modeled by \lname{JPoints.message4}
  (deviation of \lname{JPoints.message4}: the text returned, for the output on \texttt{stdout}; the fractions by \texttt{JPoints.Double.ofFloat}, which takes a NaN to \texttt{0x7ff8000000000000} and would print \texttt{nan} for \texttt{-{}nan} (they are quotients by positive counts, never NaNs).)
\item \csource{j-points.c}{2358}{2362}: modeled by \lname{JPoints.message}
  (deviation of \lname{JPoints.message}: the text returned, for the output on \texttt{stdout} that \texttt{fflush} flushes; the \texttt{long} \texttt{n} an integer; the \texttt{long} \texttt{total} an integer, which message \texttt{6} prints, and message \texttt{5} passes as a natural number to \texttt{print\_poly}, as the degree, at least \texttt{5}, that \texttt{main} passes; the global \texttt{coeffs} read from the curve of \texttt{g.choice}; \texttt{num\_primes}, which \texttt{read\_input} leaves in \texttt{[1, 53]}, as a natural number; the messages \texttt{4}, \texttt{7} and \texttt{8} by the functions of their cases.)
\item \csource{j-points.c}{2363}{2381}: modeled by \lname{JPoints.message7}
  (deviation of \lname{JPoints.message7}: the text returned, for the output on \texttt{stdout}; the index \texttt{i} that the three loops share as their bounds: each loop starts where the one before ended, at the larger of the numbers before, which is the number of the stages before it for a plan of \texttt{JPoints.choosePrimes}.)
\item \csource{j-points.c}{2382}{2401}: modeled by \lname{JPoints.message8}
  (deviation of \lname{JPoints.message8}: the text returned, for the output on \texttt{stdout}; natural numbers for the \texttt{long} indices; the linear search for the record of a prime bounded by the size of \texttt{prec}, which holds it (its records hold the indices of all primes considered), where the C code searches without a bound.)
\item \csource{j-points.c}{2402}{2408}: modeled by \lname{JPoints.message}
  (deviation of \lname{JPoints.message}: the text returned, for the output on \texttt{stdout} that \texttt{fflush} flushes; the \texttt{long} \texttt{n} an integer; the \texttt{long} \texttt{total} an integer, which message \texttt{6} prints, and message \texttt{5} passes as a natural number to \texttt{print\_poly}, as the degree, at least \texttt{5}, that \texttt{main} passes; the global \texttt{coeffs} read from the curve of \texttt{g.choice}; \texttt{num\_primes}, which \texttt{read\_input} leaves in \texttt{[1, 53]}, as a natural number; the messages \texttt{4}, \texttt{7} and \texttt{8} by the functions of their cases.)
\item \csource{j-points.c}{2410}{2412}: modeled by \lname{JPoints.Cli.Failure.message}
  (deviation of \lname{JPoints.Cli.Failure.message}: the text returned, for the output on \texttt{stdout}; the failure \texttt{f}, for the \texttt{long} \texttt{n}, with only the errors that a command line can cause (the messages of errors 1, 7 and 8 are excluded).)
\item \csource{j-points.c}{2413}{2413}: modeled by \lname{JPoints.Cli.Failure.message}
  (deviation of \lname{JPoints.Cli.Failure.message}: the text returned, for the output on \texttt{stdout}; the failure \texttt{f}, for the \texttt{long} \texttt{n}, with only the errors that a command line can cause (the messages of errors 1, 7 and 8 are excluded).)
\item \csource{j-points.c}{2414}{2416}: excluded (other): the message of error 1, for an unsigned long of neither 16, 32 nor 64 bits, which only the test of the platform in main calls (line 341); its label case 1 shares line 2413 with the brace of the switch, which stays with the model
\item \csource{j-points.c}{2417}{2421}: modeled by \lname{JPoints.Cli.Failure.message}
  (deviation of \lname{JPoints.Cli.Failure.message}: the text returned, for the output on \texttt{stdout}; the failure \texttt{f}, for the \texttt{long} \texttt{n}, with only the errors that a command line can cause (the messages of errors 1, 7 and 8 are excluded).)
\item \csource{j-points.c}{2422}{2422}: excluded (memory): the message of error 7, for lack of memory, which only failed allocations call; the model allocates nothing itself and cannot meet it
\item \csource{j-points.c}{2423}{2423}: excluded (threads): the message of error 8, which only the start of the threads calls (line 1408)
\item \csource{j-points.c}{2424}{2424}: modeled by \lname{JPoints.Cli.Failure.message}
  (deviation of \lname{JPoints.Cli.Failure.message}: the text returned, for the output on \texttt{stdout}; the failure \texttt{f}, for the \texttt{long} \texttt{n}, with only the errors that a command line can cause (the messages of errors 1, 7 and 8 are excluded).)
\item \csource{j-points.c}{2425}{2432}: modeled by \lname{JPoints.Cli.costSyntaxMessage}
  (deviation of \lname{JPoints.Cli.costSyntaxMessage}: the text returned, for the output on \texttt{stdout}; a natural number for the \texttt{size\_t} index \texttt{k}.)
\item \csource{j-points.c}{2433}{2435}: modeled by \lname{JPoints.Cli.Failure.message}
  (deviation of \lname{JPoints.Cli.Failure.message}: the text returned, for the output on \texttt{stdout}; the failure \texttt{f}, for the \texttt{long} \texttt{n}, with only the errors that a command line can cause (the messages of errors 1, 7 and 8 are excluded).)
\item \csource{j-points.c}{2436}{2439}: modeled by \lname{JPoints.Cli.usage}
  (deviation of \lname{JPoints.Cli.usage}: the text returned, for the output of the four \texttt{printf}s.)
\item \csource{j-points.c}{2440}{2441}: modeled by \lname{JPoints.Cli.Failure.message}
  (deviation of \lname{JPoints.Cli.Failure.message}: the text returned, for the output on \texttt{stdout}; the failure \texttt{f}, for the \texttt{long} \texttt{n}, with only the errors that a command line can cause (the messages of errors 1, 7 and 8 are excluded).)
\item \csource{j-points.c}{2442}{2444}: modeled by \lname{JPoints.Cli.Failure.status}
  (deviation of \lname{JPoints.Cli.Failure.status}: the status returned, for \texttt{exit(n)}; the flush of \texttt{stdout} made by \texttt{JPoints.Program.main} when it writes the text.)
\end{itemize}

\section*{\texttt{j-sift.c}}
\begin{itemize}
\item \csource{j-sift.c}{31}{36}: excluded (other): the includes of the headers of the C library, of GMP and of j-points.h; the model has no headers
\item \csource{j-sift.c}{47}{53}: modeled by \lname{JPoints.WalkSpec}
  (deviation of \lname{JPoints.WalkSpec}: natural numbers for the \texttt{long}s, which the C code keeps nonnegative.)
\item \csource{j-sift.c}{55}{56}: modeled by \lname{JPoints.ChunkSpec}
  (deviation of \lname{JPoints.ChunkSpec}: arrays of their own for the pointers \texttt{walk} and \texttt{res} into the arrays that \texttt{init\_sift} allocates for all the chunks, the parts of this chunk; integers for the \texttt{long}s \texttt{w\_low} and \texttt{w\_high}, which may be negative, and natural numbers for the residues; without the field \texttt{nvec}, which only the vector walk reads.)
\item \csource{j-sift.c}{57}{57}: excluded (other): the field nvec of a chunk, the number of vectors of VW words that cover it, which only the vector walk (VW > 1) reads (lines 1006, 1064)
\item \csource{j-sift.c}{58}{60}: modeled by \lname{JPoints.ChunkSpec}
  (deviation of \lname{JPoints.ChunkSpec}: arrays of their own for the pointers \texttt{walk} and \texttt{res} into the arrays that \texttt{init\_sift} allocates for all the chunks, the parts of this chunk; integers for the \texttt{long}s \texttt{w\_low} and \texttt{w\_high}, which may be negative, and natural numbers for the residues; without the field \texttt{nvec}, which only the vector walk reads.)
\item \csource{j-sift.c}{67}{71}: modeled by \lname{JPoints.sievePeriod}
  (deviation of \lname{JPoints.sievePeriod}: \texttt{1} for the macro \texttt{VW}, the width of the walk in the scalar build modeled; natural numbers for the \texttt{long}s \texttt{p} and \texttt{P}, here \texttt{period}.)
\item \csource{j-sift.c}{72}{72}: modeled by \lname{JPoints.sieveRowlen}
  (deviation of \lname{JPoints.sieveRowlen}: \texttt{1} for the macro \texttt{VW}, the width of the walk in the scalar build modeled; natural numbers for the \texttt{long}s.)
\item \csource{j-sift.c}{73}{77}: excluded (other): the vector types vec and vecu of gcc's vector extension for the vector walk (VW > 1)
\item \csource{j-sift.c}{79}{80}: modeled by \lname{JPoints.SiftData}, \lname{JPoints.SiftData.initial}
  (deviation of \lname{JPoints.SiftData}: arrays for the arrays of the C code and for the pointers \texttt{chunks} and \texttt{res0} to the arrays that \texttt{init\_sift} allocates, the empty array for a null pointer; natural numbers for the nonnegative \texttt{long}s and for the \texttt{unsigned short}s of \texttt{res0}, integers for the \texttt{long}s \texttt{w\_low\_all} and \texttt{w\_high\_all}; an array of Booleans for the \texttt{unsigned char}s of \texttt{all\_ones}, which hold \texttt{0} and \texttt{1}; a structure, for the statics.)
\item \csource{j-sift.c}{82}{83}: modeled by \lname{JPoints.SiftData}, \lname{JPoints.SiftData.initial}
  (deviation of \lname{JPoints.SiftData}: arrays for the arrays of the C code and for the pointers \texttt{chunks} and \texttt{res0} to the arrays that \texttt{init\_sift} allocates, the empty array for a null pointer; natural numbers for the nonnegative \texttt{long}s and for the \texttt{unsigned short}s of \texttt{res0}, integers for the \texttt{long}s \texttt{w\_low\_all} and \texttt{w\_high\_all}; an array of Booleans for the \texttt{unsigned char}s of \texttt{all\_ones}, which hold \texttt{0} and \texttt{1}; a structure, for the statics.)
\item \csource{j-sift.c}{97}{97}: modeled by \lname{JPoints.SiftData}, \lname{JPoints.SiftData.initial}
  (deviation of \lname{JPoints.SiftData}: arrays for the arrays of the C code and for the pointers \texttt{chunks} and \texttt{res0} to the arrays that \texttt{init\_sift} allocates, the empty array for a null pointer; natural numbers for the nonnegative \texttt{long}s and for the \texttt{unsigned short}s of \texttt{res0}, integers for the \texttt{long}s \texttt{w\_low\_all} and \texttt{w\_high\_all}; an array of Booleans for the \texttt{unsigned char}s of \texttt{all\_ones}, which hold \texttt{0} and \texttt{1}; a structure, for the statics.)
\item \csource{j-sift.c}{98}{99}: modeled by \lname{JPoints.Row3Cache}, \lname{JPoints.Row3Cache.initial}
  (deviation of \lname{JPoints.Row3Cache}: arrays for the arrays of the C code, \texttt{JPoints.Tab3Ptr} for the pointers of \texttt{tab3}; integers for the \texttt{long}s; a structure, for the thread-local globals of the one thread.)
\item \csource{j-sift.c}{100}{100}: modeled by \lname{JPoints.SiftData}, \lname{JPoints.SiftData.initial}
  (deviation of \lname{JPoints.SiftData}: arrays for the arrays of the C code and for the pointers \texttt{chunks} and \texttt{res0} to the arrays that \texttt{init\_sift} allocates, the empty array for a null pointer; natural numbers for the nonnegative \texttt{long}s and for the \texttt{unsigned short}s of \texttt{res0}, integers for the \texttt{long}s \texttt{w\_low\_all} and \texttt{w\_high\_all}; an array of Booleans for the \texttt{unsigned char}s of \texttt{all\_ones}, which hold \texttt{0} and \texttt{1}; a structure, for the statics.)
\item \csource{j-sift.c}{105}{105}: modeled by \lname{JPoints.mod3}
  (deviation of \lname{JPoints.mod3}: the static data and \texttt{sieve\_primes2} passed, for the globals that the macro reads; integers for the \texttt{long}s \texttt{x} and \texttt{pr[…]}, with the conversion of \texttt{x} to an \texttt{unsigned long} modulo \texttt{2 \textasciicircum{} 64} as C converts it; the quotient is below \texttt{2 \textasciicircum{} 32}, so that the conversion back to a \texttt{long} keeps it; where \texttt{x} is an \texttt{unsigned long}, the subtraction is made modulo \texttt{2 \textasciicircum{} 64}, which gives the integer for the operands of the calls, below \texttt{2 \textasciicircum{} 16}, whose remainder is nonnegative.)
\item \csource{j-sift.c}{106}{106}: modeled by \lname{JPoints.Row3Cache}, \lname{JPoints.Row3Cache.initial}
  (deviation of \lname{JPoints.Row3Cache}: arrays for the arrays of the C code, \texttt{JPoints.Tab3Ptr} for the pointers of \texttt{tab3}; integers for the \texttt{long}s; a structure, for the thread-local globals of the one thread.)
\item \csource{j-sift.c}{111}{111}: modeled by \lname{JPoints.SiftData}, \lname{JPoints.SiftData.initial}
  (deviation of \lname{JPoints.SiftData}: arrays for the arrays of the C code and for the pointers \texttt{chunks} and \texttt{res0} to the arrays that \texttt{init\_sift} allocates, the empty array for a null pointer; natural numbers for the nonnegative \texttt{long}s and for the \texttt{unsigned short}s of \texttt{res0}, integers for the \texttt{long}s \texttt{w\_low\_all} and \texttt{w\_high\_all}; an array of Booleans for the \texttt{unsigned char}s of \texttt{all\_ones}, which hold \texttt{0} and \texttt{1}; a structure, for the statics.)
\item \csource{j-sift.c}{112}{112}: excluded (other): the array period of the periods of the walks, sieve\_period(p) for each sieving prime p, which only the vector walk (VW > 1) reads (lines 186-187, 910-911, 944); with VW = 1 the period is p
\item \csource{j-sift.c}{113}{113}: modeled by \lname{JPoints.SiftState}, \lname{JPoints.SiftState.initial}
  (deviation of \lname{JPoints.SiftState}: an array for the bit array that \texttt{survivors} points to; indices into the sieve tables for the pointers \texttt{rowptr[n]} into \texttt{sieve\_tab[n]}; natural numbers for the counters, which start at \texttt{0} and count up, and for the index \texttt{cur\_chunk}; a structure, for the thread-local globals of the one thread.)
  (deviation of \lname{JPoints.SiftState.initial}: a fresh array of \texttt{arraySize + VW} zeros, \texttt{VW = 1}, for the bit array that \texttt{init\_thread\_sieve} allocates (memory management, not modeled), every word of which the sieve writes before it reads it; index \texttt{0} for the null pointers of \texttt{rowptr}, which the sieve reads only after \texttt{row\_setup} has set them.)
\item \csource{j-sift.c}{114}{116}: modeled by \lname{JPoints.SiftData}, \lname{JPoints.SiftData.initial}
  (deviation of \lname{JPoints.SiftData}: arrays for the arrays of the C code and for the pointers \texttt{chunks} and \texttt{res0} to the arrays that \texttt{init\_sift} allocates, the empty array for a null pointer; natural numbers for the nonnegative \texttt{long}s and for the \texttt{unsigned short}s of \texttt{res0}, integers for the \texttt{long}s \texttt{w\_low\_all} and \texttt{w\_high\_all}; an array of Booleans for the \texttt{unsigned char}s of \texttt{all\_ones}, which hold \texttt{0} and \texttt{1}; a structure, for the statics.)
\item \csource{j-sift.c}{120}{120}: modeled by \lname{JPoints.tubeBand}
\item \csource{j-sift.c}{121}{121}: modeled by \lname{JPoints.tubeRanges}
\item \csource{j-sift.c}{122}{124}: modeled by \lname{JPoints.TubeBands}, \lname{JPoints.TubeBands.initial}
  (deviation of \lname{JPoints.TubeBands}: arrays for the arrays that \texttt{band\_nr} and \texttt{band\_rg} point to, one flat array of \texttt{2 TUBE\_RANGES TUBE\_BAND} entries for the array of the rows of \texttt{band\_rg}, the \texttt{k}-th range of the \texttt{row}-th row at the indices \texttt{2 TUBE\_RANGES row + 2 k} and \texttt{2 TUBE\_RANGES row + 2 k + 1}; integers for the \texttt{long}s of the first row and of the ends of the ranges, natural numbers for the counts; a structure, for the thread-local globals of the one thread.)
  (deviation of \lname{JPoints.TubeBands.initial}: fresh arrays of zeros, for the arrays that \texttt{init\_thread\_sieve} allocates (memory management, not modeled), whose pointers start as null pointers; the C code sets each count of a band to \texttt{0} before it reads it and reads only the ranges that the count covers.)
\item \csource{j-sift.c}{125}{125}: modeled by \lname{JPoints.TubeBands.blo}, \lname{JPoints.TubeBands.setBlo}
  (deviation of \lname{JPoints.TubeBands.blo}: the index \texttt{2 TUBE\_RANGES row + 2 k} of the flat array \texttt{band\_rg}.)
  (deviation of \lname{JPoints.TubeBands.setBlo}: the index \texttt{2 TUBE\_RANGES row + 2 k} of the flat array \texttt{band\_rg}; the bands returned.)
\item \csource{j-sift.c}{126}{126}: modeled by \lname{JPoints.TubeBands.bhi}, \lname{JPoints.TubeBands.setBhi}
  (deviation of \lname{JPoints.TubeBands.bhi}: the index \texttt{2 TUBE\_RANGES row + 2 k + 1} of the flat array \texttt{band\_rg}.)
  (deviation of \lname{JPoints.TubeBands.setBhi}: the index \texttt{2 TUBE\_RANGES row + 2 k + 1} of the flat array \texttt{band\_rg}; the bands returned.)
\item \csource{j-sift.c}{133}{142}: modeled by \lname{JPoints.initSift}
  (deviation of \lname{JPoints.initSift}: the static data taken and returned, for the statics; the globals that it reads passed; \texttt{>>>} on integers for C's \texttt{>>} on the \texttt{long}s, which gcc shifts arithmetically also for the negative \texttt{-{}height}, where ISO C leaves the result to the implementation; \texttt{JPoints.ceilDiv} for the macro \texttt{CEIL}, and a natural number for \texttt{num\_chunks}, which is positive; \texttt{JPoints.Double.ofInt} for the conversion of the \texttt{long} \texttt{dbound} to a double, rounded to nearest; a fresh array for the array of the chunks that \texttt{init\_sift} allocates (\texttt{j-{}sift.c:157}, memory management, not modeled), filled in the order of the chunks.)
\item \csource{j-sift.c}{143}{144}: modeled by \lname{JPoints.initSiftPrimes}
  (deviation of \lname{JPoints.initSiftPrimes}: the static data taken and returned, for the statics; the globals that it reads passed; natural numbers for the nonnegative \texttt{long}s, and \texttt{UInt64}s for the \texttt{unsigned long}s of the reciprocals.)
\item \csource{j-sift.c}{145}{145}: excluded (other): the period of the walk of each sieving prime stored in the array period, which only the vector walk (VW > 1) reads
\item \csource{j-sift.c}{146}{153}: modeled by \lname{JPoints.initSiftPrimes}
  (deviation of \lname{JPoints.initSiftPrimes}: the static data taken and returned, for the statics; the globals that it reads passed; natural numbers for the nonnegative \texttt{long}s, and \texttt{UInt64}s for the \texttt{unsigned long}s of the reciprocals.)
\item \csource{j-sift.c}{154}{155}: modeled by \lname{JPoints.initSift}
  (deviation of \lname{JPoints.initSift}: the static data taken and returned, for the statics; the globals that it reads passed; \texttt{>>>} on integers for C's \texttt{>>} on the \texttt{long}s, which gcc shifts arithmetically also for the negative \texttt{-{}height}, where ISO C leaves the result to the implementation; \texttt{JPoints.ceilDiv} for the macro \texttt{CEIL}, and a natural number for \texttt{num\_chunks}, which is positive; \texttt{JPoints.Double.ofInt} for the conversion of the \texttt{long} \texttt{dbound} to a double, rounded to nearest; a fresh array for the array of the chunks that \texttt{init\_sift} allocates (\texttt{j-{}sift.c:157}, memory management, not modeled), filled in the order of the chunks.)
\item \csource{j-sift.c}{156}{156}: excluded (memory): the second allocation of the state of the calling thread (init\_thread\_sieve), a new bit array and new lists of the tube's bands in place of those that tube\_sample allocated (line 738), which it frees; the model makes the bit array as a fresh array (JPoints.SiftState.initial) and keeps the lists of the bands that tube\_sample wrote, which nothing reads before it writes them (each count of a band is set to 0 before it is read, and only the ranges that a count covers are read)
\item \csource{j-sift.c}{157}{162}: excluded (memory): the allocation of the chunks, of their walks and residues and of the table res0 on the heap (the last three one entry longer, so that none is empty), and error 7 when one of them cannot be had; the model makes them as fresh arrays, and the loops below write every entry that is read
\item \csource{j-sift.c}{163}{163}: modeled by \lname{JPoints.initSift}
  (deviation of \lname{JPoints.initSift}: the static data taken and returned, for the statics; the globals that it reads passed; \texttt{>>>} on integers for C's \texttt{>>} on the \texttt{long}s, which gcc shifts arithmetically also for the negative \texttt{-{}height}, where ISO C leaves the result to the implementation; \texttt{JPoints.ceilDiv} for the macro \texttt{CEIL}, and a natural number for \texttt{num\_chunks}, which is positive; \texttt{JPoints.Double.ofInt} for the conversion of the \texttt{long} \texttt{dbound} to a double, rounded to nearest; a fresh array for the array of the chunks that \texttt{init\_sift} allocates (\texttt{j-{}sift.c:157}, memory management, not modeled), filled in the order of the chunks.)
\item \csource{j-sift.c}{164}{169}: modeled by \lname{JPoints.initSiftChunk}
  (deviation of \lname{JPoints.initSiftChunk}: the chunk returned, for the entry of the array \texttt{chunks} that the C code writes; the globals and the static data that it reads passed; arrays of its own, filled in the order in which the C code writes them, for the parts of the arrays of the walks and of the residues of all the chunks that the pointers \texttt{walk} and \texttt{res} of the chunk point to; integers for the \texttt{long}s that may be negative, with \texttt{Int.tmod} for C's \texttt{\%}, and natural numbers for the nonnegative ones.)
\item \csource{j-sift.c}{170}{170}: excluded (other): the number of vectors of VW words that cover the chunk, which only the vector walk (VW > 1) reads
\item \csource{j-sift.c}{171}{174}: modeled by \lname{JPoints.initSiftChunk}
  (deviation of \lname{JPoints.initSiftChunk}: the chunk returned, for the entry of the array \texttt{chunks} that the C code writes; the globals and the static data that it reads passed; arrays of its own, filled in the order in which the C code writes them, for the parts of the arrays of the walks and of the residues of all the chunks that the pointers \texttt{walk} and \texttt{res} of the chunk point to; integers for the \texttt{long}s that may be negative, with \texttt{Int.tmod} for C's \texttt{\%}, and natural numbers for the nonnegative ones.)
\item \csource{j-sift.c}{175}{184}: modeled by \lname{JPoints.initSiftWalk}
  (deviation of \lname{JPoints.initSiftWalk}: the walk returned, for the entry of the array of the walks that the C code writes; the arm \texttt{VW == 1} of the test of the constant \texttt{VW}; integers for the \texttt{long}s \texttt{w\_low} of the chunk and \texttt{start}, with \texttt{Int.tmod} for C's \texttt{\%}, and natural numbers for the nonnegative \texttt{long}s.)
\item \csource{j-sift.c}{185}{190}: excluded (other): the arm of the vector walk (VW > 1) of the walk of a first-stage prime over the chunk, an if on the constant VW that the build modeled never takes
\item \csource{j-sift.c}{191}{196}: modeled by \lname{JPoints.initSiftChunk}
  (deviation of \lname{JPoints.initSiftChunk}: the chunk returned, for the entry of the array \texttt{chunks} that the C code writes; the globals and the static data that it reads passed; arrays of its own, filled in the order in which the C code writes them, for the parts of the arrays of the walks and of the residues of all the chunks that the pointers \texttt{walk} and \texttt{res} of the chunk point to; integers for the \texttt{long}s that may be negative, with \texttt{Int.tmod} for C's \texttt{\%}, and natural numbers for the nonnegative ones.)
\item \csource{j-sift.c}{197}{197}: modeled by \lname{JPoints.initSift}
  (deviation of \lname{JPoints.initSift}: the static data taken and returned, for the statics; the globals that it reads passed; \texttt{>>>} on integers for C's \texttt{>>} on the \texttt{long}s, which gcc shifts arithmetically also for the negative \texttt{-{}height}, where ISO C leaves the result to the implementation; \texttt{JPoints.ceilDiv} for the macro \texttt{CEIL}, and a natural number for \texttt{num\_chunks}, which is positive; \texttt{JPoints.Double.ofInt} for the conversion of the \texttt{long} \texttt{dbound} to a double, rounded to nearest; a fresh array for the array of the chunks that \texttt{init\_sift} allocates (\texttt{j-{}sift.c:157}, memory management, not modeled), filled in the order of the chunks.)
\item \csource{j-sift.c}{198}{202}: modeled by \lname{JPoints.initSiftResidues}
  (deviation of \lname{JPoints.initSiftResidues}: the static data taken and returned, for the statics; the globals that it reads passed; \texttt{res0} a fresh array of \texttt{(array\_size + 1) npr} zeros, for the array that \texttt{init\_sift} allocates (\texttt{j-{}sift.c:160}, memory management, not modeled), every entry of which the loop writes; natural numbers for the nonnegative \texttt{long}s and for the \texttt{unsigned short}s of \texttt{res0}; \texttt{JPoints.Double.ofInt} for the exact conversion of \texttt{p} to a double.)
\item \csource{j-sift.c}{203}{204}: modeled by \lname{JPoints.initSift}
  (deviation of \lname{JPoints.initSift}: the static data taken and returned, for the statics; the globals that it reads passed; \texttt{>>>} on integers for C's \texttt{>>} on the \texttt{long}s, which gcc shifts arithmetically also for the negative \texttt{-{}height}, where ISO C leaves the result to the implementation; \texttt{JPoints.ceilDiv} for the macro \texttt{CEIL}, and a natural number for \texttt{num\_chunks}, which is positive; \texttt{JPoints.Double.ofInt} for the conversion of the \texttt{long} \texttt{dbound} to a double, rounded to nearest; a fresh array for the array of the chunks that \texttt{init\_sift} allocates (\texttt{j-{}sift.c:157}, memory management, not modeled), filled in the order of the chunks.)
\item \csource{j-sift.c}{206}{217}: excluded (memory): init\_thread\_sieve, which frees the arrays of an earlier call of the thread and allocates the bit array of a thread (aligned and padded by VW words for the vector walk) and its lists of the tube's bands, with error 7 when it cannot; the model makes them as fresh arrays (each word of the bit array is written before it is read, each count of a band set to 0 before it is read)
\item \csource{j-sift.c}{222}{222}: modeled by \lname{JPoints.Row3Cache}, \lname{JPoints.Row3Cache.initial}
  (deviation of \lname{JPoints.Row3Cache}: arrays for the arrays of the C code, \texttt{JPoints.Tab3Ptr} for the pointers of \texttt{tab3}; integers for the \texttt{long}s; a structure, for the thread-local globals of the one thread.)
\item \csource{j-sift.c}{223}{226}: modeled by \lname{JPoints.row3Setup}
  (deviation of \lname{JPoints.row3Setup}: the data taken and returned, for the thread-local globals; the globals that it reads passed; integers for the \texttt{long}s \texttt{a} and \texttt{b}.)
\item \csource{j-sift.c}{227}{231}: modeled by \lname{JPoints.row3SetupInverses}
  (deviation of \lname{JPoints.row3SetupInverses}: the array taken and returned, for the thread-local global; the globals that it reads passed; integers for the \texttt{long}s, with \texttt{Int.tmod} for C's \texttt{\%}; the entry of \texttt{inverses} read at the natural number of the index, which is the index for the \texttt{a ≥ 0} of the calls (for a negative \texttt{a} the C code reads before the row of the prime); natural numbers for the indices \texttt{m} and \texttt{n}.)
\item \csource{j-sift.c}{232}{248}: modeled by \lname{JPoints.row3SetupFactors}
  (deviation of \lname{JPoints.row3SetupFactors}: the arrays taken and returned, for the thread-local globals; the globals that it reads passed; integers for the \texttt{long}s, with \texttt{Int.tmod} for C's \texttt{\%}; \texttt{JPoints.Tab3Ptr} for the pointers of \texttt{tab3}, the indices of the rows of the tables of the curve for the addresses; natural numbers for the indices \texttt{m}, \texttt{n} and \texttt{pn}.)
\item \csource{j-sift.c}{249}{251}: modeled by \lname{JPoints.row3Setup}
  (deviation of \lname{JPoints.row3Setup}: the data taken and returned, for the thread-local globals; the globals that it reads passed; integers for the \texttt{long}s \texttt{a} and \texttt{b}.)
\item \csource{j-sift.c}{257}{270}: modeled by \lname{JPoints.checkBit}
  (deviation of \lname{JPoints.checkBit}: the data taken and returned, for the thread-local globals; the globals that it reads passed; an index into \texttt{res0} for the pointer \texttt{r0}, and the array of the chunk for the pointer \texttt{res}; integers for the \texttt{long}s \texttt{a}, \texttt{b}, \texttt{c} and \texttt{x}, natural numbers for the residues \texttt{r} and \texttt{j}, which the C code keeps nonnegative; \texttt{c \% 64}, the remainder in \texttt{[0, 64)}, for \texttt{c \& LONG\_MASK}, which agree in two's complement; a Boolean for the \texttt{int}.)
\item \csource{j-sift.c}{272}{272}: excluded (other): the forward declaration of relprime3; the model has no forward declarations
\item \csource{j-sift.c}{277}{293}: modeled by \lname{JPoints.checkBits}
  (deviation of \lname{JPoints.checkBits}: the state taken and returned, for the thread-local globals; the globals that it reads passed; an index into \texttt{res0} for the pointer \texttt{r0}, and the array of the chunk for the pointer \texttt{res}; integers for the \texttt{long}s \texttt{a}, \texttt{b}, \texttt{i} and \texttt{c}; \texttt{nums \&\&\& 1 != 0} for \texttt{nums \& 1}; a Boolean for the \texttt{int}.)
\item \csource{j-sift.c}{302}{308}: modeled by \lname{JPoints.relprime2}
  (deviation of \lname{JPoints.relprime2}: integers for the \texttt{long}s; \texttt{x \% 2 == 0} for \texttt{!(x \& 1)} and \texttt{x >>> 1} for \texttt{x >>= 1}, which agree on the positive \texttt{x} of the loops; the value of \texttt{n} after the loop, which \texttt{JPoints.oddGcd} returns, for the value of \texttt{m} that the C code compares with \texttt{1}, the same; a Boolean for the \texttt{int}.)
\item \csource{j-sift.c}{312}{317}: modeled by \lname{JPoints.oddGcd}
  (deviation of \lname{JPoints.oddGcd}: one function for the three loops of the C code, which are the same on the variables \texttt{(n, m)}, \texttt{(a, b)} and \texttt{(a, c)}, returning the value of the first; each loop \texttt{do \{ n >>= 1; \} while(!(n \& 1));} as one shift followed by a \texttt{while} loop of shifts; integers for the \texttt{long}s; \texttt{x \% 2 == 0} for \texttt{!(x \& 1)} and \texttt{x >>> 1} for \texttt{x >>= 1}, which agree on the positive \texttt{x} of the loop.)
\item \csource{j-sift.c}{318}{319}: modeled by \lname{JPoints.relprime2}
  (deviation of \lname{JPoints.relprime2}: integers for the \texttt{long}s; \texttt{x \% 2 == 0} for \texttt{!(x \& 1)} and \texttt{x >>> 1} for \texttt{x >>= 1}, which agree on the positive \texttt{x} of the loops; the value of \texttt{n} after the loop, which \texttt{JPoints.oddGcd} returns, for the value of \texttt{m} that the C code compares with \texttt{1}, the same; a Boolean for the \texttt{int}.)
\item \csource{j-sift.c}{322}{335}: modeled by \lname{JPoints.relprime3}
  (deviation of \lname{JPoints.relprime3}: integers for the \texttt{long}s, so that the negations of \texttt{b} and \texttt{c} are exact (the callers pass \texttt{|b| ≤ height} and \texttt{|c| < height + 64}); the test \texttt{((a | b | c) \& 1) == 0} as the conjunction of \texttt{x \% 2 == 0} for the three, which agree on every \texttt{long} in two's complement; \texttt{x \% 2 == 0} for \texttt{!(x \& 1)} and \texttt{x >>> 1} for \texttt{x >>= 1}, which agree on the positive \texttt{x} of the loops; a Boolean for the \texttt{int}.)
\item \csource{j-sift.c}{339}{344}: modeled by \lname{JPoints.oddGcd}
  (deviation of \lname{JPoints.oddGcd}: one function for the three loops of the C code, which are the same on the variables \texttt{(n, m)}, \texttt{(a, b)} and \texttt{(a, c)}, returning the value of the first; each loop \texttt{do \{ n >>= 1; \} while(!(n \& 1));} as one shift followed by a \texttt{while} loop of shifts; integers for the \texttt{long}s; \texttt{x \% 2 == 0} for \texttt{!(x \& 1)} and \texttt{x >>> 1} for \texttt{x >>= 1}, which agree on the positive \texttt{x} of the loop.)
\item \csource{j-sift.c}{346}{346}: modeled by \lname{JPoints.relprime3}
  (deviation of \lname{JPoints.relprime3}: integers for the \texttt{long}s, so that the negations of \texttt{b} and \texttt{c} are exact (the callers pass \texttt{|b| ≤ height} and \texttt{|c| < height + 64}); the test \texttt{((a | b | c) \& 1) == 0} as the conjunction of \texttt{x \% 2 == 0} for the three, which agree on every \texttt{long} in two's complement; \texttt{x \% 2 == 0} for \texttt{!(x \& 1)} and \texttt{x >>> 1} for \texttt{x >>= 1}, which agree on the positive \texttt{x} of the loops; a Boolean for the \texttt{int}.)
\item \csource{j-sift.c}{348}{353}: modeled by \lname{JPoints.oddGcd}
  (deviation of \lname{JPoints.oddGcd}: one function for the three loops of the C code, which are the same on the variables \texttt{(n, m)}, \texttt{(a, b)} and \texttt{(a, c)}, returning the value of the first; each loop \texttt{do \{ n >>= 1; \} while(!(n \& 1));} as one shift followed by a \texttt{while} loop of shifts; integers for the \texttt{long}s; \texttt{x \% 2 == 0} for \texttt{!(x \& 1)} and \texttt{x >>> 1} for \texttt{x >>= 1}, which agree on the positive \texttt{x} of the loop.)
\item \csource{j-sift.c}{354}{355}: modeled by \lname{JPoints.relprime3}
  (deviation of \lname{JPoints.relprime3}: integers for the \texttt{long}s, so that the negations of \texttt{b} and \texttt{c} are exact (the callers pass \texttt{|b| ≤ height} and \texttt{|c| < height + 64}); the test \texttt{((a | b | c) \& 1) == 0} as the conjunction of \texttt{x \% 2 == 0} for the three, which agree on every \texttt{long} in two's complement; \texttt{x \% 2 == 0} for \texttt{!(x \& 1)} and \texttt{x >>> 1} for \texttt{x >>= 1}, which agree on the positive \texttt{x} of the loops; a Boolean for the \texttt{int}.)
\item \csource{j-sift.c}{360}{373}: modeled by \lname{JPoints.checkPoint}
  (deviation of \lname{JPoints.checkPoint}: the state taken and returned, for the thread-local globals; the globals that it reads passed; an index into \texttt{res0} for the pointer \texttt{r0}, the index \texttt{rowptr[n]} into \texttt{sieve\_tab[n]} for the pointer, and the array of the chunk for the pointer \texttt{res}; integers for the \texttt{long}s \texttt{a}, \texttt{b} and \texttt{i}, natural numbers for \texttt{k}, \texttt{n} and the residue \texttt{r}, which the C code keeps nonnegative; the test of \texttt{nums} at the head of the loop as a \texttt{break}; a Boolean for the \texttt{int}.)
\item \csource{j-sift.c}{427}{427}: modeled by \lname{JPoints.tubeEps}
\item \csource{j-sift.c}{428}{428}: modeled by \lname{JPoints.tubeBig}
\item \csource{j-sift.c}{429}{429}: modeled by \lname{JPoints.tubeRows}
\item \csource{j-sift.c}{430}{430}: modeled by \lname{JPoints.coeffRow}
\item \csource{j-sift.c}{431}{431}: modeled by \lname{JPoints.numCoeffs}
\item \csource{j-sift.c}{432}{432}: modeled by \lname{JPoints.coeffIndex}
\item \csource{j-sift.c}{437}{439}: modeled by \lname{JPoints.TubeState}, \lname{JPoints.TubeState.initial}
  (deviation of \lname{JPoints.TubeState}: arrays of \texttt{JPoints.numCoeffs} doubles for the arrays of the coefficients; a Boolean for the \texttt{int} \texttt{tube\_ok}; a natural number for the \texttt{long} \texttt{tube\_ncells}, which starts at \texttt{0} and counts up; a structure, for the thread-local globals of the one thread.)
\item \csource{j-sift.c}{441}{441}: excluded (other): the forward declaration of tube\_eval; the model has no forward declarations
\item \csource{j-sift.c}{442}{446}: modeled by \lname{JPoints.tubeAInit}
  (deviation of \lname{JPoints.tubeAInit}: the state taken and returned, for the globals that it writes; the globals that it reads passed (the coefficients \texttt{kd}, the height bound, \texttt{th} and \texttt{th2}); \texttt{JPoints.Double.ofInt} for the conversions of the \texttt{long}s \texttt{a} and \texttt{height} and of the \texttt{unsigned long} \texttt{height + LONG\_LENGTH}, rounded to nearest, exact for the values of the program; \texttt{JPoints.Double.lt} for the comparison \texttt{<}, false when either side is a NaN.)
\item \csource{j-sift.c}{447}{453}: modeled by \lname{JPoints.tubeAInitK0}
  (deviation of \lname{JPoints.tubeAInitK0}: the array returned, for the global \texttt{K0} that the C code writes; the coefficients \texttt{kd} passed, for the globals that it reads, and \texttt{A}, \texttt{A2}, \texttt{A3} and \texttt{A4}, for the locals of \texttt{tube\_a\_init}.)
\item \csource{j-sift.c}{454}{457}: modeled by \lname{JPoints.tubeAInitK1K2}
  (deviation of \lname{JPoints.tubeAInitK1K2}: the arrays returned, for the globals \texttt{K1} and \texttt{K2} that the C code writes; the coefficients \texttt{kd} passed, for the globals that it reads, and \texttt{A}, \texttt{A2} and \texttt{A3}, for the locals of \texttt{tube\_a\_init}; the literals \texttt{-{}4.0} and \texttt{-{}2.0} as the negations of \texttt{4.0} and \texttt{2.0}, which they are.)
\item \csource{j-sift.c}{458}{467}: modeled by \lname{JPoints.tubeAInit}
  (deviation of \lname{JPoints.tubeAInit}: the state taken and returned, for the globals that it writes; the globals that it reads passed (the coefficients \texttt{kd}, the height bound, \texttt{th} and \texttt{th2}); \texttt{JPoints.Double.ofInt} for the conversions of the \texttt{long}s \texttt{a} and \texttt{height} and of the \texttt{unsigned long} \texttt{height + LONG\_LENGTH}, rounded to nearest, exact for the values of the program; \texttt{JPoints.Double.lt} for the comparison \texttt{<}, false when either side is a NaN.)
\item \csource{j-sift.c}{481}{481}: modeled by \lname{JPoints.rowRanges}
\item \csource{j-sift.c}{482}{483}: modeled by \lname{JPoints.rowBeyond}
\item \csource{j-sift.c}{484}{490}: modeled by \lname{JPoints.rowRegion}
  (deviation of \lname{JPoints.rowRegion}: the ranges returned as arrays of their number, for the arrays \texttt{lo} and \texttt{hi} of the caller, which have room for \texttt{ROW\_RANGES} (\texttt{JPoints.rowRanges}); the intervals and the static data passed, for the globals and the statics that it reads; the arrays of the excluded intervals returned by the functions that fill them, for the local arrays \texttt{xlo} and \texttt{xhi}, with their size for the count \texttt{n}; integers for the \texttt{long}s \texttt{a} and \texttt{b} and for the ends of the ranges, natural numbers for the counts.)
\item \csource{j-sift.c}{491}{508}: modeled by \lname{JPoints.rowRegionPositive}
  (deviation of \lname{JPoints.rowRegionPositive}: the ends returned as two arrays of the number \texttt{n} of the intervals, filled in the order in which the C code writes them, for the local arrays \texttt{xlo} and \texttt{xhi} of eight doubles and the count; the intervals passed, for the globals; \texttt{JPoints.Double.ofInt} for the conversions of the \texttt{long}s \texttt{a} and \texttt{b}, exact for the values of the program; \texttt{JPoints.Double.lt} for \texttt{>} and \texttt{<} and \texttt{JPoints.Double.beq} for \texttt{==}, false when either side is a NaN.)
\item \csource{j-sift.c}{509}{510}: modeled by \lname{JPoints.rowRegion}
  (deviation of \lname{JPoints.rowRegion}: the ranges returned as arrays of their number, for the arrays \texttt{lo} and \texttt{hi} of the caller, which have room for \texttt{ROW\_RANGES} (\texttt{JPoints.rowRanges}); the intervals and the static data passed, for the globals and the statics that it reads; the arrays of the excluded intervals returned by the functions that fill them, for the local arrays \texttt{xlo} and \texttt{xhi}, with their size for the count \texttt{n}; integers for the \texttt{long}s \texttt{a} and \texttt{b} and for the ends of the ranges, natural numbers for the counts.)
\item \csource{j-sift.c}{511}{517}: modeled by \lname{JPoints.rowRegionZero}
  (deviation of \lname{JPoints.rowRegionZero}: the ends returned as two arrays of the number \texttt{n} of the intervals, filled in the order in which the C code writes them, for the local arrays \texttt{xlo} and \texttt{xhi} of eight doubles and the count; the intervals passed, for the globals; \texttt{JPoints.Double.ofInt} for the conversion of the \texttt{long} \texttt{b}, exact for the values of the program; \texttt{JPoints.Double.beq} for \texttt{==}.)
\item \csource{j-sift.c}{518}{518}: modeled by \lname{JPoints.rowRegion}
  (deviation of \lname{JPoints.rowRegion}: the ranges returned as arrays of their number, for the arrays \texttt{lo} and \texttt{hi} of the caller, which have room for \texttt{ROW\_RANGES} (\texttt{JPoints.rowRanges}); the intervals and the static data passed, for the globals and the statics that it reads; the arrays of the excluded intervals returned by the functions that fill them, for the local arrays \texttt{xlo} and \texttt{xhi}, with their size for the count \texttt{n}; integers for the \texttt{long}s \texttt{a} and \texttt{b} and for the ends of the ranges, natural numbers for the counts.)
\item \csource{j-sift.c}{521}{524}: modeled by \lname{JPoints.rowRegionMargins}
  (deviation of \lname{JPoints.rowRegionMargins}: the arrays taken and returned, for the local arrays that the C code writes; \texttt{JPoints.Double.lt} for \texttt{>} and \texttt{<}, false when either side is a NaN.)
\item \csource{j-sift.c}{525}{529}: modeled by \lname{JPoints.rowRegionSort}
  (deviation of \lname{JPoints.rowRegionSort}: the arrays taken and returned, for the local arrays that the C code writes; natural numbers for the indices \texttt{i} and \texttt{j}, which the C code keeps nonnegative; the inner loop, which ends on the data, a \texttt{while} loop; \texttt{JPoints.Double.lt} for \texttt{>}, false when either side is a NaN.)
\item \csource{j-sift.c}{531}{536}: modeled by \lname{JPoints.rowRegionWords}
  (deviation of \lname{JPoints.rowRegionWords}: the ranges returned as arrays of their number, filled in the order in which the C code writes them, for the arrays \texttt{lo} and \texttt{hi} of the caller, which have room for \texttt{ROW\_RANGES} ranges (\texttt{JPoints.rowRanges}), and there are at most \texttt{2 * 4 + 1} of them; the static data passed, for the statics \texttt{w\_low\_all} and \texttt{w\_high\_all}; integers for the \texttt{long}s, a natural number for the count \texttt{nr} and the index \texttt{i}; \texttt{JPoints.Double.le} for \texttt{<=} and \texttt{>=} and \texttt{JPoints.Double.lt} for \texttt{<}, false when either side is a NaN; \texttt{JPoints.Double.ofInt} for the conversion of \texttt{LONG\_LENGTH}.)
\item \csource{j-sift.c}{543}{547}: modeled by \lname{JPoints.rowRegionEnds}
  (deviation of \lname{JPoints.rowRegionEnds}: the words returned as a pair, for the local variables \texttt{wl} and \texttt{wh}; the static data passed, for the statics \texttt{w\_low\_all} and \texttt{w\_high\_all}; integers for the \texttt{long}s, which they are, since an end that is truncated lies in \texttt{(-{}ROW\_BEYOND, ROW\_BEYOND)} (a NaN, which the ends of the program never are, would truncate to \texttt{-{}2 \textasciicircum{} 63}, from which \texttt{(long)xh -{} 2} overflows); \texttt{JPoints.Double.toInt64} for the conversion \texttt{(long)}, which truncates towards zero; \texttt{>>>} on integers for C's \texttt{>>} on \texttt{long}s, which gcc shifts arithmetically, also for a negative \texttt{c}; \texttt{JPoints.Double.le} for \texttt{<=} and \texttt{>=}, false when either side is a NaN.)
\item \csource{j-sift.c}{548}{554}: modeled by \lname{JPoints.rowRegionWords}
  (deviation of \lname{JPoints.rowRegionWords}: the ranges returned as arrays of their number, filled in the order in which the C code writes them, for the arrays \texttt{lo} and \texttt{hi} of the caller, which have room for \texttt{ROW\_RANGES} ranges (\texttt{JPoints.rowRanges}), and there are at most \texttt{2 * 4 + 1} of them; the static data passed, for the statics \texttt{w\_low\_all} and \texttt{w\_high\_all}; integers for the \texttt{long}s, a natural number for the count \texttt{nr} and the index \texttt{i}; \texttt{JPoints.Double.le} for \texttt{<=} and \texttt{>=} and \texttt{JPoints.Double.lt} for \texttt{<}, false when either side is a NaN; \texttt{JPoints.Double.ofInt} for the conversion of \texttt{LONG\_LENGTH}.)
\item \csource{j-sift.c}{555}{556}: modeled by \lname{JPoints.rowRegion}
  (deviation of \lname{JPoints.rowRegion}: the ranges returned as arrays of their number, for the arrays \texttt{lo} and \texttt{hi} of the caller, which have room for \texttt{ROW\_RANGES} (\texttt{JPoints.rowRanges}); the intervals and the static data passed, for the globals and the statics that it reads; the arrays of the excluded intervals returned by the functions that fill them, for the local arrays \texttt{xlo} and \texttt{xhi}, with their size for the count \texttt{n}; integers for the \texttt{long}s \texttt{a} and \texttt{b} and for the ends of the ranges, natural numbers for the counts.)
\item \csource{j-sift.c}{560}{566}: modeled by \lname{JPoints.shift1}
  (deviation of \lname{JPoints.shift1}: the array and the index \texttt{base} in it, for the pointer \texttt{p} into the array of the caller (\texttt{Q + k} or \texttt{Q + L j}), and the array returned, for the array that the C code writes; natural numbers for the \texttt{long}s; the loop of \texttt{j} down from \texttt{deg -{} 1} to \texttt{i} as a loop up over \texttt{t < deg -{} i} with \texttt{j = deg -{} 1 -{} t}.)
\item \csource{j-sift.c}{569}{575}: modeled by \lname{JPoints.shift2}
  (deviation of \lname{JPoints.shift2}: the array returned, for the array that the C code shifts in place; a natural number for the \texttt{long} \texttt{deg}.)
\item \csource{j-sift.c}{579}{587}: modeled by \lname{JPoints.tubeSum}
  (deviation of \lname{JPoints.tubeSum}: the arrays passed, for the pointers to them; a natural number for the \texttt{long} \texttt{deg}.)
\item \csource{j-sift.c}{590}{600}: modeled by \lname{JPoints.tubeEval}
  (deviation of \lname{JPoints.tubeEval}: the array passed, for the pointer to it; a natural number for the \texttt{long} \texttt{deg}.)
\item \csource{j-sift.c}{602}{609}: modeled by \lname{JPoints.polySign}
  (deviation of \lname{JPoints.polySign}: an integer for the \texttt{int} sign; \texttt{JPoints.Double.lt} for the comparisons \texttt{>}, false when either side is a NaN.)
\item \csource{j-sift.c}{614}{626}: modeled by \lname{JPoints.tubeEmit}
  (deviation of \lname{JPoints.tubeEmit}: the bands passed and returned, for the thread-local globals that the C code reads and writes; integers for the \texttt{long}s \texttt{b}, \texttt{i\_lo} and \texttt{i\_hi} and for the ends of the ranges, a natural number for the index \texttt{row} of the row in the band, which the C code keeps nonnegative; the test \texttt{nr > 0 \&\& …} as a conjunction, the second part read only when the first holds.)
\item \csource{j-sift.c}{629}{630}: modeled by \lname{JPoints.tubeCell}
  (deviation of \lname{JPoints.tubeCell}: the state taken and returned, for the globals and the statics that it writes; the bound and its square passed, for the statics that it reads; integers for the \texttt{long}s, and Lean's \texttt{/} for C's \texttt{/} on the positive \texttt{words} and \texttt{rows -{} 1}, where it truncates the same way; the recursion well founded on the number of the rows and the words of the cell, which each recursive call lowers (\texttt{JPoints.tubeRows\_pos}).)
\item \csource{j-sift.c}{631}{631}: modeled by \lname{JPoints.TubeCellCache}, \lname{JPoints.TubeCellCache.initial}
  (deviation of \lname{JPoints.TubeCellCache}: an array for the array \texttt{rp}; a structure, for the thread-local statics of the one thread.)
\item \csource{j-sift.c}{632}{632}: modeled by \lname{JPoints.tubeCell}
  (deviation of \lname{JPoints.tubeCell}: the state taken and returned, for the globals and the statics that it writes; the bound and its square passed, for the statics that it reads; integers for the \texttt{long}s, and Lean's \texttt{/} for C's \texttt{/} on the positive \texttt{words} and \texttt{rows -{} 1}, where it truncates the same way; the recursion well founded on the number of the rows and the words of the cell, which each recursive call lowers (\texttt{JPoints.tubeRows\_pos}).)
\item \csource{j-sift.c}{633}{633}: modeled by \lname{JPoints.tubeCellAnalyze}
  (deviation of \lname{JPoints.tubeCellAnalyze}: the bound \texttt{th} and its square \texttt{th2} passed, for the statics, and the polynomials of the state and the products \texttt{rp}, for the globals and the static that it reads; the shifted copies of \texttt{K0}, \texttt{K1} and \texttt{K2} returned by \texttt{JPoints.shift2}, for the copies into the local arrays \texttt{p0}, \texttt{p1} and \texttt{p2} that it shifts in place; \texttt{inside} and the returns of the C code returned as a pair of Booleans; integers for the \texttt{int} signs; \texttt{JPoints.Double.lt} for the comparison \texttt{>}, false when either side is a NaN, which keeps the cell.)
\item \csource{j-sift.c}{634}{636}: modeled by \lname{JPoints.tubeCellDecide}
  (deviation of \lname{JPoints.tubeCellDecide}: the bound \texttt{th} and its square \texttt{th2} passed, for the statics; the statics taken and returned in the state, for the thread-local statics that the C code writes; the center and the half-sizes computed after the test of \texttt{tube\_ok}, which reads none of them; \texttt{JPoints.Double.ofInt} for the exact conversions of the \texttt{long}s to doubles.)
\item \csource{j-sift.c}{637}{640}: modeled by \lname{JPoints.tubeCellAnalyze}
  (deviation of \lname{JPoints.tubeCellAnalyze}: the bound \texttt{th} and its square \texttt{th2} passed, for the statics, and the polynomials of the state and the products \texttt{rp}, for the globals and the static that it reads; the shifted copies of \texttt{K0}, \texttt{K1} and \texttt{K2} returned by \texttt{JPoints.shift2}, for the copies into the local arrays \texttt{p0}, \texttt{p1} and \texttt{p2} that it shifts in place; \texttt{inside} and the returns of the C code returned as a pair of Booleans; integers for the \texttt{int} signs; \texttt{JPoints.Double.lt} for the comparison \texttt{>}, false when either side is a NaN, which keeps the cell.)
\item \csource{j-sift.c}{641}{643}: modeled by \lname{JPoints.tubeCell}
  (deviation of \lname{JPoints.tubeCell}: the state taken and returned, for the globals and the statics that it writes; the bound and its square passed, for the statics that it reads; integers for the \texttt{long}s, and Lean's \texttt{/} for C's \texttt{/} on the positive \texttt{words} and \texttt{rows -{} 1}, where it truncates the same way; the recursion well founded on the number of the rows and the words of the cell, which each recursive call lowers (\texttt{JPoints.tubeRows\_pos}).)
\item \csource{j-sift.c}{644}{644}: modeled by \lname{JPoints.tubeEmitRows}
  (deviation of \lname{JPoints.tubeEmitRows}: one function for the three loops, which differ only in the words; the loop \texttt{for(b = b\_lo; b <= b\_hi; b++)} over \texttt{b = b\_lo + k} for \texttt{k < b\_hi -{} b\_lo + 1}, none when \texttt{b\_hi < b\_lo}; integers for the \texttt{long}s \texttt{b\_lo}, \texttt{b\_hi}, \texttt{i\_lo} and \texttt{i\_hi}; the bands taken out of the state while they are written, so that they are not shared, and the state returned.)
\item \csource{j-sift.c}{645}{647}: modeled by \lname{JPoints.tubeCell}
  (deviation of \lname{JPoints.tubeCell}: the state taken and returned, for the globals and the statics that it writes; the bound and its square passed, for the statics that it reads; integers for the \texttt{long}s, and Lean's \texttt{/} for C's \texttt{/} on the positive \texttt{words} and \texttt{rows -{} 1}, where it truncates the same way; the recursion well founded on the number of the rows and the words of the cell, which each recursive call lowers (\texttt{JPoints.tubeRows\_pos}).)
\item \csource{j-sift.c}{648}{655}: modeled by \lname{JPoints.tubeCellProducts}
  (deviation of \lname{JPoints.tubeCellProducts}: the statics taken and returned, for the thread-local statics that the C code writes; \texttt{!rb.beq last\_rb} for the comparison \texttt{rb != last\_rb}.)
\item \csource{j-sift.c}{656}{666}: modeled by \lname{JPoints.tubeCellAnalyze}
  (deviation of \lname{JPoints.tubeCellAnalyze}: the bound \texttt{th} and its square \texttt{th2} passed, for the statics, and the polynomials of the state and the products \texttt{rp}, for the globals and the static that it reads; the shifted copies of \texttt{K0}, \texttt{K1} and \texttt{K2} returned by \texttt{JPoints.shift2}, for the copies into the local arrays \texttt{p0}, \texttt{p1} and \texttt{p2} that it shifts in place; \texttt{inside} and the returns of the C code returned as a pair of Booleans; integers for the \texttt{int} signs; \texttt{JPoints.Double.lt} for the comparison \texttt{>}, false when either side is a NaN, which keeps the cell.)
\item \csource{j-sift.c}{667}{670}: modeled by \lname{JPoints.tubeCellCombine}
  (deviation of \lname{JPoints.tubeCellCombine}: the bound \texttt{th} and its square \texttt{th2} passed, for the statics, and the shifted coefficients, for the local arrays \texttt{p0}, \texttt{p1} and \texttt{p2}; fresh arrays filled by the loop, returned, for the local arrays of \texttt{tube\_cell}.)
\item \csource{j-sift.c}{671}{683}: modeled by \lname{JPoints.tubeCellAnalyze}
  (deviation of \lname{JPoints.tubeCellAnalyze}: the bound \texttt{th} and its square \texttt{th2} passed, for the statics, and the polynomials of the state and the products \texttt{rp}, for the globals and the static that it reads; the shifted copies of \texttt{K0}, \texttt{K1} and \texttt{K2} returned by \texttt{JPoints.shift2}, for the copies into the local arrays \texttt{p0}, \texttt{p1} and \texttt{p2} that it shifts in place; \texttt{inside} and the returns of the C code returned as a pair of Booleans; integers for the \texttt{int} signs; \texttt{JPoints.Double.lt} for the comparison \texttt{>}, false when either side is a NaN, which keeps the cell.)
\item \csource{j-sift.c}{684}{686}: modeled by \lname{JPoints.tubeCell}
  (deviation of \lname{JPoints.tubeCell}: the state taken and returned, for the globals and the statics that it writes; the bound and its square passed, for the statics that it reads; integers for the \texttt{long}s, and Lean's \texttt{/} for C's \texttt{/} on the positive \texttt{words} and \texttt{rows -{} 1}, where it truncates the same way; the recursion well founded on the number of the rows and the words of the cell, which each recursive call lowers (\texttt{JPoints.tubeRows\_pos}).)
\item \csource{j-sift.c}{687}{687}: modeled by \lname{JPoints.tubeEmitRows}
  (deviation of \lname{JPoints.tubeEmitRows}: one function for the three loops, which differ only in the words; the loop \texttt{for(b = b\_lo; b <= b\_hi; b++)} over \texttt{b = b\_lo + k} for \texttt{k < b\_hi -{} b\_lo + 1}, none when \texttt{b\_hi < b\_lo}; integers for the \texttt{long}s \texttt{b\_lo}, \texttt{b\_hi}, \texttt{i\_lo} and \texttt{i\_hi}; the bands taken out of the state while they are written, so that they are not shared, and the state returned.)
\item \csource{j-sift.c}{688}{705}: modeled by \lname{JPoints.tubeCell}
  (deviation of \lname{JPoints.tubeCell}: the state taken and returned, for the globals and the statics that it writes; the bound and its square passed, for the statics that it reads; integers for the \texttt{long}s, and Lean's \texttt{/} for C's \texttt{/} on the positive \texttt{words} and \texttt{rows -{} 1}, where it truncates the same way; the recursion well founded on the number of the rows and the words of the cell, which each recursive call lowers (\texttt{JPoints.tubeRows\_pos}).)
\item \csource{j-sift.c}{706}{706}: modeled by \lname{JPoints.tubeEmitRows}
  (deviation of \lname{JPoints.tubeEmitRows}: one function for the three loops, which differ only in the words; the loop \texttt{for(b = b\_lo; b <= b\_hi; b++)} over \texttt{b = b\_lo + k} for \texttt{k < b\_hi -{} b\_lo + 1}, none when \texttt{b\_hi < b\_lo}; integers for the \texttt{long}s \texttt{b\_lo}, \texttt{b\_hi}, \texttt{i\_lo} and \texttt{i\_hi}; the bands taken out of the state while they are written, so that they are not shared, and the state returned.)
\item \csource{j-sift.c}{707}{709}: modeled by \lname{JPoints.tubeCell}
  (deviation of \lname{JPoints.tubeCell}: the state taken and returned, for the globals and the statics that it writes; the bound and its square passed, for the statics that it reads; integers for the \texttt{long}s, and Lean's \texttt{/} for C's \texttt{/} on the positive \texttt{words} and \texttt{rows -{} 1}, where it truncates the same way; the recursion well founded on the number of the rows and the words of the cell, which each recursive call lowers (\texttt{JPoints.tubeRows\_pos}).)
\item \csource{j-sift.c}{714}{724}: modeled by \lname{JPoints.rowCut}
  (deviation of \lname{JPoints.rowCut}: the ranges returned as arrays of their number, filled in the order in which the C code writes them, for the arrays \texttt{rlo} and \texttt{rhi} of the caller; the bands and the arrays of the region passed, for the globals and the pointers; integers for the \texttt{long}s of the ends of the ranges, natural numbers for the counts and the indices, which the C code keeps nonnegative.)
\item \csource{j-sift.c}{731}{737}: modeled by \lname{JPoints.tubeSample}
  (deviation of \lname{JPoints.tubeSample}: the static data and the state of the analysis taken and returned, for the statics and the thread-local globals that it writes; the globals that it reads passed; the statistics returned, for the globals that the C code writes through the pointers of its arguments; the sums in a structure, for the locals (\texttt{JPoints.TubeSampleSums}), \texttt{Float}s for the doubles, an integer for the \texttt{long} \texttt{rows} and natural numbers for the \texttt{long}s \texttt{left} and \texttt{tleft}, which count up from \texttt{0}; \texttt{>>>} on integers for C's \texttt{>>} on the \texttt{long}s, which gcc shifts arithmetically also for the negative \texttt{-{}height}; \texttt{JPoints.Double.ofInt} for the conversion of the \texttt{long} \texttt{dbound}, rounded to nearest; the loop \texttt{for(b = b\_lo; b <= b\_hi; b++)} over \texttt{b = b\_lo + j} for \texttt{j < b\_hi -{} b\_lo + 1}; \texttt{Float.ofNat} for the conversion of the \texttt{long} \texttt{tube\_ncells}, exact.)
\item \csource{j-sift.c}{738}{738}: excluded (memory): the allocation of the state of the calling thread (init\_thread\_sieve) when tube\_sample runs before init\_sift has made it; the model makes the lists of the bands as fresh arrays
\item \csource{j-sift.c}{739}{739}: modeled by \lname{JPoints.tubeSample}
  (deviation of \lname{JPoints.tubeSample}: the static data and the state of the analysis taken and returned, for the statics and the thread-local globals that it writes; the globals that it reads passed; the statistics returned, for the globals that the C code writes through the pointers of its arguments; the sums in a structure, for the locals (\texttt{JPoints.TubeSampleSums}), \texttt{Float}s for the doubles, an integer for the \texttt{long} \texttt{rows} and natural numbers for the \texttt{long}s \texttt{left} and \texttt{tleft}, which count up from \texttt{0}; \texttt{>>>} on integers for C's \texttt{>>} on the \texttt{long}s, which gcc shifts arithmetically also for the negative \texttt{-{}height}; \texttt{JPoints.Double.ofInt} for the conversion of the \texttt{long} \texttt{dbound}, rounded to nearest; the loop \texttt{for(b = b\_lo; b <= b\_hi; b++)} over \texttt{b = b\_lo + j} for \texttt{j < b\_hi -{} b\_lo + 1}; \texttt{Float.ofNat} for the conversion of the \texttt{long} \texttt{tube\_ncells}, exact.)
\item \csource{j-sift.c}{740}{746}: modeled by \lname{JPoints.tubeSampleRange}
  (deviation of \lname{JPoints.tubeSampleRange}: the curve and the height bound passed, for the globals that it reads; the first coordinate and the rows returned, for the locals of \texttt{tube\_sample}; integers for the \texttt{long}s, with \texttt{Int.tmod} for C's \texttt{\%} on the nonnegative operands; \texttt{JPoints.Double.ofInt} for the conversion of the \texttt{long} \texttt{height}, exact, \texttt{JPoints.Double.sqrt}, \texttt{JPoints.Double.floor} and the sum with \texttt{1.0} for \texttt{floor(sqrt((double)height) + 1)}, and \texttt{JPoints.Double.toInt64} for the cast to a \texttt{long}, which truncates.)
\item \csource{j-sift.c}{747}{751}: modeled by \lname{JPoints.tubeSampleBand}
  (deviation of \lname{JPoints.tubeSampleBand}: the state of the analysis taken and returned, for the thread-local globals of the bands and the globals of the analysis; the globals and the statics that it reads passed; integers for the \texttt{long}s \texttt{a}, \texttt{b\_lo}, \texttt{b\_hi}, \texttt{w\_low} and \texttt{w\_high}; the loop \texttt{for(b = b\_lo; b <= b\_hi; b++)} over \texttt{b = b\_lo + k} for \texttt{k < b\_hi -{} b\_lo + 1}; the bands taken out of the state while their counts are written, so that they are not shared.)
\item \csource{j-sift.c}{752}{752}: modeled by \lname{JPoints.tubeSample}
  (deviation of \lname{JPoints.tubeSample}: the static data and the state of the analysis taken and returned, for the statics and the thread-local globals that it writes; the globals that it reads passed; the statistics returned, for the globals that the C code writes through the pointers of its arguments; the sums in a structure, for the locals (\texttt{JPoints.TubeSampleSums}), \texttt{Float}s for the doubles, an integer for the \texttt{long} \texttt{rows} and natural numbers for the \texttt{long}s \texttt{left} and \texttt{tleft}, which count up from \texttt{0}; \texttt{>>>} on integers for C's \texttt{>>} on the \texttt{long}s, which gcc shifts arithmetically also for the negative \texttt{-{}height}; \texttt{JPoints.Double.ofInt} for the conversion of the \texttt{long} \texttt{dbound}, rounded to nearest; the loop \texttt{for(b = b\_lo; b <= b\_hi; b++)} over \texttt{b = b\_lo + j} for \texttt{j < b\_hi -{} b\_lo + 1}; \texttt{Float.ofNat} for the conversion of the \texttt{long} \texttt{tube\_ncells}, exact.)
\item \csource{j-sift.c}{753}{761}: modeled by \lname{JPoints.tubeSampleRow}
  (deviation of \lname{JPoints.tubeSampleRow}: the sums taken and returned, for the locals of the caller; the ranges of the real region and of the cut returned as arrays of their number, for the local arrays \texttt{reglo}, \texttt{reghi}, \texttt{rlo} and \texttt{rhi}; the intervals, the static data and the bands passed, for the globals and the statics that it reads; integers for the \texttt{long}s \texttt{a}, \texttt{b} and \texttt{b\_lo} and for the ends of the ranges; a natural number for the index \texttt{b -{} b\_lo} of the row in the band and for the counts; \texttt{Float.ofInt} and \texttt{Float.ofNat} for the conversions of the \texttt{long}s to doubles, exact.)
\item \csource{j-sift.c}{762}{765}: modeled by \lname{JPoints.tubeSample}
  (deviation of \lname{JPoints.tubeSample}: the static data and the state of the analysis taken and returned, for the statics and the thread-local globals that it writes; the globals that it reads passed; the statistics returned, for the globals that the C code writes through the pointers of its arguments; the sums in a structure, for the locals (\texttt{JPoints.TubeSampleSums}), \texttt{Float}s for the doubles, an integer for the \texttt{long} \texttt{rows} and natural numbers for the \texttt{long}s \texttt{left} and \texttt{tleft}, which count up from \texttt{0}; \texttt{>>>} on integers for C's \texttt{>>} on the \texttt{long}s, which gcc shifts arithmetically also for the negative \texttt{-{}height}; \texttt{JPoints.Double.ofInt} for the conversion of the \texttt{long} \texttt{dbound}, rounded to nearest; the loop \texttt{for(b = b\_lo; b <= b\_hi; b++)} over \texttt{b = b\_lo + j} for \texttt{j < b\_hi -{} b\_lo + 1}; \texttt{Float.ofNat} for the conversion of the \texttt{long} \texttt{tube\_ncells}, exact.)
\item \csource{j-sift.c}{766}{772}: modeled by \lname{JPoints.tubeSampleStats}
  (deviation of \lname{JPoints.tubeSampleStats}: the statistics returned, for the globals that the C code writes through the pointers of its arguments; \texttt{Float.ofInt} and \texttt{Float.ofNat} for the conversions of the \texttt{long}s, exact.)
\item \csource{j-sift.c}{773}{774}: modeled by \lname{JPoints.tubeSample}
  (deviation of \lname{JPoints.tubeSample}: the static data and the state of the analysis taken and returned, for the statics and the thread-local globals that it writes; the globals that it reads passed; the statistics returned, for the globals that the C code writes through the pointers of its arguments; the sums in a structure, for the locals (\texttt{JPoints.TubeSampleSums}), \texttt{Float}s for the doubles, an integer for the \texttt{long} \texttt{rows} and natural numbers for the \texttt{long}s \texttt{left} and \texttt{tleft}, which count up from \texttt{0}; \texttt{>>>} on integers for C's \texttt{>>} on the \texttt{long}s, which gcc shifts arithmetically also for the negative \texttt{-{}height}; \texttt{JPoints.Double.ofInt} for the conversion of the \texttt{long} \texttt{dbound}, rounded to nearest; the loop \texttt{for(b = b\_lo; b <= b\_hi; b++)} over \texttt{b = b\_lo + j} for \texttt{j < b\_hi -{} b\_lo + 1}; \texttt{Float.ofNat} for the conversion of the \texttt{long} \texttt{tube\_ncells}, exact.)
\item \csource{j-sift.c}{780}{780}: modeled by \lname{JPoints.SiftState}, \lname{JPoints.SiftState.initial}
  (deviation of \lname{JPoints.SiftState}: an array for the bit array that \texttt{survivors} points to; indices into the sieve tables for the pointers \texttt{rowptr[n]} into \texttt{sieve\_tab[n]}; natural numbers for the counters, which start at \texttt{0} and count up, and for the index \texttt{cur\_chunk}; a structure, for the thread-local globals of the one thread.)
  (deviation of \lname{JPoints.SiftState.initial}: a fresh array of \texttt{arraySize + VW} zeros, \texttt{VW = 1}, for the bit array that \texttt{init\_thread\_sieve} allocates (memory management, not modeled), every word of which the sieve writes before it reads it; index \texttt{0} for the null pointers of \texttt{rowptr}, which the sieve reads only after \texttt{row\_setup} has set them.)
\item \csource{j-sift.c}{781}{786}: modeled by \lname{JPoints.tubeWord}
  (deviation of \lname{JPoints.tubeWord}: the state taken and returned, for the thread-local globals; the globals and the static data that it reads passed; the chunk read from the array of the chunks, for the pointer \texttt{ch}; the index \texttt{k * npr} into \texttt{res0} for the pointer \texttt{r0}, the index \texttt{rowptr[n]} into \texttt{sieve\_tab[n]} for the pointer, and the array of the chunk for the pointer \texttt{res}; integers for the \texttt{long}s \texttt{a}, \texttt{b}, \texttt{i} and \texttt{c0}, natural numbers for \texttt{k}, \texttt{n} and the residue \texttt{r}, which the C code keeps nonnegative; \texttt{c0 >>> 6} for \texttt{c0 >> LONG\_SHIFT} and \texttt{c0 \% 64}, in \texttt{[0, 64)}, for \texttt{c0 \& LONG\_MASK}, the same for the nonnegative \texttt{c0} of the calls; a Boolean for the \texttt{int}.)
\item \csource{j-sift.c}{787}{788}: modeled by \lname{JPoints.seekChunk}
  (deviation of \lname{JPoints.seekChunk}: one function for the two pairs of loops, which are the same on \texttt{cur\_chunk} and the word, \texttt{i} in \texttt{tube\_word} and \texttt{lo} in \texttt{sift\_bands}, returning \texttt{cur\_chunk}; the static data passed, for the static \texttt{chunks}; a natural number for the index \texttt{cur\_chunk}, which stays in \texttt{[0, num\_chunks)} for a word of the row (for a word outside the row the C code reads beyond the array of the chunks, and the loops of the model need not end); integers for the \texttt{long}s of the words.)
\item \csource{j-sift.c}{789}{805}: modeled by \lname{JPoints.tubeWord}
  (deviation of \lname{JPoints.tubeWord}: the state taken and returned, for the thread-local globals; the globals and the static data that it reads passed; the chunk read from the array of the chunks, for the pointer \texttt{ch}; the index \texttt{k * npr} into \texttt{res0} for the pointer \texttt{r0}, the index \texttt{rowptr[n]} into \texttt{sieve\_tab[n]} for the pointer, and the array of the chunk for the pointer \texttt{res}; integers for the \texttt{long}s \texttt{a}, \texttt{b}, \texttt{i} and \texttt{c0}, natural numbers for \texttt{k}, \texttt{n} and the residue \texttt{r}, which the C code keeps nonnegative; \texttt{c0 >>> 6} for \texttt{c0 >> LONG\_SHIFT} and \texttt{c0 \% 64}, in \texttt{[0, 64)}, for \texttt{c0 \& LONG\_MASK}, the same for the nonnegative \texttt{c0} of the calls; a Boolean for the \texttt{int}.)
\item \csource{j-sift.c}{807}{808}: excluded (other): the forward declarations of row\_setup and sift\_range; the model has no forward declarations
\item \csource{j-sift.c}{815}{826}: modeled by \lname{JPoints.siftBands}
  (deviation of \lname{JPoints.siftBands}: the state taken and returned, for the thread-local globals; the globals and the static data that it reads passed; integers for the \texttt{long}s, with \texttt{Int.tmod} for C's \texttt{\%}; the loop \texttt{for(b\_lo = b\_first; b\_lo <= height; b\_lo += step)} over \texttt{b\_lo = b\_first + t step} for \texttt{t < (height -{} b\_first) / step + 1}, with Lean's \texttt{/}, which rounds down, none when \texttt{b\_first > height}, and the loop \texttt{for(b = b\_lo; b <= b\_hi; b += m)} over \texttt{b = b\_lo + k m} for \texttt{k < (b\_hi -{} b\_lo) / m + 1}; Booleans for the \texttt{int} \texttt{tube} and for the \texttt{int} returned.)
\item \csource{j-sift.c}{827}{831}: modeled by \lname{JPoints.siftBandsBand}
  (deviation of \lname{JPoints.siftBandsBand}: the state of the analysis taken and returned, for the thread-local globals of the bands and the globals of the analysis; the static data passed, for the statics; integers for the \texttt{long}s \texttt{b\_lo} and \texttt{b\_hi}; the loop \texttt{for(k = 0; k <= b\_hi -{} b\_lo; k++)} over \texttt{k < b\_hi -{} b\_lo + 1}; the bands taken out of the state while their counts are written, so that they are not shared.)
\item \csource{j-sift.c}{832}{832}: modeled by \lname{JPoints.siftBands}
  (deviation of \lname{JPoints.siftBands}: the state taken and returned, for the thread-local globals; the globals and the static data that it reads passed; integers for the \texttt{long}s, with \texttt{Int.tmod} for C's \texttt{\%}; the loop \texttt{for(b\_lo = b\_first; b\_lo <= height; b\_lo += step)} over \texttt{b\_lo = b\_first + t step} for \texttt{t < (height -{} b\_first) / step + 1}, with Lean's \texttt{/}, which rounds down, none when \texttt{b\_first > height}, and the loop \texttt{for(b = b\_lo; b <= b\_hi; b += m)} over \texttt{b = b\_lo + k m} for \texttt{k < (b\_hi -{} b\_lo) / m + 1}; Booleans for the \texttt{int} \texttt{tube} and for the \texttt{int} returned.)
\item \csource{j-sift.c}{833}{837}: modeled by \lname{JPoints.siftBandsRow}
  (deviation of \lname{JPoints.siftBandsRow}: the state taken and returned, for the thread-local globals; the globals and the static data that it reads passed; integers for the \texttt{long}s, a natural number for the index \texttt{b -{} b\_lo} of the row in the band; \texttt{a \% 64} and \texttt{b \% 64}, in \texttt{[0, 64)}, for \texttt{a \& 63} and \texttt{b \& 63}, the same in two's complement, and the entry \texttt{(a \& 63) * 64 + (b \& 63)} of the flat array \texttt{alive2}; the \texttt{continue}s of the C code as returns of \texttt{false}, and whether \texttt{sift\_bands} returns \texttt{1} from the row, for its return.)
\item \csource{j-sift.c}{838}{849}: modeled by \lname{JPoints.siftBandsRanges}
  (deviation of \lname{JPoints.siftBandsRanges}: the ranges returned as arrays of their number, for the arrays \texttt{reglo}, \texttt{reghi}, \texttt{clo} and \texttt{chi} and the pointers \texttt{rlo} and \texttt{rhi} to the arrays sieved; the bands and the globals that it reads passed; integers for the \texttt{long}s \texttt{a}, \texttt{b} and \texttt{b\_lo} and for the ends of the ranges, a natural number for the index \texttt{b -{} b\_lo} of the row in the band and for the count.)
\item \csource{j-sift.c}{850}{851}: modeled by \lname{JPoints.siftBandsRow}
  (deviation of \lname{JPoints.siftBandsRow}: the state taken and returned, for the thread-local globals; the globals and the static data that it reads passed; integers for the \texttt{long}s, a natural number for the index \texttt{b -{} b\_lo} of the row in the band; \texttt{a \% 64} and \texttt{b \% 64}, in \texttt{[0, 64)}, for \texttt{a \& 63} and \texttt{b \& 63}, the same in two's complement, and the entry \texttt{(a \& 63) * 64 + (b \& 63)} of the flat array \texttt{alive2}; the \texttt{continue}s of the C code as returns of \texttt{false}, and whether \texttt{sift\_bands} returns \texttt{1} from the row, for its return.)
\item \csource{j-sift.c}{852}{856}: modeled by \lname{JPoints.siftMask}
  (deviation of \lname{JPoints.siftMask}: one function for the two copies in \texttt{sift\_bands} and \texttt{sift}, which are the same on \texttt{a}, \texttt{b}, \texttt{c0} and \texttt{mask} (in \texttt{sift} the call of \texttt{row\_setup} between them reads neither); the pair returned, for the local variables; the globals that it reads passed; integers for the \texttt{long}s, with \texttt{Int.tmod} and \texttt{Int.tdiv} for C's \texttt{\%} and \texttt{/}, and \texttt{height + 64} for \texttt{height + LONG\_LENGTH}, which the C code computes as a \texttt{size\_t} and converts back to a \texttt{long}, the same value; \texttt{a \% 64} and \texttt{b \% 64}, in \texttt{[0, 64)}, for \texttt{a \& 63} and \texttt{b \& 63}, the same in two's complement, and the entry \texttt{(a \& 63) * 64 + (b \& 63)} of the flat array \texttt{mask2}; the test that \texttt{a} and \texttt{b} are even for \texttt{!((a|b)\&1)}.)
\item \csource{j-sift.c}{857}{859}: modeled by \lname{JPoints.siftBandsRow}
  (deviation of \lname{JPoints.siftBandsRow}: the state taken and returned, for the thread-local globals; the globals and the static data that it reads passed; integers for the \texttt{long}s, a natural number for the index \texttt{b -{} b\_lo} of the row in the band; \texttt{a \% 64} and \texttt{b \% 64}, in \texttt{[0, 64)}, for \texttt{a \& 63} and \texttt{b \& 63}, the same in two's complement, and the entry \texttt{(a \& 63) * 64 + (b \& 63)} of the flat array \texttt{alive2}; the \texttt{continue}s of the C code as returns of \texttt{false}, and whether \texttt{sift\_bands} returns \texttt{1} from the row, for its return.)
\item \csource{j-sift.c}{860}{863}: modeled by \lname{JPoints.siftBandsTube}
  (deviation of \lname{JPoints.siftBandsTube}: the state taken and returned, for the thread-local globals; the globals and the static data that it reads passed; integers for the \texttt{long}s, and the number \texttt{hi -{} lo} of the words, which is positive for a range of the tube, added as a natural number; the loop over \texttt{i} from \texttt{lo} below \texttt{hi} as a loop over \texttt{i = lo + k} for \texttt{k < hi -{} lo}; whether \texttt{sift\_bands} returns \texttt{1} here, for its return.)
\item \csource{j-sift.c}{864}{864}: modeled by \lname{JPoints.siftBandsRow}
  (deviation of \lname{JPoints.siftBandsRow}: the state taken and returned, for the thread-local globals; the globals and the static data that it reads passed; integers for the \texttt{long}s, a natural number for the index \texttt{b -{} b\_lo} of the row in the band; \texttt{a \% 64} and \texttt{b \% 64}, in \texttt{[0, 64)}, for \texttt{a \& 63} and \texttt{b \& 63}, the same in two's complement, and the entry \texttt{(a \& 63) * 64 + (b \& 63)} of the flat array \texttt{alive2}; the \texttt{continue}s of the C code as returns of \texttt{false}, and whether \texttt{sift\_bands} returns \texttt{1} from the row, for its return.)
\item \csource{j-sift.c}{865}{865}: modeled by \lname{JPoints.siftBandsRegion}
  (deviation of \lname{JPoints.siftBandsRegion}: the state taken and returned, for the thread-local globals; the globals and the static data that it reads passed; integers for the \texttt{long}s \texttt{a}, \texttt{b}, \texttt{lo}, \texttt{hi}, \texttt{l}, \texttt{h} and \texttt{c0}, a natural number for the index \texttt{c} of the chunk; the second condition of the loop over the chunks, \texttt{chunks[c].w\_low < hi}, tested at the head of its body, with a \texttt{break}; whether \texttt{sift\_bands} returns \texttt{1} here, for its return.)
\item \csource{j-sift.c}{866}{867}: modeled by \lname{JPoints.seekChunk}
  (deviation of \lname{JPoints.seekChunk}: one function for the two pairs of loops, which are the same on \texttt{cur\_chunk} and the word, \texttt{i} in \texttt{tube\_word} and \texttt{lo} in \texttt{sift\_bands}, returning \texttt{cur\_chunk}; the static data passed, for the static \texttt{chunks}; a natural number for the index \texttt{cur\_chunk}, which stays in \texttt{[0, num\_chunks)} for a word of the row (for a word outside the row the C code reads beyond the array of the chunks, and the loops of the model need not end); integers for the \texttt{long}s of the words.)
\item \csource{j-sift.c}{868}{873}: modeled by \lname{JPoints.siftBandsRegion}
  (deviation of \lname{JPoints.siftBandsRegion}: the state taken and returned, for the thread-local globals; the globals and the static data that it reads passed; integers for the \texttt{long}s \texttt{a}, \texttt{b}, \texttt{lo}, \texttt{hi}, \texttt{l}, \texttt{h} and \texttt{c0}, a natural number for the index \texttt{c} of the chunk; the second condition of the loop over the chunks, \texttt{chunks[c].w\_low < hi}, tested at the head of its body, with a \texttt{break}; whether \texttt{sift\_bands} returns \texttt{1} here, for its return.)
\item \csource{j-sift.c}{874}{875}: modeled by \lname{JPoints.siftBandsRow}
  (deviation of \lname{JPoints.siftBandsRow}: the state taken and returned, for the thread-local globals; the globals and the static data that it reads passed; integers for the \texttt{long}s, a natural number for the index \texttt{b -{} b\_lo} of the row in the band; \texttt{a \% 64} and \texttt{b \% 64}, in \texttt{[0, 64)}, for \texttt{a \& 63} and \texttt{b \& 63}, the same in two's complement, and the entry \texttt{(a \& 63) * 64 + (b \& 63)} of the flat array \texttt{alive2}; the \texttt{continue}s of the C code as returns of \texttt{false}, and whether \texttt{sift\_bands} returns \texttt{1} from the row, for its return.)
\item \csource{j-sift.c}{876}{878}: modeled by \lname{JPoints.siftBands}
  (deviation of \lname{JPoints.siftBands}: the state taken and returned, for the thread-local globals; the globals and the static data that it reads passed; integers for the \texttt{long}s, with \texttt{Int.tmod} for C's \texttt{\%}; the loop \texttt{for(b\_lo = b\_first; b\_lo <= height; b\_lo += step)} over \texttt{b\_lo = b\_first + t step} for \texttt{t < (height -{} b\_first) / step + 1}, with Lean's \texttt{/}, which rounds down, none when \texttt{b\_first > height}, and the loop \texttt{for(b = b\_lo; b <= b\_hi; b += m)} over \texttt{b = b\_lo + k m} for \texttt{k < (b\_hi -{} b\_lo) / m + 1}; Booleans for the \texttt{int} \texttt{tube} and for the \texttt{int} returned.)
\item \csource{j-sift.c}{899}{910}: modeled by \lname{JPoints.walkParts}
  (deviation of \lname{JPoints.walkParts}: the parts returned as a walk, for the arguments \texttt{f1}, \texttt{hd}, \texttt{np} and \texttt{tl} through which the C code returns them; the walk passed, for the pointer \texttt{w}; the static data that it reads passed; natural numbers for the nonnegative \texttt{long}s; \texttt{JPoints.Double.ofInt} for the exact conversion of \texttt{rem -{} r} to a double, the operations of the doubles of \texttt{JPoints.Model.DoubleArith}, and \texttt{JPoints.Double.toInt64} for the conversion back to a \texttt{long}, which truncates; the arm \texttt{VW == 1} of the test of the constant \texttt{VW}.)
\item \csource{j-sift.c}{911}{915}: excluded (other): the arm of the vector walk (VW > 1) of walk\_parts, an if on the constant VW that the build modeled never takes
\item \csource{j-sift.c}{916}{917}: modeled by \lname{JPoints.walkParts}
  (deviation of \lname{JPoints.walkParts}: the parts returned as a walk, for the arguments \texttt{f1}, \texttt{hd}, \texttt{np} and \texttt{tl} through which the C code returns them; the walk passed, for the pointer \texttt{w}; the static data that it reads passed; natural numbers for the nonnegative \texttt{long}s; \texttt{JPoints.Double.ofInt} for the exact conversion of \texttt{rem -{} r} to a double, the operations of the doubles of \texttt{JPoints.Model.DoubleArith}, and \texttt{JPoints.Double.toInt64} for the conversion back to a \texttt{long}, which truncates; the arm \texttt{VW == 1} of the test of the constant \texttt{VW}.)
\item \csource{j-sift.c}{918}{922}: modeled by \lname{JPoints.walk}
  (deviation of \lname{JPoints.walk}: a function for the macro, the variables that it reads of the function in which it is expanded passed (the bit array and the rows \texttt{rowptr} of the state, \texttt{surv0}, \texttt{n}, \texttt{p}, and \texttt{f1}, \texttt{hd}, \texttt{np} and \texttt{tl} as the walk \texttt{w}), and the bit array returned; the operation an argument; indices into the bit array and into \texttt{sieve\_tab[n]} for the pointers \texttt{surv} and \texttt{siv1}, without the pointer \texttt{siv0}, at which the loop over a period ends after \texttt{p} words (\texttt{JPoints.walkWords}); the loop \texttt{for(j = np; j; j-{}-{})}, whose counter the body does not read, over \texttt{nper} steps, and the test \texttt{if((j = tl))}, which sets the counter of the tail, as \texttt{tail ≠ 0}; natural numbers for the counts; the arm \texttt{VW == 1} of the preprocessor.)
\item \csource{j-sift.c}{923}{923}: modeled by \lname{JPoints.walkWords}
  (deviation of \lname{JPoints.walkWords}: one function for the three loops of the macro, which are the same on their counts \texttt{hd}, \texttt{p} and \texttt{tl}: the loop \texttt{do … while(siv1 != siv0)}, which starts at \texttt{siv1 = siv0 -{} p} and moves \texttt{siv1} on by one word, runs \texttt{p} times; the loops \texttt{for(j = hd; j; j-{}-{})} and \texttt{for( ; j; j-{}-{})}, whose counter the body does not read, over \texttt{cnt} steps; indices for the pointers \texttt{surv} and \texttt{siv1}, and the table that \texttt{siv1} points into passed; natural numbers for the counts.)
\item \csource{j-sift.c}{924}{928}: modeled by \lname{JPoints.walk}
  (deviation of \lname{JPoints.walk}: a function for the macro, the variables that it reads of the function in which it is expanded passed (the bit array and the rows \texttt{rowptr} of the state, \texttt{surv0}, \texttt{n}, \texttt{p}, and \texttt{f1}, \texttt{hd}, \texttt{np} and \texttt{tl} as the walk \texttt{w}), and the bit array returned; the operation an argument; indices into the bit array and into \texttt{sieve\_tab[n]} for the pointers \texttt{surv} and \texttt{siv1}, without the pointer \texttt{siv0}, at which the loop over a period ends after \texttt{p} words (\texttt{JPoints.walkWords}); the loop \texttt{for(j = np; j; j-{}-{})}, whose counter the body does not read, over \texttt{nper} steps, and the test \texttt{if((j = tl))}, which sets the counter of the tail, as \texttt{tail ≠ 0}; natural numbers for the counts; the arm \texttt{VW == 1} of the preprocessor.)
\item \csource{j-sift.c}{929}{929}: modeled by \lname{JPoints.walkWords}
  (deviation of \lname{JPoints.walkWords}: one function for the three loops of the macro, which are the same on their counts \texttt{hd}, \texttt{p} and \texttt{tl}: the loop \texttt{do … while(siv1 != siv0)}, which starts at \texttt{siv1 = siv0 -{} p} and moves \texttt{siv1} on by one word, runs \texttt{p} times; the loops \texttt{for(j = hd; j; j-{}-{})} and \texttt{for( ; j; j-{}-{})}, whose counter the body does not read, over \texttt{cnt} steps; indices for the pointers \texttt{surv} and \texttt{siv1}, and the table that \texttt{siv1} points into passed; natural numbers for the counts.)
\item \csource{j-sift.c}{930}{932}: modeled by \lname{JPoints.walk}
  (deviation of \lname{JPoints.walk}: a function for the macro, the variables that it reads of the function in which it is expanded passed (the bit array and the rows \texttt{rowptr} of the state, \texttt{surv0}, \texttt{n}, \texttt{p}, and \texttt{f1}, \texttt{hd}, \texttt{np} and \texttt{tl} as the walk \texttt{w}), and the bit array returned; the operation an argument; indices into the bit array and into \texttt{sieve\_tab[n]} for the pointers \texttt{surv} and \texttt{siv1}, without the pointer \texttt{siv0}, at which the loop over a period ends after \texttt{p} words (\texttt{JPoints.walkWords}); the loop \texttt{for(j = np; j; j-{}-{})}, whose counter the body does not read, over \texttt{nper} steps, and the test \texttt{if((j = tl))}, which sets the counter of the tail, as \texttt{tail ≠ 0}; natural numbers for the counts; the arm \texttt{VW == 1} of the preprocessor.)
\item \csource{j-sift.c}{933}{933}: modeled by \lname{JPoints.walkWords}
  (deviation of \lname{JPoints.walkWords}: one function for the three loops of the macro, which are the same on their counts \texttt{hd}, \texttt{p} and \texttt{tl}: the loop \texttt{do … while(siv1 != siv0)}, which starts at \texttt{siv1 = siv0 -{} p} and moves \texttt{siv1} on by one word, runs \texttt{p} times; the loops \texttt{for(j = hd; j; j-{}-{})} and \texttt{for( ; j; j-{}-{})}, whose counter the body does not read, over \texttt{cnt} steps; indices for the pointers \texttt{surv} and \texttt{siv1}, and the table that \texttt{siv1} points into passed; natural numbers for the counts.)
\item \csource{j-sift.c}{934}{935}: modeled by \lname{JPoints.walk}
  (deviation of \lname{JPoints.walk}: a function for the macro, the variables that it reads of the function in which it is expanded passed (the bit array and the rows \texttt{rowptr} of the state, \texttt{surv0}, \texttt{n}, \texttt{p}, and \texttt{f1}, \texttt{hd}, \texttt{np} and \texttt{tl} as the walk \texttt{w}), and the bit array returned; the operation an argument; indices into the bit array and into \texttt{sieve\_tab[n]} for the pointers \texttt{surv} and \texttt{siv1}, without the pointer \texttt{siv0}, at which the loop over a period ends after \texttt{p} words (\texttt{JPoints.walkWords}); the loop \texttt{for(j = np; j; j-{}-{})}, whose counter the body does not read, over \texttt{nper} steps, and the test \texttt{if((j = tl))}, which sets the counter of the tail, as \texttt{tail ≠ 0}; natural numbers for the counts; the arm \texttt{VW == 1} of the preprocessor.)
\item \csource{j-sift.c}{936}{936}: modeled by \lname{JPoints.andOp}
  (deviation of \lname{JPoints.andOp}: the value returned, for the assignment \texttt{d \&= s} of the macro, which \texttt{JPoints.walkWords} stores in the word of the bit array.)
\item \csource{j-sift.c}{937}{937}: modeled by \lname{JPoints.fillOp}
  (deviation of \lname{JPoints.fillOp}: the mask passed, for the local variable of the function in which the macro is expanded; the value returned, for the assignment to the word \texttt{d} of the macro, which \texttt{JPoints.walkWords} stores, and the old value of \texttt{d} an argument that it does not read, so that both operations have the type of the operation that \texttt{JPoints.walk} takes.)
\item \csource{j-sift.c}{938}{958}: excluded (other): the macros of the vector walk (VW > 1), selected by the \#else of the scalar ones: VSHIFT, WALK over vectors of VW words, AND\_OP and FILL\_OP on vectors
\item \csource{j-sift.c}{963}{973}: modeled by \lname{JPoints.scanWords}
  (deviation of \lname{JPoints.scanWords}: the state taken and returned, for the thread-local globals; the globals and the static data that it reads passed; the chunk passed, for the pointer \texttt{ch}; an index into the bit array for the pointer \texttt{surv}; integers for the \texttt{long}s \texttt{a} and \texttt{b}, natural numbers for the positions; the arm \texttt{VW == 1} of the preprocessor; a Boolean for the \texttt{int}.)
\item \csource{j-sift.c}{974}{992}: excluded (other): the arm of the vector walk (VW > 1) of scan\_words, which tests a vector at a time
\item \csource{j-sift.c}{993}{994}: modeled by \lname{JPoints.scanWords}
  (deviation of \lname{JPoints.scanWords}: the state taken and returned, for the thread-local globals; the globals and the static data that it reads passed; the chunk passed, for the pointer \texttt{ch}; an index into the bit array for the pointer \texttt{surv}; integers for the \texttt{long}s \texttt{a} and \texttt{b}, natural numbers for the positions; the arm \texttt{VW == 1} of the preprocessor; a Boolean for the \texttt{int}.)
\item \csource{j-sift.c}{1001}{1005}: modeled by \lname{JPoints.sift0}
  (deviation of \lname{JPoints.sift0}: the state taken and returned, for the thread-local globals, its bit array taken out of it while the passes write it; the globals and the static data that it reads passed; the chunk read from the array of the chunks, for the pointer \texttt{ch}; index \texttt{0} for the pointer \texttt{surv0}; the walks passed whole, for the copies of their fields in \texttt{f1}, \texttt{hd}, \texttt{np} and \texttt{tl}; integers for the \texttt{long}s \texttt{a}, \texttt{b} and \texttt{c0}, natural numbers for \texttt{k}, \texttt{range} and \texttt{n}; a Boolean for the \texttt{int}.)
\item \csource{j-sift.c}{1006}{1010}: excluded (other): the vector of the mask and the number of vectors of the chunk in sift0, for the vector walk (VW > 1)
\item \csource{j-sift.c}{1012}{1025}: modeled by \lname{JPoints.sift0Fill}
  (deviation of \lname{JPoints.sift0Fill}: the bit array and the variable \texttt{n} of \texttt{sift0} returned, for the bit array that the C code writes and the local variable; the globals and the static data that it reads passed, and the rows \texttt{rowptr} of the state, for the thread-local global; the chunk passed, for the pointer \texttt{ch}; an index for the pointer \texttt{surv}; the loop \texttt{for(i = range; i; i-{}-{})}, whose counter the body does not read, over \texttt{range} steps; natural numbers for the counts and for \texttt{n}.)
\item \csource{j-sift.c}{1026}{1029}: modeled by \lname{JPoints.sift0Masks}
  (deviation of \lname{JPoints.sift0Masks}: the bit array taken and returned; the globals and the static data that it reads passed; the chunk passed, for the pointer \texttt{ch}; integers for the \texttt{long} \texttt{c0} and the index of its word, \texttt{c0 \% 64}, in \texttt{[0, 64)}, for \texttt{c0 \& LONG\_MASK}, the same in two's complement; natural numbers for \texttt{k} and \texttt{range}.)
\item \csource{j-sift.c}{1030}{1032}: excluded (instrumentation): the bounds of the chunk printed by sift0 under DEBUG >= 1
\item \csource{j-sift.c}{1034}{1039}: modeled by \lname{JPoints.sift0}
  (deviation of \lname{JPoints.sift0}: the state taken and returned, for the thread-local globals, its bit array taken out of it while the passes write it; the globals and the static data that it reads passed; the chunk read from the array of the chunks, for the pointer \texttt{ch}; index \texttt{0} for the pointer \texttt{surv0}; the walks passed whole, for the copies of their fields in \texttt{f1}, \texttt{hd}, \texttt{np} and \texttt{tl}; integers for the \texttt{long}s \texttt{a}, \texttt{b} and \texttt{c0}, natural numbers for \texttt{k}, \texttt{range} and \texttt{n}; a Boolean for the \texttt{int}.)
\item \csource{j-sift.c}{1040}{1044}: excluded (other): the padding beyond the chunk cleared for the scan of the vector walk (VW > 1) in sift0
\item \csource{j-sift.c}{1045}{1046}: modeled by \lname{JPoints.sift0}
  (deviation of \lname{JPoints.sift0}: the state taken and returned, for the thread-local globals, its bit array taken out of it while the passes write it; the globals and the static data that it reads passed; the chunk read from the array of the chunks, for the pointer \texttt{ch}; index \texttt{0} for the pointer \texttt{surv0}; the walks passed whole, for the copies of their fields in \texttt{f1}, \texttt{hd}, \texttt{np} and \texttt{tl}; integers for the \texttt{long}s \texttt{a}, \texttt{b} and \texttt{c0}, natural numbers for \texttt{k}, \texttt{range} and \texttt{n}; a Boolean for the \texttt{int}.)
\item \csource{j-sift.c}{1053}{1058}: modeled by \lname{JPoints.siftRange}
  (deviation of \lname{JPoints.siftRange}: the state taken and returned, for the thread-local globals, its bit array taken out of it while the passes write it; the globals and the static data that it reads passed; the chunk read from the array of the chunks, for the pointer \texttt{ch}; the index \texttt{off} for the pointer \texttt{surv0}; integers for the \texttt{long}s \texttt{a}, \texttt{b}, \texttt{lo}, \texttt{hi}, \texttt{c0} and \texttt{cw}, and natural numbers for \texttt{k}, \texttt{range}, \texttt{off}, \texttt{len}, \texttt{lo0}, \texttt{hi0} and \texttt{n}, which the C code keeps nonnegative; the parts of a walk as a walk, for the variables \texttt{f1}, \texttt{hd}, \texttt{np} and \texttt{tl}; a Boolean for the \texttt{int}.)
\item \csource{j-sift.c}{1059}{1069}: excluded (other): the vector of the mask and the range widened to whole vectors in sift\_range, for the vector walk (VW > 1)
\item \csource{j-sift.c}{1070}{1070}: modeled by \lname{JPoints.siftRange}
  (deviation of \lname{JPoints.siftRange}: the state taken and returned, for the thread-local globals, its bit array taken out of it while the passes write it; the globals and the static data that it reads passed; the chunk read from the array of the chunks, for the pointer \texttt{ch}; the index \texttt{off} for the pointer \texttt{surv0}; integers for the \texttt{long}s \texttt{a}, \texttt{b}, \texttt{lo}, \texttt{hi}, \texttt{c0} and \texttt{cw}, and natural numbers for \texttt{k}, \texttt{range}, \texttt{off}, \texttt{len}, \texttt{lo0}, \texttt{hi0} and \texttt{n}, which the C code keeps nonnegative; the parts of a walk as a walk, for the variables \texttt{f1}, \texttt{hd}, \texttt{np} and \texttt{tl}; a Boolean for the \texttt{int}.)
\item \csource{j-sift.c}{1072}{1085}: modeled by \lname{JPoints.siftRangeFill}
  (deviation of \lname{JPoints.siftRangeFill}: the bit array and the variable \texttt{n} of \texttt{sift\_range} returned, for the bit array that the C code writes and the local variable; the globals and the static data that it reads passed, and the rows \texttt{rowptr} of the state, for the thread-local global; the chunk passed, for the pointer \texttt{ch}; an index for the pointer \texttt{surv}; the parts of the walk as a walk, for the variables \texttt{f1}, \texttt{hd}, \texttt{np} and \texttt{tl}; the loop \texttt{for(i = len; i; i-{}-{})}, whose counter the body does not read, over \texttt{len} steps; natural numbers for the counts and for \texttt{n}.)
\item \csource{j-sift.c}{1086}{1086}: excluded (other): the casts of hd, np and tl to void, which silence the warning of the compiler that the vector walk (VW > 1) does not read them; they do nothing
\item \csource{j-sift.c}{1087}{1089}: modeled by \lname{JPoints.siftRangeMasks}
  (deviation of \lname{JPoints.siftRangeMasks}: the bit array taken and returned; the globals and the static data that it reads passed; integers for the \texttt{long}s \texttt{c0} and \texttt{cw}, \texttt{c0 \% 64}, in \texttt{[0, 64)}, for \texttt{c0 \& LONG\_MASK}, the same in two's complement; natural numbers for \texttt{k}, \texttt{range}, \texttt{off} and \texttt{len}.)
\item \csource{j-sift.c}{1090}{1092}: excluded (instrumentation): the bounds of the range printed by sift\_range under DEBUG >= 1
\item \csource{j-sift.c}{1093}{1098}: modeled by \lname{JPoints.siftRange}
  (deviation of \lname{JPoints.siftRange}: the state taken and returned, for the thread-local globals, its bit array taken out of it while the passes write it; the globals and the static data that it reads passed; the chunk read from the array of the chunks, for the pointer \texttt{ch}; the index \texttt{off} for the pointer \texttt{surv0}; integers for the \texttt{long}s \texttt{a}, \texttt{b}, \texttt{lo}, \texttt{hi}, \texttt{c0} and \texttt{cw}, and natural numbers for \texttt{k}, \texttt{range}, \texttt{off}, \texttt{len}, \texttt{lo0}, \texttt{hi0} and \texttt{n}, which the C code keeps nonnegative; the parts of a walk as a walk, for the variables \texttt{f1}, \texttt{hd}, \texttt{np} and \texttt{tl}; a Boolean for the \texttt{int}.)
\item \csource{j-sift.c}{1099}{1103}: excluded (other): the padding beyond the range cleared for the scan of the vector walk (VW > 1) in sift\_range
\item \csource{j-sift.c}{1104}{1105}: modeled by \lname{JPoints.siftRange}
  (deviation of \lname{JPoints.siftRange}: the state taken and returned, for the thread-local globals, its bit array taken out of it while the passes write it; the globals and the static data that it reads passed; the chunk read from the array of the chunks, for the pointer \texttt{ch}; the index \texttt{off} for the pointer \texttt{surv0}; integers for the \texttt{long}s \texttt{a}, \texttt{b}, \texttt{lo}, \texttt{hi}, \texttt{c0} and \texttt{cw}, and natural numbers for \texttt{k}, \texttt{range}, \texttt{off}, \texttt{len}, \texttt{lo0}, \texttt{hi0} and \texttt{n}, which the C code keeps nonnegative; the parts of a walk as a walk, for the variables \texttt{f1}, \texttt{hd}, \texttt{np} and \texttt{tl}; a Boolean for the \texttt{int}.)
\item \csource{j-sift.c}{1113}{1114}: modeled by \lname{JPoints.rowSetup}
  (deviation of \lname{JPoints.rowSetup}: the state taken and returned, for the thread-local global \texttt{rowptr} and the statics of \texttt{row\_setup} (\texttt{JPoints.RowSetupCache}); the globals and the static data that it reads passed; integers for the \texttt{long}s \texttt{a}, \texttt{b} and \texttt{d}.)
\item \csource{j-sift.c}{1115}{1116}: modeled by \lname{JPoints.RowSetupCache}, \lname{JPoints.RowSetupCache.initial}
  (deviation of \lname{JPoints.RowSetupCache}: integers for the \texttt{long}s \texttt{last\_a}, \texttt{last\_b} and \texttt{last\_d}, natural numbers for the residues and the offsets, which the C code keeps nonnegative; a structure, for the thread-local statics of the one thread.)
\item \csource{j-sift.c}{1117}{1119}: modeled by \lname{JPoints.rowSetup}
  (deviation of \lname{JPoints.rowSetup}: the state taken and returned, for the thread-local global \texttt{rowptr} and the statics of \texttt{row\_setup} (\texttt{JPoints.RowSetupCache}); the globals and the static data that it reads passed; integers for the \texttt{long}s \texttt{a}, \texttt{b} and \texttt{d}.)
\item \csource{j-sift.c}{1120}{1124}: modeled by \lname{JPoints.rowSetupAddSmall}
  (deviation of \lname{JPoints.rowSetupAddSmall}: the rows and the residues taken and returned, for the thread-local global \texttt{rowptr} and the static \texttt{bres}; the globals and the static data that it reads passed; indices into \texttt{sieve\_tab[n]} for the pointers \texttt{rowptr[n]}; natural numbers for \texttt{d} and the residues.)
\item \csource{j-sift.c}{1125}{1128}: modeled by \lname{JPoints.rowSetup}
  (deviation of \lname{JPoints.rowSetup}: the state taken and returned, for the thread-local global \texttt{rowptr} and the statics of \texttt{row\_setup} (\texttt{JPoints.RowSetupCache}); the globals and the static data that it reads passed; integers for the \texttt{long}s \texttt{a}, \texttt{b} and \texttt{d}.)
\item \csource{j-sift.c}{1129}{1138}: modeled by \lname{JPoints.rowSetupResidues}
  (deviation of \lname{JPoints.rowSetupResidues}: the residues and the offsets taken and returned, for the statics; the globals and the static data that it reads passed; integers for the \texttt{long}s \texttt{d}, \texttt{p} and \texttt{r}, and natural numbers for the residues and the offsets that the C code stores, the residue \texttt{r} in \texttt{[0, p)} after the corrections; \texttt{JPoints.Double.ofInt} for the exact conversion of \texttt{d} to a double, the product of \texttt{JPoints.Model.DoubleArith}, and \texttt{JPoints.Double.toInt64} for the conversion back to a \texttt{long}, which truncates.)
\item \csource{j-sift.c}{1139}{1140}: modeled by \lname{JPoints.rowSetup}
  (deviation of \lname{JPoints.rowSetup}: the state taken and returned, for the thread-local global \texttt{rowptr} and the statics of \texttt{row\_setup} (\texttt{JPoints.RowSetupCache}); the globals and the static data that it reads passed; integers for the \texttt{long}s \texttt{a}, \texttt{b} and \texttt{d}.)
\item \csource{j-sift.c}{1141}{1145}: modeled by \lname{JPoints.rowSetupAddResidues}
  (deviation of \lname{JPoints.rowSetupAddResidues}: the rows and the residues taken and returned, for the thread-local global \texttt{rowptr} and the static \texttt{bres}; the globals and the static data that it reads passed; indices into \texttt{sieve\_tab[n]} for the pointers \texttt{rowptr[n]}; natural numbers for the residues and the offsets.)
\item \csource{j-sift.c}{1146}{1147}: modeled by \lname{JPoints.rowSetup}
  (deviation of \lname{JPoints.rowSetup}: the state taken and returned, for the thread-local global \texttt{rowptr} and the statics of \texttt{row\_setup} (\texttt{JPoints.RowSetupCache}); the globals and the static data that it reads passed; integers for the \texttt{long}s \texttt{a}, \texttt{b} and \texttt{d}.)
\item \csource{j-sift.c}{1148}{1153}: modeled by \lname{JPoints.rowSetupDivide}
  (deviation of \lname{JPoints.rowSetupDivide}: the rows and the residues taken and returned, for the thread-local global \texttt{rowptr} and the static \texttt{bres}; the globals and the static data that it reads passed; integers for the \texttt{long}s, with \texttt{Int.tmod} for C's \texttt{\%}; the index into \texttt{sieve\_tab[n]} for the pointer \texttt{rowptr[n]}, a natural number, which is the index for the \texttt{a ≥ 0} of the calls; a natural number for the residue \texttt{bp}.)
\item \csource{j-sift.c}{1154}{1158}: modeled by \lname{JPoints.rowSetup}
  (deviation of \lname{JPoints.rowSetup}: the state taken and returned, for the thread-local global \texttt{rowptr} and the statics of \texttt{row\_setup} (\texttt{JPoints.RowSetupCache}); the globals and the static data that it reads passed; integers for the \texttt{long}s \texttt{a}, \texttt{b} and \texttt{d}.)
\item \csource{j-sift.c}{1160}{1164}: modeled by \lname{JPoints.sift}
  (deviation of \lname{JPoints.sift}: the state taken and returned, for the thread-local globals; the globals and the static data that it reads passed; integers for the \texttt{long}s; \texttt{c0} and the mask computed together before \texttt{row\_setup} (\texttt{JPoints.siftMask}), which reads neither; a Boolean for the \texttt{int}.)
\item \csource{j-sift.c}{1165}{1165}: modeled by \lname{JPoints.siftMask}
  (deviation of \lname{JPoints.siftMask}: one function for the two copies in \texttt{sift\_bands} and \texttt{sift}, which are the same on \texttt{a}, \texttt{b}, \texttt{c0} and \texttt{mask} (in \texttt{sift} the call of \texttt{row\_setup} between them reads neither); the pair returned, for the local variables; the globals that it reads passed; integers for the \texttt{long}s, with \texttt{Int.tmod} and \texttt{Int.tdiv} for C's \texttt{\%} and \texttt{/}, and \texttt{height + 64} for \texttt{height + LONG\_LENGTH}, which the C code computes as a \texttt{size\_t} and converts back to a \texttt{long}, the same value; \texttt{a \% 64} and \texttt{b \% 64}, in \texttt{[0, 64)}, for \texttt{a \& 63} and \texttt{b \& 63}, the same in two's complement, and the entry \texttt{(a \& 63) * 64 + (b \& 63)} of the flat array \texttt{mask2}; the test that \texttt{a} and \texttt{b} are even for \texttt{!((a|b)\&1)}.)
\item \csource{j-sift.c}{1166}{1166}: modeled by \lname{JPoints.sift}
  (deviation of \lname{JPoints.sift}: the state taken and returned, for the thread-local globals; the globals and the static data that it reads passed; integers for the \texttt{long}s; \texttt{c0} and the mask computed together before \texttt{row\_setup} (\texttt{JPoints.siftMask}), which reads neither; a Boolean for the \texttt{int}.)
\item \csource{j-sift.c}{1167}{1167}: modeled by \lname{JPoints.siftMask}
  (deviation of \lname{JPoints.siftMask}: one function for the two copies in \texttt{sift\_bands} and \texttt{sift}, which are the same on \texttt{a}, \texttt{b}, \texttt{c0} and \texttt{mask} (in \texttt{sift} the call of \texttt{row\_setup} between them reads neither); the pair returned, for the local variables; the globals that it reads passed; integers for the \texttt{long}s, with \texttt{Int.tmod} and \texttt{Int.tdiv} for C's \texttt{\%} and \texttt{/}, and \texttt{height + 64} for \texttt{height + LONG\_LENGTH}, which the C code computes as a \texttt{size\_t} and converts back to a \texttt{long}, the same value; \texttt{a \% 64} and \texttt{b \% 64}, in \texttt{[0, 64)}, for \texttt{a \& 63} and \texttt{b \& 63}, the same in two's complement, and the entry \texttt{(a \& 63) * 64 + (b \& 63)} of the flat array \texttt{mask2}; the test that \texttt{a} and \texttt{b} are even for \texttt{!((a|b)\&1)}.)
\item \csource{j-sift.c}{1168}{1170}: excluded (instrumentation): the row (a, b) printed by sift under DEBUG >= 1
\item \csource{j-sift.c}{1171}{1171}: modeled by \lname{JPoints.sift}
  (deviation of \lname{JPoints.sift}: the state taken and returned, for the thread-local globals; the globals and the static data that it reads passed; integers for the \texttt{long}s; \texttt{c0} and the mask computed together before \texttt{row\_setup} (\texttt{JPoints.siftMask}), which reads neither; a Boolean for the \texttt{int}.)
\item \csource{j-sift.c}{1173}{1175}: modeled by \lname{JPoints.siftMask}
  (deviation of \lname{JPoints.siftMask}: one function for the two copies in \texttt{sift\_bands} and \texttt{sift}, which are the same on \texttt{a}, \texttt{b}, \texttt{c0} and \texttt{mask} (in \texttt{sift} the call of \texttt{row\_setup} between them reads neither); the pair returned, for the local variables; the globals that it reads passed; integers for the \texttt{long}s, with \texttt{Int.tmod} and \texttt{Int.tdiv} for C's \texttt{\%} and \texttt{/}, and \texttt{height + 64} for \texttt{height + LONG\_LENGTH}, which the C code computes as a \texttt{size\_t} and converts back to a \texttt{long}, the same value; \texttt{a \% 64} and \texttt{b \% 64}, in \texttt{[0, 64)}, for \texttt{a \& 63} and \texttt{b \& 63}, the same in two's complement, and the entry \texttt{(a \& 63) * 64 + (b \& 63)} of the flat array \texttt{mask2}; the test that \texttt{a} and \texttt{b} are even for \texttt{!((a|b)\&1)}.)
\item \csource{j-sift.c}{1176}{1180}: modeled by \lname{JPoints.sift}
  (deviation of \lname{JPoints.sift}: the state taken and returned, for the thread-local globals; the globals and the static data that it reads passed; integers for the \texttt{long}s; \texttt{c0} and the mask computed together before \texttt{row\_setup} (\texttt{JPoints.siftMask}), which reads neither; a Boolean for the \texttt{int}.)
\end{itemize}

\section*{\texttt{j-points.h}}
\begin{itemize}
\item \csource{j-points.h}{27}{29}: excluded (instrumentation): the default of DEBUG, 0, the build modeled; DEBUG > 0 selects the debug build, with the small table of primes (lines 35-40, and j-points.c:54-57) and the output of internals in init\_sieve, sift0, sift\_range and sift, all excluded
\item \csource{j-points.h}{35}{40}: excluded (instrumentation): the sizes of the table of primes of the debug build (DEBUG > 0), five odd primes up to 13, with the \#if that selects them and the \#else of the sizes of the default build (lines 41-44)
\item \csource{j-points.h}{41}{41}: modeled by \lname{JPoints.numPrimes}
\item \csource{j-points.h}{42}{42}: modeled by \lname{JPoints.maxPrime}
\item \csource{j-points.h}{43}{43}: modeled by \lname{JPoints.maxPrimeEven}
\item \csource{j-points.h}{44}{44}: modeled by \lname{JPoints.maxTablePrime}
\item \csource{j-points.h}{45}{45}: excluded (instrumentation): the end of the selection of the sizes of the table of primes by DEBUG
\item \csource{j-points.h}{47}{47}: excluded (other): the macro FLOOR, which no code of j-points uses (j-sift.c:135 computes the floor of -height / LONG\_LENGTH by a shift)
\item \csource{j-points.h}{48}{48}: modeled by \lname{JPoints.ceilDiv}
\item \csource{j-points.h}{50}{50}: modeled by \lname{JPoints.longLength}
\item \csource{j-points.h}{52}{54}: modeled by \lname{JPoints.longShift}
\item \csource{j-points.h}{55}{55}: modeled by \lname{JPoints.longMask}
\item \csource{j-points.h}{56}{58}: modeled by \lname{JPoints.maxHeight}
\item \csource{j-points.h}{59}{59}: modeled by \lname{JPoints.maxSearchHeight}
\item \csource{j-points.h}{64}{64}: modeled by \lname{JPoints.halfMask}
\item \csource{j-points.h}{66}{67}: excluded (other): the type bit\_array of the words of the bit arrays, unsigned long, and the macro zero, the word 0; the words of the model are UInt64, and zero is 0
\item \csource{j-points.h}{75}{77}: excluded (other): the default of VW, the width in words of the passes of the sieve: the build modeled has VW = 1, the scalar walk; the model writes 1 where the scalar code reads VW (sieve\_period, sieve\_rowlen, and the length array\_size + VW of the bit array in JPoints.SiftState.initial), takes the arm VW == 1 where the code tests VW (lines 175, 898, 914 and 961 of j-sift.c) and the costs of the \#else of VW == 4 (lines 85-91 of j-points.c); the vector builds, VW = 2 and VW = 4, are not modeled
\item \csource{j-points.h}{79}{129}: excluded (other): the declarations of the globals and the functions that j-points.c and j-sift.c share; the model has no headers, and each global and function is recorded where the C code defines it
\end{itemize}
